% -----------------------------------------------------------
\section{1834}
% -----------------------------------------------------------
	\label{EMS-1834}
	
	% ----------------------
	\subsection{Joseph Smith}
	% ----------------------
		\label{EMS-1834-JS}

		% -----
		\subsubsection{Introduction to 1834}
		% -----
		
			sdf
		
		% -----
		\subsubsection{Prayer, 11 January 1834}
		% -----
			\label{EMS-1834-JS-11-JAN-1834}
			% \label{EMS-1830-JS-DAY-3LETTERMONTH-YEAR}
			
			\begin{description}
					\item[Print Sources]
				\end{description}

					\begin{tabular}{p{0.5in}p{1in}p{0.5in}p{1in}}
					JSP	 	& D3:403--407; J1:25--26		& JSR 		& --- \\
					DC	 	& --- 				& HC 		& 2:2--3 \\
					HDDC 	& ---				& TCHC 		& 2:6--7 \\
					\end{tabular}
					% HDDC = woodford's DC diss; JSR = Marquardt's JS Revelations; TCHC = Vogel HC
				
				\begin{description}
					\item[Online Access] \href{https://www.josephsmithpapers.org/paper-summary/prayer-11-january-1834/}{Here}
					% \item[Analysis] \hyperref[XX]{Here}
				\end{description}
				
				\begin{description}
					\item[Background]
				\end{description}
				
					% It might also be useful to give a two sentence context of the period: e.g., "This speech was given in Nauvoo to a small congregation one month prior to the destruction of the Nauvoo Expositor. These and other tensions may have been on the prophet's mind when he gave the speech."
				
				\begin{description}
					\item[Original Source]
				\end{description}
				
					% brief description of original source material, with reference to JSP or somewhere else for more info
					% Background material: it might be helpful to have a note that says whether a given document was presented as a speech, was written in a journal, etc.
					% Also a comment here about how reliable the source might be, or what shape it is in, if there are large omissions or parts that are hard to understand, and so on.
				
				\begin{description}
					\item[Text]
				\end{description}
				
						\hypertarget{EMS-1834-JS-11-JAN-1834-a}%
					January
						\marginnote{a}%
					11, 1834. This evening Joseph Smith Jr, Frederick G. Williams, Newel K. Whitney, John Johnson, Oliver Cowdery, and Orson Hyde united in prayer and asked the Lord to grant the following petition:
					
					Firstly, That the Lord would grant that our lives might be precious in his sight, that he would watch over our persons and give his angels charge concerning us and our families that no evil nor unseen hand might be permitted to harm us.
					
						\hypertarget{EMS-1834-JS-11-JAN-1834-b}%
					Secondly,
						\marginnote{b}%
					That the Lord would also hold the lives of all the United Firm, and not suffer that any of them shall be taken.
					
					Thirdly, That the Lord would grant that our brother Joseph might prevail over his enemy, even Docter P. [Doctor Philastus] Hurlbut, who has threatened his life, whom brother Joseph has <caused to be> taken with a precept; that the Lord would fill the heart of the Court with a spirit to do justice, and cause that the law of the land may be magnified in bringing him to justice.
					
						\hypertarget{EMS-1834-JS-11-JAN-1834-c}%
					Fourthly,
						\marginnote{c}%
					That the Lord would provide, in the order of his Providence, the bishop of this Church with means sufficient to discharge every debt that the Firm owes, in due season, that the Church may not be braught into disrepute, and the saints be afflicted by the hands of their enemies.
					Fifthly, That the Lord would protect our printing press from the hands of evil men, and give us means to send forth his word, even his gospel that the ears of all may hear it, and also that we may print his scriptures; and also that he would give those who were appointed to conduct the press, wisdom sufficient that the cause may not be hindered, but that men's eyes may thereby be opened to see the truth.
					
						\hypertarget{EMS-1834-JS-11-JAN-1834-d}%
					Sixthly,
						\marginnote{d}%
					That the Lord would deliver Zion, and gather in his scattered people, to possess it in peace; and also, while in their dispersion, that he would provide for them that they perish not with hunger nor cold. And finally, that God in the name of Jesus would gather his elect speedily, and unveil his face that his saints might behold his glory and dwell with him; Amen.
					
		% -----
		\subsubsection{Letter to the Church in Clay County, Missouri, 22 January 1834}
		% -----
			\label{EMS-1834-JS-22-JAN-1834}
			% \label{EMS-1830-JS-DAY-3LETTERMONTH-YEAR}
			
			\begin{description}
					\item[Print Sources]
				\end{description}

					\begin{tabular}{p{0.5in}p{1in}p{0.5in}p{1in}}
					JSP	 	& D3:407--412		& JSR 		& --- \\
					DC	 	& --- 		& HC 		& --- \\
					HDDC 	& ---		& TCHC 		& --- \\
					\end{tabular}
					% HDDC = woodford's DC diss; JSR = Marquardt's JS Revelations; TCHC = Vogel HC
				
				\begin{description}
					\item[Online Access] \href{https://www.josephsmithpapers.org/paper-summary/letter-to-the-church-in-clay-county-missouri-22-january-1834/1}{Here}
					% \item[Analysis] \hyperref[XX]{Here}
				\end{description}
				
				\begin{description}
					\item[Background]
				\end{description}
				
					% It might also be useful to give a two sentence context of the period: e.g., "This speech was given in Nauvoo to a small congregation one month prior to the destruction of the Nauvoo Expositor. These and other tensions may have been on the prophet's mind when he gave the speech."
				
				\begin{description}
					\item[Original Source]
				\end{description}
				
					% brief description of original source material, with reference to JSP or somewhere else for more info
					% Background material: it might be helpful to have a note that says whether a given document was presented as a speech, was written in a journal, etc.
					% Also a comment here about how reliable the source might be, or what shape it is in, if there are large omissions or parts that are hard to understand, and so on.
				
				\begin{description}
					\item[Text]
				\end{description}
				
						\hypertarget{EMS-1834-JS-22-JAN-1834-a}%
					...Inquire
						\marginnote{a}%
					of Bro [Thomas B.] Marsh and find out the entire secret of mixing or compounding lead and Antimo[n]y so as to make type mettle [metal] and write us concerning it Bro Joseph tells me that he has sent another \$50= note making \$100, to you, write us concerning it, there is a prospect of the eastern churches doing something pretty handsome toward the deliverance of Zion in the course of a year if Zion is not delivered otherwise
						\hypertarget{EMS-1834-JS-22-JAN-1834-b}%
					Tho
						\marginnote{b}%
					the Lord says this affliction came upon you because of your sin polluting your inheritances \&c yet there is an exception of some namely the heads of Zion for the Lord said your brethren in Zion begin to repent, and the Angels rejoice over them \&c you will also see an exception at the top of the 2d collum of this revelation Therefore this affliction came upon the Church to chastin those in transgression, and prepare the hearts of those who had repented for an end[o]wment from the Lord. We shall not be able to send you any more money at present unless the Lord put it into our \sout{power} hands unexpectedly.
						\hypertarget{EMS-1834-JS-22-JAN-1834-c}%
					our
						\marginnote{c}%
					calls for money are many and pressing. There is not quite so much danger of a Mob upon us as there has been. The hand of the Lord has thus far been stretched out to protect us, Doctor P[hilastus] Hurlbut an apostate Elder from this Church has been to the State of New York and gathered up all the rediculous stories that could be invented and some affidavits respecting the character of Bro Joseph and the Smith family and he exhibeted them to numerous congregations in Chagrin Kirtland Mentor and Painesville and fired the minds of the people with much indignation against Bro Joseph and the Church.
						\hypertarget{EMS-1834-JS-22-JAN-1834-d}%
					Hurlbut
						\marginnote{d}%
					also made many harsh threats \&c that he would take the life of Bro Joseph if he could not distroy Mormonism without. Bro Joseph took him with a peace warrant and after 3 days trial and investigating the merits of our religion in the town of Painesville by able attorney[s] on both sides he was bound over to the County Court. thus his influence was pritty much distroyed, and since the trial the spirit of hostility seames to be broken down in a good degree but how long it will continue so we cannot say
					
						\hypertarget{EMS-1834-JS-22-JAN-1834-e}%
					You
						\marginnote{e}%
					purchaced your inheritence with money, Therefore behold you are blessed. you have not purchaced your lands by the shedding of blood consiquently you do not come under the censure of that commandment which says ``if by blood lo your enemies are upon you and ye shall be driven from city to city'' give youselves no uneasiness on this account...
					
		% -----
		\subsubsection{Letter to the Church, circa February 1834}
		% -----
			\label{EMS-1834-JS-CIR-FEB-1834}
			% \label{EMS-1830-JS-DAY-3LETTERMONTH-YEAR}
			
			\begin{description}
					\item[Print Sources]
				\end{description}

					\begin{tabular}{p{0.5in}p{1in}p{0.5in}p{1in}}
					JSP	 	& D3:412--422		& JSR 		& --- \\
					DC	 	& --- 				& HC 		& 2:4--13 \\
					HDDC 	& ---				& TCHC 		& --- \\
					TPJS 	& 47--56			& 	 		&  \\
					\end{tabular}
					% HDDC = woodford's DC diss; JSR = Marquardt's JS Revelations; TCHC = Vogel HC
				
				\begin{description}
					\item[Online Access] \href{https://www.josephsmithpapers.org/paper-summary/letter-to-the-church-circa-february-1834/1}{Here}
					% \item[Analysis] \hyperref[XX]{Here}
				\end{description}
				
				\begin{description}
					\item[Background]
				\end{description}
				
					% It might also be useful to give a two sentence context of the period: e.g., "This speech was given in Nauvoo to a small congregation one month prior to the destruction of the Nauvoo Expositor. These and other tensions may have been on the prophet's mind when he gave the speech."
				
				\begin{description}
					\item[Original Source]
				\end{description}
				
					% brief description of original source material, with reference to JSP or somewhere else for more info
					% Background material: it might be helpful to have a note that says whether a given document was presented as a speech, was written in a journal, etc.
					% Also a comment here about how reliable the source might be, or what shape it is in, if there are large omissions or parts that are hard to understand, and so on.
				
				\begin{description}
					\item[Text]
				\end{description}
				
					\begin{tightcenter}
							\hypertarget{EMS-1834-JS-CIR-FEB-1834-a}%
						THE
							\marginnote{a}
						ELDERS OF THE CHURCH IN KIRTLAND, TO THEIR BRETHREN ABROAD.
					\end{tightcenter}
					
					\begin{tightright}
						\textit{Dear brethren in Christ, and companions in tribulation:}
					\end{tightright}

					WHEN we call to remembrance the ties with which we are bound to those who embrace the everlasting covenant, and the fellowship and love with which the hearts of the children of our Lord's kingdom should be united, we cherish a belief, that you will bear with us, when we take this course to communicate to you some of the many thoughts which occupy our minds, and press with continued weight upon our hearts, as we reflect upon the vast importance and responsibility of your callings, in the sight of the Master of the vineyard.
						\hypertarget{EMS-1834-JS-CIR-FEB-1834-b}%
					And
						\marginnote{b}%
					though our communications to you may be frequent, yet we believe they will be received on your part with brotherly feelings; and, that from us your unworthy brethren, you will suffer a word of exhortation to have place in your hearts, as you see the great extent of the power and dominion of the prince of darkness, and realize how vast the numbers are who are crowding the road to death without ever giving heed to the cheering sound of the gospel of our Lord Jesus Christ!

						\hypertarget{EMS-1834-JS-CIR-FEB-1834-c}%
					Consider
						\marginnote{c}%
					for a moment, brethren, the fufillment of the words of the prophet; for we behold that darkness covers the earth, and gross darkness the minds of the inhabitants thereof---that crimes of every discription are increasing among men---vices of every enormity are practiced---the rising generation growing up in the fulness of pride and arrogance---the aged losing every sense of conviction, and seemingly banishing every thought of a day of retribution---intemperence, immorality, extravigance, pride, blindness of heart, idolatry, the loss of natural affection, the love of this world, and indifference toward the things of eternity increasing among those who profess a belief in the religion of heaven, and infidelity spreading itself in consequence of the same---%
						\hypertarget{EMS-1834-JS-CIR-FEB-1834-d}%
					men
						\marginnote{d}%
					giving themselves up to commit acts of the foulest kind, and deeds of the blackest dye; lying, blaspheming, stealing, robbing, murdering, defaming, defrauding, blasting the reputation of neighbors, advocating error and opposing the truth, forsaking the covenant of heaven, and denying the faith of Jesus---and in the midst of all this, the day of the Lord fast approaching when none except those who have on the wedding garment will be permitted to eat and drink in the presence of the Bridegroom, the Prince of peace!

						\hypertarget{EMS-1834-JS-CIR-FEB-1834-e}%
					Impressed
						\marginnote{e}%
					with the truth of these facts, what can be the feelings of those who have been made partakers of the heavenly gift, and have tasted the good word of God, and the powers of the world to come? Who but those who can see the awful precipice upon which the world of mankind stand in this generation, can labor in the vineyard of the Lord with a feeling sense of their deplorable situation?
						\hypertarget{EMS-1834-JS-CIR-FEB-1834-f}%
					Who
						\marginnote{f}%
					but those who have duly considered the condesention of the Father of our spirits, in providing a sacrifice for his creatures, a plan of redemption, a power of atonement, a scheme of salvation, having as one of its great objects, to bring men back into the presence of the King of heaven, crown them in the celestial glory, and make them heirs with his Son to that inheritance which is incorruptible, undefiled, and which fadeth not away---can realize the importance of a perfect walk before all men, and a diligence in calling upon all men to partake of these blessings? How indescribably glorious are these tidings to mankind!
						\hypertarget{EMS-1834-JS-CIR-FEB-1834-g}%
					Of
						\marginnote{g}%
					a truth they may be considered tidings of great joy to all people; and tidings too that ought to fill the earth and cheer the heart of every one when sounded in his ears.--- And the reflection, that every one is to receive according to his own diligence and perseverance while in the vineyard, ought to inspire every one who is called to be a minister of these glad tidings, to so improve upon their talent that they may gain other talents, that when the Master sits down to take an account of the conduct of his servants, that it may be said, Well done, good and faithful servant: thou hast been faithful over a few things; I will now make thee ruler over many things: enter thou into the joy of thy Lord.

						\hypertarget{EMS-1834-JS-CIR-FEB-1834-h}%
					Some
						\marginnote{h}%
					may presume to say, that the world in this age is fast increasing in righteousness; that the dark ages of superstition and blindness have passed over, when the faith of Christ was known and practiced only by a few, when ecclesiastic power held an almost universal control over christendom, and when the consciences of men were held bound by the strong chains of priestly power;
						\hypertarget{EMS-1834-JS-CIR-FEB-1834-i}%
					but
						\marginnote{i}%
					now, the gloomy cloud is burst, and the gospel is shining with all the resplendent glory of an apostolic day; and that the kingdom of the Messiah is greatly spreading, that the gospel of our Lord is carried to divers nations of the earth, the scriptures translating into different tongues, the ministers of truth crossing the vast deep to proclaim to men in darkness a risen Savior, and to erect the standard of Emmanuel where light has never shone, and that the idol is destroyed, the temple of images forsaken;
						\hypertarget{EMS-1834-JS-CIR-FEB-1834-j}%
					and
						\marginnote{j}%
					those who but a short time previous followed the traditions of their fathers and sacrificed their own flesh to appease the wrath of some imaginary god, are now raising their voices in the worship of the Most High, and are lifting their thoughts up to him with the full expectation, that one day they will meet with a joyful reception into his everlasting kingdom!

						\hypertarget{EMS-1834-JS-CIR-FEB-1834-k}%
					But,
						\marginnote{k}%
					a moment's candid reflection upon the principles of these systems, the manner they are conducted, the individuals employed, the apparent object held out as an inducement to cause them to act, we think, is sufficient for every candid man to draw a conclusion in his own bosom, whether this is the order of heaven or not. We deem it a just principle, and it is one the force of which we believe ought to be duly considered by every individual, that all men are created \textit{equal}, and that all have the privilege of thinking for themselves upon all matters relative to conscience.
						\hypertarget{EMS-1834-JS-CIR-FEB-1834-l}%
					Consequently,
						\marginnote{l}%
					then, we are not disposed, had we the power, to deprive any one from exercising that free independence of mind which heaven has so graciously bestowed upon the human family as one of its choicest gifts; but we take the liberty, (and this we have a right to do,) of looking at this order of things a few moments, and contrasting it with the order of God as we find it in the sacred scriptures.
						\hypertarget{EMS-1834-JS-CIR-FEB-1834-m}%
					In
						\marginnote{m}%
					this review, however, we shall present the points as we consider they were really designed by the great Giver to be understood, and the happy result arising from a performance of the requirements of heaven, as therein revealed, to every one who obeys them; and the consequence attending a false construction, a misrepresentation, or a forced meaning that was never designed in the mind of the Lord when he condescended to speak from the heavens to men for their salvation.

						\hypertarget{EMS-1834-JS-CIR-FEB-1834-n}%
					Previous
						\marginnote{n}%
					to entering upon a subject of so great a moment to the human family, there is a prominent item which suggests itself to our minds which, here, in few words we wish to discuss: All regularly organized and well established governments, have certain laws by which, more or less, the innocent are protected and the guilty punished. The fact admitted, that certain laws are good, equitable and just, ought to be binding upon the individual who admits this fact, to observe in the strictest manner an obedience to those laws.
						\hypertarget{EMS-1834-JS-CIR-FEB-1834-o}%
					These
						\marginnote{o}%
					laws when violated, or broken by that individual, must, in justice convict his mind with a double force, if possible, of the extent and magnitude of his crime; because he could have no plea of ignorance to produce; and his act of transgression was openly committed againts light and knowledge. But the individual who may be ignorant, and imperceptibly transgresses or violates these laws, though the voice of the country requires that he should suffer, yet he will never feel that remorse of conscience that the other will, and that keen-cutting reflection will never rise in his brest that otherwise would, had he done the deed, or committed the offence in full conviction that he was breaking the law of his country, and having previously acknowledged the same to be just.
						\hypertarget{EMS-1834-JS-CIR-FEB-1834-p}%
					It
						\marginnote{p}%
					is not our intention by these remarks, to attempt to place the law of man on a parallel with the law of heaven; because we do not consider that it is formed in that wisdom and propriety; neither do we consider that it is sufficient in itself to bestow any thing in comparison to the law of heaven, even should it promise it. The law of men may guarantee to a people protection in the honorable pursuits of this life, and the temporal happiness arising from a protection against unjust insults and injuries; and when this is said, all is said, that can be in truth, of the power, extent, and influence of the law of men, exclusive of the law of God.
						\hypertarget{EMS-1834-JS-CIR-FEB-1834-q}%
					The
						\marginnote{q}%
					law of heaven is presented to man, and as such guarantees to all who obey it a reward far beyond any earthly consideration: it does not promise that the believer in every age should be exempt from the afflictions and troubles arising from different sources in consequence of wicked men on earth; though in the midst of all this there is a promise predicated upon the fact that it is the law of heaven, which transcends the law of man, as far as eternal life is prefferable to temporal; and the blessings which God is able to give, greater than those which can be given by man!
						\hypertarget{EMS-1834-JS-CIR-FEB-1834-r}%
					Then,
						\marginnote{r}%
					certainly, if the law of man is binding upon man when acknowledged, much more must the law of heaven be. And as much as the law of heaven is perfect, more than the law of man, so much greater must be the reward if obeyed. The law of man promises safety in temporal life; but the law of God promises that life which is eternal, even an inheritance at his own right hand, secure from all the powers of the wicked one.

						\hypertarget{EMS-1834-JS-CIR-FEB-1834-s}%
					We
						\marginnote{s}%
					consider that God has created man with a mind capable of instruction, and a faculty which may be enlarged in proportion to the heed and diligence given to the light communicated from heaven to the intellect; and that the nearer man approaches perfection, the more conspicuous are his views, \& the greater his enjoyments, until he has overcome the evils of this life and lost every desire of sin; and like the ancients, arrives to that point of faith that he is wrapped in the glory and power of his Maker and is caught up to dwell with him.
						\hypertarget{EMS-1834-JS-CIR-FEB-1834-t}%
					But
						\marginnote{t}%
					we consider that this is a station to which no man ever arrived in a moment: he must have been instructed into the government and laws of that kingdom by proper degrees, till his mind was capable in some measure of comprehending the propriety, justice, equity, and consistency of the same. For further instruction we refer you to Deut. xxxii. where the Lord says, that Jacob is the lot of his inheritance. He found him in a desert land, and in the waste howling wilderness; he led him about, he instructed him, he kept him as the apple of his eye, \&c. which will show the force of the last item advanced, that it is necessary for men to receive an understanding concerning the laws of the heavenly kingdom, before they are permitted to enter it: we mean the celestial glory.
						\hypertarget{EMS-1834-JS-CIR-FEB-1834-u}%
					So
						\marginnote{u}%
					dissimilar are the governments of men, and so divers are their laws, from the government and laws of heaven, that a man, for instance, hearing that there was a country on this globe called the United States of North America, could take his journey to this place without first learning the laws of this government; but the conditions of God's kingdom are such, that all who are made partakers of that glory, are under the necessity of first learning something respe[c]ting it previous to their entering into it. But the foreignor can come to this country without knowing a syllable of its laws, or even subscribing to obey them after he arrives. Why? Because the government of the United State does not require it: it only requires an obedience to its laws after the individual has arrived within its jurisdiction.

						\hypertarget{EMS-1834-JS-CIR-FEB-1834-v}%
					As
						\marginnote{v}%
					we previously remarked, we do not attempt to place the law of man on a parallel with the law of heaven; but we will bring forward another item, to further urge the propriety of yielding obedience to the law of heaven, after the fact is admitted, that the laws of man are binding upon man. Were a king to extend his dominion over the habitable earth, and send forth his laws which were of the most perfect kind, and command his subjects one and all to yield obedience to the same;
						\hypertarget{EMS-1834-JS-CIR-FEB-1834-w}%
					and
						\marginnote{w}%
					annex as a reward to those who obeyed them, that at a certain period they should be called to attend the marriage of his son, who in due time was to receive the kingdom, and they should be made equal with him in the same; and annex as a penalty for disobedience that every individual should be cast out at the marriage feast, and have no part nor portion with his government; and what rational mind could for a moment accuse the king with injustice for punishing such rebellious subjects?
						\hypertarget{EMS-1834-JS-CIR-FEB-1834-x}%
					In
						\marginnote{x}%
					the first place his laws were just, easy and perfect: nothing was required in them of a tyranical nature; but their very construction was equity and beauty; and when obeyed would produce the happiest situation possible to all who adheard to them, beside the last great benefit of sitting down with a royal robe in the presence of the king at the great grand marriage supper of his son, and be made equal with him in all the affairs of the kingdom.

						\hypertarget{EMS-1834-JS-CIR-FEB-1834-y}%
					When
						\marginnote{y}%
					these royal laws were issued, and promulgated throughout the vast dominion, every subject, when interrogated whether he believed them to be from his sovereign answered, Yes, I know they are, I am acquainted with the signature, for it is as usual, \textit{THUS SAITH THE KING!} This admitted, the subject is bound by every consideration of honor to his country, his king, and his own personal character, to observe in the strictest sense every requisition in the royal edict. Should any escape the search of the embassadors of the king, and never hear these last laws, giving his subjects such exalted privileges, an excuse might be urged in their behalf, and they escape the censure of the king.
						\hypertarget{EMS-1834-JS-CIR-FEB-1834-z}%
					But
						\marginnote{z}%
					for those who had heard, who had admitted, and who had promised obedience to these just laws no excuse could be urged, and when brought into the presence of the king, certainly, justice would require that they should suffer a penalty! Could that king be just in admitting these rebellious individuals into the full enjoyment and privileges with his son, and those who had been obedient to his commandments? Certainly not. Because they disregarded the voice of their lawful king; they had no regard for his virtuous laws, for his dignity, nor for the honor of his name; neither for their own country's sake, nor their private virtue!
						\hypertarget{EMS-1834-JS-CIR-FEB-1834-aa}%
					They
						\marginnote{aa}%
					neither regarded his authority enough to obey him, neither did they regard the immediate advantages and blessings arising from these laws if kept, to observe them, so destitute were they of virtue and goodness; and above all, they regarded so little the joy and satisfaction of a legal seat in the presence of the king's only son, and to be made equal with him in all the blessings, honors, comforts, and felicities of his kingdom, that they turned away from an anticipation of them, and considered that they were beneath their present notice, though they had no doubt as to the real authenticity of the royal edict.
		
						\hypertarget{EMS-1834-JS-CIR-FEB-1834-bb}%
					We
						\marginnote{bb}%
					ask, again, would the king be just in admitting these rebels to all the privileges of his kingdom, with those who had served him with the strictest integrity? We again answer, No! such individuals would be dangerous characters in any government, good \& wholesome laws they dispised; just and perfect principles they trampled under their feet as something beneath their notice, and disregarded those commands of their sovereign entirely which they had once acknowledged to be equitable!
						\hypertarget{EMS-1834-JS-CIR-FEB-1834-cc}%
					How
						\marginnote{cc}%
					could a government be conducted with harmony if its administrators were possessed with such different dispositions and different principles? Could it prosper? Could it flourish? Would harmony prevail? Would order be established, and could justice be executed in righteousness in all branches of its department? No! In it were two classes of men as discimilar as light is from darkness, virtue from vice, justice from injustice, truth from falsehood, and holiness from sin! One class were perfectly harmless and virtuous: they knew what virtue was for they had lived in the fullest enjoyment of it, and their fidelity to truth fairly tested by a series of years of faithful obedience to all its heavenly precepts.
						\hypertarget{EMS-1834-JS-CIR-FEB-1834-dd}%
					They
						\marginnote{dd}%
					knew what good order was, for they had been orderly and obedient to the laws imposed on them by their wise sovereign, and had experienced the benefits arising from a life spent in his government till he had now seen proper to make them equal with his son.--- Such individuals would indeed adorn any court where perfection was one of its main springs of action, and shine far more brilliant than the richest gem in the diadem of the prince.

						\hypertarget{EMS-1834-JS-CIR-FEB-1834-ee}%
					The
						\marginnote{ee}%
					other class were a set of individuals who disregarded every principle of justice and equity, whatever: and this is demonstrated from the fact, that when just laws were issued by the king, which were perfectly equitable, they were so lost to a sense of righteousness that they disregarded those laws, notwithstanding an obedience to them would have produced the happiest result possible, at the time, as regarded their own personal comfort and advantage. They were entirely destitute of harmony and virtue, so much so that virtuous laws they dispised. They had proven themselves unworthy a place in the joys of the prince, because they had for a series of years lived in open violation of his government.
						\hypertarget{EMS-1834-JS-CIR-FEB-1834-ff}%
					Certainly,
						\marginnote{ff}%
					then, those two clases of men could not hold the reins of the same government at the same time in peace; for internal jars, broils, and discords would rack it to the center, were such a form of government to attempt to exist under such a system. The virtuous could not enjoy peace in the constant and unceasing schemes and evil plans of the wicked; neither could the wicked have enjoyment in the constant perseverance of the righteous to do justly. And that there must be an agreement in this government, or it could not stand, must be admitted by all.
						\hypertarget{EMS-1834-JS-CIR-FEB-1834-gg}%
					Should
						\marginnote{gg}%
					the king convey the reins into the hands of the rebellious the government must soon fall; for every government, from the creation to the present; when it ceased to be virtuous, and failed to execute justice, sooner or later has been overthrown. And without virtuous principles to actuate a government all care for justice is soon lost, and the only motive which prompts it to act is, ambition and selfishness. Should the king admit these rebels into his house to make them equal with the others, would be condescending beneath his character;
						\hypertarget{EMS-1834-JS-CIR-FEB-1834-hh}%
					because
						\marginnote{hh}%
					he once issued virtuous laws which were received by a part of his subjects, and the reward annexed was a seat at the marriage feast, and an adoption into his own family as lawful heirs. So should he now offer any thing differently he would blast forever his own reputation, and destroy forever that government which he once so diligently labored to establish and preserve, and which he once had wisdom to organize. Such individuals as the last named, would be a bane to a virtuous government, and would prove its overthrow if suffered to hold a part in conducting its h[e]lm!

						\hypertarget{EMS-1834-JS-CIR-FEB-1834-ii}%
					We
						\marginnote{ii}%
					take the sacred writings into our hands. and admit that they were given by direct inspiration for the good of man. We believe that God condescended to speak from the heavens and declare his will concerning the human family: give to them just and holy laws to regulate their conduct, and guide them in a direct way, that in due time he might take them to himself, and make them joint heirs with his Son. But when this fact is admitted, that the immediate will of heaven is here contained, are we not bound as rational creatures to live in accordance to all its precepts? Will the mere admision, that this is the will of heaven ever benefit us if we do not comply with all its teachings?
						\hypertarget{EMS-1834-JS-CIR-FEB-1834-jj}%
					Do
						\marginnote{jj}%
					we not offer violence to the Supreme Intelligence of heaven, when we admit the truth of its teachings, and do not obey them? Do we not condescend beneath our own character, and the better wisdom which heaven has endowed us with, by such a course of conduct? For these reasons, if we have direct revelations given us from heaven, surely, those revelations were never given to be trifled with, without the triflers incuring displeasure, and assuring vengeance upon their own heads, if there is any justice in heaven; and that there is, must be admitted by every individual who admits the truth and force of its teachings; its blessings and cursings, as contained in the sacred volume.

						\hypertarget{EMS-1834-JS-CIR-FEB-1834-kk}%
					Here,
						\marginnote{kk}%
					then, we have this part of our subject immediately before us for consideration: God has in reserve a time, or period appointed in his own bosom, when he will bring all his subjects, who have obeyed his voice and kept his commandments, into his celestial rest. This rest is of such perfection and glory, that man has need of a preparation before he can, according to the laws of that kingdom enter it and enjoy its blessings.--- This being the fact, God has given certain laws to the human family, which, if observed, are sufficient to prepare them to inherit this rest.
						\hypertarget{EMS-1834-JS-CIR-FEB-1834-ll}%
					This,
						\marginnote{ll}%
					then, we conclude, was the purpose of God in giving his laws to us: if not, why, or for what were they given? If the whole family of man were as well off without them as they might be with them, for what purpose or intent were they ever given? Was it that God wanted to merely show that he could talk? This would be nonsense, to suppose that he would condescend to talk in vain; for it would be in vain, and to no purpose whatever:
						\hypertarget{EMS-1834-JS-CIR-FEB-1834-mm}%
					because,
						\marginnote{mm}%
					all the commandments contained in the law of the Lord, have the sure promise annexed of a reward to all who obey; predicated upon the fact, that they are really the promises of a Being who cannot lie, and who is abundantly able to fulfil every tittle of his word: and if men were as well prepared, or could be as well prepared, to meet God without their ever having been given in the first instance, why were they ever given? for certainly, in that case they can now do him no good.

						\hypertarget{EMS-1834-JS-CIR-FEB-1834-nn}%
					As
						\marginnote{nn}%
					we previously remarked, all well established and properly organized governments have certain fixed and prominent laws for the regulation and management of the same.--- If man has grown to wisdom and is capable of discerning the propriety of laws to govern nations, what less can we expect from the Ruler and Upholder of the universe? Can we suppose that he has a kingdom without laws? Or do we believe that it is composed of an innumerable company of beings who are entirely beyond all law?
						\hypertarget{EMS-1834-JS-CIR-FEB-1834-oo}%
					Consequently
						\marginnote{oo}%
					have need of nothing to govern or regulate them? Would not such ideas be reproachful to our Great Parent, and an attempt to cast a stigma upon his glorious character! Would it not be asserting, that we had found out a secret beyond Deity? that we had learned that it was good to have laws, and yet He, after existing from eternity, and having power to create man, had not found out the fact, that it was proper to have laws for his government!
						\hypertarget{EMS-1834-JS-CIR-FEB-1834-pp}%
					We
						\marginnote{pp}%
					admit that God is the great source and fountain from whence proceeds all good; that he is perfect intelligence, and that \textit{his} wisdom is alone sufficient to govern and regulate the mighty creations and worlds which shine and blaze with such magnificence and splendor over our heads, as though touched with his finger and moved by his Almighty word. And if so, it is done and regulated by law; for without law all must certainly fall into chaos. If, then, we admit that God is the source of all wisdom and understanding, we must admit that by his direct inspiration he has taught man that law was necessary in order to govern and regulate his own immediate interest and welfare:
						\hypertarget{EMS-1834-JS-CIR-FEB-1834-qq}%
					For
						\marginnote{qq}%
					this reason, it is beneficial to promote peace and happiness among men: And as before remarked, God is the source from whence proceeds all good; and if man is benefitted by law, then certainly, law is good; and if law is good, it, or the principle of it emanated from God: for God is the source of all good; consequently, then, he was the first Author of law, or the principle of it, to mankind.

					\textsc{To be continued.}
					
		% -----
		\subsubsection{Letter to J. G. Fosdick, 3 February 1834}
		% -----
			\label{EMS-1834-JS-3-FEB-1834}
			% \label{EMS-1830-JS-DAY-3LETTERMONTH-YEAR}
			
			\begin{description}
					\item[Print Sources]
				\end{description}

					\begin{tabular}{p{0.5in}p{1in}p{0.5in}p{1in}}
					JSP	 	& D3:422--425		& JSR 		& --- \\
					DC	 	& --- 		& HC 		& --- \\
					HDDC 	& ---		& TCHC 		& --- \\
					\end{tabular}
					% HDDC = woodford's DC diss; JSR = Marquardt's JS Revelations; TCHC = Vogel HC
				
				\begin{description}
					\item[Online Access] \href{https://www.josephsmithpapers.org/paper-summary/letter-to-j-g-fosdick-3-february-1834/1}{Here}
					% \item[Analysis] \hyperref[XX]{Here}
				\end{description}
				
				\begin{description}
					\item[Background]
				\end{description}
				
					% It might also be useful to give a two sentence context of the period: e.g., "This speech was given in Nauvoo to a small congregation one month prior to the destruction of the Nauvoo Expositor. These and other tensions may have been on the prophet's mind when he gave the speech."
				
				\begin{description}
					\item[Original Source]
				\end{description}
				
					% brief description of original source material, with reference to JSP or somewhere else for more info
					% Background material: it might be helpful to have a note that says whether a given document was presented as a speech, was written in a journal, etc.
					% Also a comment here about how reliable the source might be, or what shape it is in, if there are large omissions or parts that are hard to understand, and so on.
				
				\begin{description}
					\item[Text]
				\end{description}
				
					\begin{tightright}
							\hypertarget{EMS-1834-JS-3-FEB-1834-a}%
						Kirtland
							\marginnote{a}
						Mills Feb. 3, 1834.
					\end{tightright}
				
					Dear Bro. [J. G.] Fosdick:
					
					Your letter of the 10th. Jan. last is just recd. and this day there has been a regular Council of H. Priests \& Elders in this place, and the subject spoken of in your letter was, we believe, taken into due consideration. We were very sorry to learn that Bro. J[oseph] Wood had gone so far astray and offered such violence to the pure principles of the Gospel of Christ. But, alas! Such is the depravity of man when lost to a sense of the fear of God and of the ties which bind [e]very virtuous man to the interest and happiness of his follow man.
					
						\hypertarget{EMS-1834-JS-3-FEB-1834-b}%
					Every
						\marginnote{b}%
					principle inculcated among you which is contrary [t]o \uline{virtue}, to industry, to wisdom, to good order, to propriety, and in fine, to the pure principles of godliness as contained in the Scriptures of the old and new Testaments, the Book of Mormon and the revelations and commandments of Jesus Christ, which have been given to his Church in these last days, is entirely foreign from the feelings of our breasts, and is that upon which we look down with feelings of the utmost disapprobation; and as consc[i]encious men who expect to render an impartial account, before [th]e searcher of hearts, of all our transactions here, we cannot [lo]ok upon any principle contrary to the above with any degree [of] allowance.
					
						\hypertarget{EMS-1834-JS-3-FEB-1834-c}%
					After
						\marginnote{c}%
					some investigation of the case of Bro. Wood, in Council, [it] was decided that he should be cut off from the Church. [Ac]cordingly the Council lifted their hands against him and [he] was excluded from the church on this 3d. day of Feb. 1834. [for] indulging an \uline{idle}, \uline{partial}, \uline{overbearing} and \uline{lustful} spirit, and [not] magnifying his holy calling whereunto he had been [ord]ained. These things were plainly manifest to the satisfaction [of] [a]ll the council, and the Spirit constrained us to separate him [fro]m the church. Should bro. Joseph Wood, after learning [th]e decission of this council, truly repent of all his sins and bring forth fruit meet to the satisfaction of that branch of the Church where he has committed the offences, he can be rebaptized and come into the church again if he desire so to do. The instructions which you desire relative to church gov. \&c. the extent of the power of a high Priest over any branch of the church, are subjects which will be investigated in the next no. of the Star. Time will not allow us to write the subjects at full length now;
						\hypertarget{EMS-1834-JS-3-FEB-1834-c}%
					suffice it,
						\marginnote{c}%
					therefore, to say that there is no office in this Church which can be placed upon the head of any man that will place him beyond the power \& control of any branch of the church where he may be guilty of transgression, even if there is not another ordained member in the church, let the church appoint some brother to preside and let them do as one church did in ancient days ``\uline{try} \uline{them} \uline{who} \uline{say} \uline{they} \uline{are} \uline{apostles} \uline{and} \uline{are} \uline{not}, \uline{but} \uline{are} \uline{liars},'' then let them demand their license, raise their hands against them and thus they are expelled from the communion of the church.
						\hypertarget{EMS-1834-JS-3-FEB-1834-d}%
					It
						\marginnote{d}%
					requires all the members of the church to constitute the body of Christ. One man is not the body, nor are the children of <the> Kingdom to be tantalized by men who may hold licenses and have authority to preach the gospel; such have the more need to be discreet and humble. Should the individual, after being thus dealt with be dissatisfied with the decission of the church he can appeal to a Bishop's Court, and should he there be judged guilty, he can yet appeal to a court of high pri[ests] and this is an end of all disputes and controversies in the Church of God on earth.
					
					\begin{tightright}
					Brethren in the new covenant, Farewell.
					\end{tightright}
					
					\begin{tabular}{p{1in}p{2in}p{2in}}
					Signed & Joseph Smith Jr. & (Moderator \\
					 & Orson Hyde & (Clerk of Coun[cil]\\
					\end{tabular}
					
					Copied by direction of the Council, by Oliver Cowdery
					
		% -----
		\subsubsection{Minutes, 9 February 1834}
		% -----
			\label{EMS-1834-JS-9-FEB-1834}
			% \label{EMS-1830-JS-DAY-3LETTERMONTH-YEAR}
			
			\begin{description}
					\item[Print Sources]
				\end{description}

					\begin{tabular}{p{0.5in}p{1in}p{0.5in}p{1in}}
					JSP	 	& D3:425--427		& JSR 		& --- \\
					DC	 	& --- 		& HC 		& 2:24--25 \\
					HDDC 	& ---		& TCHC 		& 2:12--13 \\
					\end{tabular}
					% HDDC = woodford's DC diss; JSR = Marquardt's JS Revelations; TCHC = Vogel HC
				
				\begin{description}
					\item[Online Access] \href{https://www.josephsmithpapers.org/paper-summary/minutes-9-february-1834/1}{Here}
					% \item[Analysis] \hyperref[XX]{Here}
				\end{description}
				
				\begin{description}
					\item[Background]
				\end{description}
				
					% It might also be useful to give a two sentence context of the period: e.g., "This speech was given in Nauvoo to a small congregation one month prior to the destruction of the Nauvoo Expositor. These and other tensions may have been on the prophet's mind when he gave the speech."
				
				\begin{description}
					\item[Original Source]
				\end{description}
				
					% brief description of original source material, with reference to JSP or somewhere else for more info
					% Background material: it might be helpful to have a note that says whether a given document was presented as a speech, was written in a journal, etc.
					% Also a comment here about how reliable the source might be, or what shape it is in, if there are large omissions or parts that are hard to understand, and so on.
				
				\begin{description}
					\item[Text]
				\end{description}
				
					...The work of the building of the House of the Lord in Kirtland was also taken into Consid[e]ration it was decided that the brethren in this place should assist in Erecting the house all that is in their power, that the Elders of the Church may be endowed with power from on high according to the promise of God, that the work of the father may roll forth...
					
		% -----
		\subsubsection{Minutes, 12 February 1834}
		% -----
			\label{EMS-1834-JS-12-FEB-1834}
			% \label{EMS-1830-JS-DAY-3LETTERMONTH-YEAR}
			
			\begin{description}
					\item[Print Sources]
				\end{description}

					\begin{tabular}{p{0.5in}p{1in}p{0.5in}p{1in}}
					JSP	 	& D3:427--431		& JSR 		& --- \\
					DC	 	& --- 		& HC 		& 2:25--27 \\
					HDDC 	& ---		& TCHC 		& 2:13--14 \\
					TPJS 	& 69--70		& 	 		&  \\
					\end{tabular}
					% HDDC = woodford's DC diss; JSR = Marquardt's JS Revelations; TCHC = Vogel HC
				
				\begin{description}
					\item[Online Access] \href{https://www.josephsmithpapers.org/paper-summary/minutes-12-february-1834/1}{Here}
					% \item[Analysis] \hyperref[XX]{Here}
				\end{description}
				
				\begin{description}
					\item[Background]
				\end{description}
				
					% It might also be useful to give a two sentence context of the period: e.g., "This speech was given in Nauvoo to a small congregation one month prior to the destruction of the Nauvoo Expositor. These and other tensions may have been on the prophet's mind when he gave the speech."
				
				\begin{description}
					\item[Original Source]
				\end{description}
				
					% brief description of original source material, with reference to JSP or somewhere else for more info
					% Background material: it might be helpful to have a note that says whether a given document was presented as a speech, was written in a journal, etc.
					% Also a comment here about how reliable the source might be, or what shape it is in, if there are large omissions or parts that are hard to understand, and so on.
				
				\begin{description}
					\item[Text]
				\end{description}
				
						\hypertarget{EMS-1834-JS-12-FEB-1834-a}%
					Thursday
						\marginnote{a}%
					Evening, February 12. 1834. This evening the high Priests and Elders of the Church in Kirtland at the house of bro. Joseph Smith Jun. in Council for Church business. The Council was \sout{opened} organized, and opened by bro. Joseph Smith Ju[nior] in prayer. Bro. Joseph then rose and said: I shall now endeavour to set forth before this Council, the dignity of the office which has been conferred upon me by the ministering of the Angel of God, by his own voice and by the voice of this Church. I have never set before any council in all the order in which a Council ought to be conducted, which, prehaps, has deprived the Councils of some, or many blessings.
					
						\hypertarget{EMS-1834-JS-12-FEB-1834-b}%
					He
						\marginnote{b}%
					said, that no man was capable of judging a matter in council without his own heart was pure; and that we frequently, are so filled with prejudice, or have a beam in our own eye, that we are not capable of passing right descissions, \&c.
					
					But to return to the subject of the order: In ancient days councils were conducted with such strict propriety, that no one was allowed to whisper, be weary, leave the room, or get uneasy in the least, until the voice of the Lord, by revelation, or by the voice of the Council by the spirit was obtained: which has not been observed in this church to the present. It was understood in ancient days, that if one man could stay in Council another Could, and if the president could spend his time, the members could also.
						\hypertarget{EMS-1834-JS-12-FEB-1834-c}%
					But
						\marginnote{c}%
					in our Councils, generally, one would be uneasy, another asleep, one praying another not; one's mind on the business of the Council and another thinking on something else \&c. Our acts are recorded, and at a future day they will be laid before us, and if we should fail to judge right and injure our fellow beings, they may there prehaps condemn us; then, they are of great Consequence: and to me the Consequence appears to be of force beyond any thing which I am able to express \&c. Ask yourselves, brethrn, how much you have exercised yourselves in prayer since you heard of this Council; and if you are now prepared to sit in judgment upon the soul of your brother.---
						\hypertarget{EMS-1834-JS-12-FEB-1834-d}%
					Bro
						\marginnote{d}%
					Joseph then went on to give us a relation of his situation at the time he obtained the record, the persecution he met with \&C. He also told us of his transgressing at the time he was translateing the Book of Mormon. He also prophecied that he should stand and shine like the sun in the firmament when his enemies and the gainsayers of his testimony should be put down and Cut off and their names blotted out from among men. After the Council had rec[e]ived much good instruction from Bro. Joseph. The Case of Bro. Martin Harris against whom certain Charges were preferred by bro. Sidney Rigdon. One was that he told Esqr A[lpheus] C. Russell that Joseph drank too much liquor when he was translating the Book of Mormon and that he wrestled with many men and threw them \&c. Another charge was, that he exalted himself above bro. Joseph, in that he said bro. Joseph knew not the contents of the book of Mormon until it was translated.
						\hypertarget{EMS-1834-JS-12-FEB-1834-e}%
					\phantom{ }
						\marginnote{e}%
					\phantom{ }\sout{Bro.} \sout{Martin} but that he himself knew all about it before it was translated. Bro. Martin said he did not tell Esqr Russell that bro. Joseph drank too much liquor while translateing the book of Mormon, but this thing took place before the book of Mormon was translated. He confessed that his mind was darkend and that he had said many things inadvertently \sout{calculateing} calculateid to wound the feelings of his bretheren and promised to do better. The Council forgave him and gave him much good advice. Bro Rich was Called in question for transgressing the word of wisdom and for selling the revelations at an extortionary price while he was gone East with father Lions which thing Bro. Rich confessed before the Council and the Council forgave him upon his promiseing to do better and reform his life.---
					
					Council then Concluded by prayer by Bro. S. Rigdon
					
					\begin{tightright}
					Orson Hyde Clk
					\end{tightright}
					
		% -----
		\subsubsection{Minutes, 17 February 1834}
		% -----
			\label{EMS-1834-JS-17-FEB-1834}
			% \label{EMS-1830-JS-DAY-3LETTERMONTH-YEAR}
			
			\begin{description}
					\item[Print Sources]
				\end{description}

					\begin{tabular}{p{0.5in}p{1in}p{0.5in}p{1in}}
					JSP	 	& D3:435--439		& JSR 		& --- \\
					DC	 	& --- 		& HC 		& 2:28--31 \\
					HDDC 	& ---		& TCHC 		& 2:17--20 \\
					\end{tabular}
					% HDDC = woodford's DC diss; JSR = Marquardt's JS Revelations; TCHC = Vogel HC
				
				\begin{description}
					\item[Online Access] \href{https://www.josephsmithpapers.org/paper-summary/minutes-17-february-1834/1}{Here}
					% \item[Analysis] \hyperref[XX]{Here}
				\end{description}
				
				\begin{description}
					\item[Background]
				\end{description}
				
					% It might also be useful to give a two sentence context of the period: e.g., "This speech was given in Nauvoo to a small congregation one month prior to the destruction of the Nauvoo Expositor. These and other tensions may have been on the prophet's mind when he gave the speech."
				
				\begin{description}
					\item[Original Source]
				\end{description}
				
					% brief description of original source material, with reference to JSP or somewhere else for more info
					% Background material: it might be helpful to have a note that says whether a given document was presented as a speech, was written in a journal, etc.
					% Also a comment here about how reliable the source might be, or what shape it is in, if there are large omissions or parts that are hard to understand, and so on.
				
				\begin{description}
					\item[Text]
				\end{description}
				
						\hypertarget{EMS-1834-JS-17-FEB-1834-a}%
					This
						\marginnote{a}%
					day, Feb. 17 1834, a Conference of High Priests assembled \sout{at} in Kirtla[n]d at the House of bro. Joseph Smith Ju\supers{r}. They proceeded to organize the Presidents Church Council, Consisting of twelve high priests, and this according to the law of God. The names of those who were chosen \sout{as} \sout{Counsellors,} were Joseph Smith Jun\supers{r.} Sidney Rigdon and F[rederick] G. Williams Presidents. Joseph Smith Seigr [Sr.], John Smith, Joseph Coe, John Johnson, Martin Harris, John S Carter, Jared Corter [Carter], Oliver Cowdery, Sam\supers{l.} H. Smith, Orson Hyde, Sylvester Smith, and Luke Johnson, Counsellors.
						\hypertarget{EMS-1834-JS-17-FEB-1834-b}%
					Bro.
						\marginnote{b}%
					Joseph opened the Council by solem prayer. He then arose and called upon the high priests, Elders, priests, teachers and deacons that were present who had not been nominated as Counsellors to pass their vote whether they were satisfied with the appointment or nomination of the twelve to Compose the Church Council. It was the unanimous voice of all present that those who had been nominated, as above, should compose a standing Council in Kirtland. It was also voted that when any one <or more> of the standing Counsellors \sout{was} were absent, their vacancy should be filled by any high priests whom the majority of the Council should nominate or choose,
						
						\hypertarget{EMS-1834-JS-17-FEB-1834-c}%
					Provideing
						\marginnote{c}%
					that no Council shall be held unless seven of the above named Counsellors are present, or their successors. The above named Counsellors all manifested a willingness to act according to their appointment, the Lord being their helper. Bro Hyrum Smith acted in the place of John Smith. There were nine high priests present and acted in the appointment of the above named Counsellors, also seventeen Elders, and four priests with thirteen private members. Bro Joseph then said he would show the order of Councils in ancient days (see 27 \& 28 pages) as shown to him by vision.
						\hypertarget{EMS-1834-JS-17-FEB-1834-d}%
					The
						\marginnote{d}%
					law \sout{and} by which to govern the Council in the Church of Christ. Jerusalem was the seat of the Church Council in ancient days. The apostle, Peter, was the president of the Council \sout{in ancient days} and held the Keys of the Kingdom of God, <on the Earth> was appointed to this office by the voice of the Saviour and \sout{confirmed} <\sout{acknowledgement} acknowledged> in it by the voice of the Church. He had two men appointed as Counsellors with him, and in case Peter was absent, his Counsellors Could transact business. <or either one of them. The President could also transact business alone.>
						\hypertarget{EMS-1834-JS-17-FEB-1834-e}%
					It
						\marginnote{e}%
					was not the order of heaven in ancient Councils to plead for and against the guilty as in our judicial Courts (so called) but that \sout{if} every Counsellor when he arose to speak, should speak precisely according to evidence and according to the teaching of the spirit of the Lord, that no Counsellor should attempt to screen the guilty when his guilt was manifest
						\hypertarget{EMS-1834-JS-17-FEB-1834-f}%
					That
						\marginnote{f}%
					the person acused before the high council had a right to one half the memb[e]rs of the council to plead his cause, \sout{that is six,} in order that his case might be fairly presented before the President that a descission might be rendered according to truth and righteousness. If the case was not a very difficult one to investigate, two of the Counsellors only, spoke, one \sout{for the accused and one against} <on one side and one on the other> according to evidence. If the case was more difficult, according to the judgment of the Council, two were to speak on each side, and if more difficult, three might speak on each side, and three only. Those who spoke in Council were chosen by the council and that too by casting lots. Those who were thus chosen to speak, took their regular turn, in speaking.
						\hypertarget{EMS-1834-JS-17-FEB-1834-g}%
					Bro
						\marginnote{g}%
					Joseph said that this organization was an ensample to the high priests in their Councils abroad, and a copy of their proceedings be transmitted to the seat of the goverment of the Church to be recorded on the general record. In all cases, the accuser and the acused have a perfect right to speak for themselves before the Council. The Councils abroad, have a right and it is their duty to appoint a president for the time being for themselves. If in case the parties are not satisfied with the decission of the Council abroad, they have a right to an appeal to the Bishops Court, and from thence to the presidents Council which is an\sout{d} end of all strife
					
						\hypertarget{EMS-1834-JS-17-FEB-1834-h}%
					The
						\marginnote{h}%
					remaining six Counsellors who do not speak in Council, are to hear patiently the reasoning of the other and correct all errors which they may discover, and after decission is rendered by the president, if these remaining counsellors can throw any farther light upon the subject, so as to correct the decissin [decision] of the president, they have the liberty so to do, otherwise it stands and the majority of the Council must rule.
						\hypertarget{EMS-1834-JS-17-FEB-1834-i}%
					It
						\marginnote{i}%
					was then voted by all present that they desired to come under the present order of things which they all considered to be the will of God. Many questions have been asked during the time of the organization of this Council and doubtless some errors have been committed, it was, therefore, voted by all present that Bro Joseph should make all necessary corrections by the spirit of inspiration
						\hypertarget{EMS-1834-JS-17-FEB-1834-j}%
					hereafter
						\marginnote{j}%
					Oliver Cowdery drew no. one by lot. Joseph Coe drew No 2. Samuel H Smith drew No 3. Luke Johnson drew No 4. John S Carter drew No 5. Sylvester Smith drew No 6. Oliver Cowdery, Samuel H Smith and John S Carter speak for and on the part of the accuser. Joseph Coe, Luke Johnson and Sylvester Smith, speak for and on the part of the accused. The remaining six counsellors are to sit and hear patiently and correct errors if they discover them. \sout{The} \sout{Council} John Johnson drew No 7. Orson Hyde drew No 8, Jared Carter drew No 9, Joseph Smith Seignr drew No 10, John Smith drew No 11, Martin Harris drew No 12, The Council adjourned then, until wednesday at 10 oclk A.M.------
					
					\begin{tightright}
					Orson Hyde Clk
					\end{tightright}
					
		% -----
		\subsubsection{Revised Minutes, 18--19 February 1834 [D\&C 102]}
		% -----
			\label{EMS-1834-JS-18-19-FEB-1834}
			% \label{EMS-1830-JS-DAY-3LETTERMONTH-YEAR}
			
			\begin{description}
					\item[Print Sources]
				\end{description}

					\begin{tabular}{p{0.5in}p{1in}p{0.5in}p{1in}}
					JSP	 	& D3:439--444		& JSR 		& --- \\
					DC	 	& Section 102 		& HC 		& 2:28--31 \\
					HDDC 	& 1308--1226		& TCHC 		& 2:17--20 \\
					\end{tabular}
					% HDDC = woodford's DC diss; JSR = Marquardt's JS Revelations; TCHC = Vogel HC
				
				\begin{description}
					\item[Online Access] \href{https://www.josephsmithpapers.org/paper-summary/revised-minutes-18-19-february-1834-dc-102/1}{Here}
					% \item[Analysis] \hyperref[XX]{Here}
				\end{description}
				
				\begin{description}
					\item[Background]
				\end{description}
				
					% It might also be useful to give a two sentence context of the period: e.g., "This speech was given in Nauvoo to a small congregation one month prior to the destruction of the Nauvoo Expositor. These and other tensions may have been on the prophet's mind when he gave the speech."
				
				\begin{description}
					\item[Original Source]
				\end{description}
				
					% brief description of original source material, with reference to JSP or somewhere else for more info
					% Background material: it might be helpful to have a note that says whether a given document was presented as a speech, was written in a journal, etc.
					% Also a comment here about how reliable the source might be, or what shape it is in, if there are large omissions or parts that are hard to understand, and so on.
				
				\begin{description}
					\item[Text]
				\end{description}
				
					The above items have been corrected according to the resolution passed <in the same>, and the following is the correction.------

					\begin{tightright}
					Kirtland Feb 17. 1834.
					\end{tightright}

					This day a <general> council of <24> high Priests assembled at the house of Joseph Smith Jun\supers{r.} <by revelation> and proceeded to organize the high council of the Church of Christ, which \sout{is} <was> to consist of twelve high priests, and one, or three presidents, as the case \sout{may} <might> require. This <high> council \sout{is} <was> appointed by revelation, for the purpose of settleing important difficulties which \sout{may} <might> arise in the church, which \sout{cannot} <could not> be settled by the Church, or the bishop's council to the satisfaction of the parties
					
					\begin{tightcenter}
					Joseph Smith Jun\supers{r.}

					Sidney Rigdon and
					
					Frederick G Williams were acknowledged presidents, by the voice of the council; and
					\end{tightcenter}

					\begin{tabular}{p{3in}p{3in}}
					Joseph Smith Seign & Jared Carter \\
					John Smith & Oliver Cowdery \\
					Joseph Coe & Sam\supers{l.} H Smith \\
					John Johnson & Orson Hyde \\
					Martin Harris & Sylvester Smith and \\
					John S Carter & Luke Johnson, \\
					\end{tabular}

					high priests, were chosen to be a standing council for the Church, by the unanimous voice of the council.

					The above named counsellors were then asked whether they accepted their appointments, and whether they would act in that office according to the law of Heaven: to which they all answered, that they accepted their several appointments, and would fill their offices according to the grace of God bestowed upon them.

					The numbers composeing the council, who voted in the name, and for the church in appointing the above named counsellors, were forty three; As follows: Nine high priests, Seventeen elders, four priests, and thirteen members.

					\uline{Voted}, that \sout{this} <the high> council cannot have power to act without seven of the above named counsellors, or their regularly appointed successors, are present; these seven shall have power to appoint other high priests whom they may consider worthy and capable to act in the place of absent counsellors.

					Voted, that whenever any vacancy shall occur by the death, removeal from office, for transgression, or removal from the bounds of this church goverment of any one of the above named counsellors, it shall be filled by the nomination of the president, or presidents and sanctioned by the voice of a general \sout{Conference} <Council of high priests> convened for that purpose to act in the name of the Church.

					The president of the church, who is also the president of the council, is appointed by the voice of the Saviour, and acknowledged in his administration, by the voice of the Church; and it is according to the dignity of his office that he should preside over the high Council of the Church; and it is his privilege to be assisted by two other presidents, appointed after the same manner that he himself was appointed; and in case of the abscence of one or both of those who are appointed to assist him, he has power to preside over the council without an assistant: and in case that he himself is abscent, the other presidents have power to preside in his stead, both or either of them.

					Whenever a high council of the Church of Christ, is regularly organized according to the foregoing pattern, it shall be the duty of the twelve counsellors to cast lots by numbers and thereby ascertain who of the twelve shall speak first, commenceing with Number One, and so in succession to number twelve

					Whenever this council convenes to act upon any case; \sout{in} \sout{the} \sout{Church}, the twelve counsellors shall consider whether it is a difficult one or not; If it is not, two <only> of the Counsellors shall speak upon it according to the form above written; but if it is thought to be \sout{more} \sout{<a>} difficult, \sout{<one>} four shall be appointed, and if \sout{still} more difficult, six: but in no case \sout{not} \sout{over} \sout{that} \sout{number} \sout{shall} <shall \sout{be} more than six \sout{be}> be appointed to speak. The accused in all cases has a right to one half of the council to prevent insult or injustice; and the counsellors appointed to speak before the council, are to present the case after the evidence is examined, in its true light before the Council, and every man is to speak according to equity and justice.

					Those counsellors who draw even numbers, that is, 2, 4, 6, 8, 10 and 12, are the individuals who are to stand up <in> the behalf of the accused and prevent insult or injustice.

					In all cases the accuser and the accused shall have a privilege of speaking for themselves before the council, after the evidences are heared, and the Counsellors who are appointed to speak on the case, have finished their remarks.

					After the evidences are heared; the counsellors, accuser and <the> accused, have spoken, the president shall give a decision according to the understanding which he shall have of the case, and call upon the twelve Counsellors to sanction the same by their voices.

					But should the remaining Counsellors who have not spoken*, or any one of them, after hearing the evidences and pleadings impartially, discover an error in the descision of the president, they can manifest it, and the case shall have a re-hearing; and if after a careful rehearing, any additional light is thrown upon the case, the descision shall be altered accordingly; but in case no additional light is given, the first decision shall stand; the majority of the Council haveing power to determine the same.

					In cases of difficulty respecting doctrine, or principle; if there is not a sufficiency written to make the case clear to the mind of the Council, the president may inquire and obtain the mind of the Lord by revelation.

					The high priests, when abroad, have power to call and organize a Council after the manner of the foregoing, to settle difficulties when the parties, or either of them shall request it, <and the said council of high priests shall have power to appoint one of their own number to preside over such council> \sout{by} \sout{appointing} \sout{or chooseing} \sout{one} \sout{of} \sout{their number} \sout{to} \sout{preside over} \sout{the} \sout{council} for the time being.

					It shall be the duty of said Council to transmit, immediately, a copy of their proceedings, with a full statement of the testimony \sout{with} <accompanying> their decision, to the high council at the seat of the government of the Church.

					Should the parties, or either of them, be dissatisfied with the decision of said Council, they may appeal to the high Council at the seat of the general government of the Church, and have a re-hearing, which case shall there be conducted according to the former pattern written, as though no such descision had been \sout{passed} <made>.

					This Council of high priests abroad, is only to be called on the most difficult cases of Church matters; and no common or ordinary case is to be sufficient to call such Councils. The travelling or located high priests abroad, have the power to say whether it is necessary to call such a Council or not.

					*The twelve counsellors then proceeded to cast lots or ballot, to ascertain who should speak first, and the following was the result, viz:

					\begin{tabular}{p{1.5in}p{1cm}p{1cm}p{1.5in}p{1cm}p{1cm}}				
					Oliver Cowdery & drew & No. 1 & John Johnson & drew & No 7 \\
					Joseph Coe------ &  & \textquotesingle\textquotesingle\ \textquotesingle\textquotesingle\ 2 & Orson Hyde &  & \textquotesingle\textquotesingle\ \textquotesingle\textquotesingle\ 8 \\
					Sam\supers{l.} H Smith------ &  & \textquotesingle\textquotesingle\ \textquotesingle\textquotesingle\ 3 & Jared Carter &  & \textquotesingle\textquotesingle\ \textquotesingle\textquotesingle\ 9 \\
					Luke Johnson ------ &  & \textquotesingle\textquotesingle\ \textquotesingle\textquotesingle\ 4 & Joseph Smith sen &  & \textquotesingle\textquotesingle\ 10 \\
					John S Carter------ &  & \textquotesingle\textquotesingle\ \textquotesingle\textquotesingle\ 5 & John Smith &  & \textquotesingle\textquotesingle\ \textquotesingle\textquotesingle\ 11 \\
					Sylvester Smith------ &  & \textquotesingle\textquotesingle\ \textquotesingle\textquotesingle\ 6 & Martin Harris &  & \textquotesingle\textquotesingle\ \textquotesingle\textquotesingle\ 12\\
					\end{tabular}

					Council then adjourned to meet on wednesday the 19\supers{th.} Inst. at 10 Oclk A.M.

					\begin{tightright}
					Orson Hyde Clk------
					\end{tightright}

					*Resolved, that the president\sout{s} or presidents at the seat of general church government, shall have power to determine whether any such case as may be appealed, is justly entitled to a re-hearing after examineing the appeal and the evidences and statements accompanying it.
					
		% -----
		\subsubsection{Minutes, 19 February 1834}
		% -----
			\label{EMS-1834-JS-19-FEB-1834}
			% \label{EMS-1830-JS-DAY-3LETTERMONTH-YEAR}
			
			\begin{description}
					\item[Print Sources]
				\end{description}

					\begin{tabular}{p{0.5in}p{1in}p{0.5in}p{1in}}
					JSP	 	& D3:444--448		& JSR 		& --- \\
					DC	 	& --- 		& HC 		& 2:31--34 \\
					HDDC 	& ---		& TCHC 		& 2:20--23 \\
					\end{tabular}
					% HDDC = woodford's DC diss; JSR = Marquardt's JS Revelations; TCHC = Vogel HC
				
				\begin{description}
					\item[Online Access] \href{https://www.josephsmithpapers.org/paper-summary/minutes-19-february-1834/1}{Here}
					% \item[Analysis] \hyperref[XX]{Here}
				\end{description}
				
				\begin{description}
					\item[Background]
				\end{description}
				
					% It might also be useful to give a two sentence context of the period: e.g., "This speech was given in Nauvoo to a small congregation one month prior to the destruction of the Nauvoo Expositor. These and other tensions may have been on the prophet's mind when he gave the speech."
				
				\begin{description}
					\item[Original Source]
				\end{description}
				
					% brief description of original source material, with reference to JSP or somewhere else for more info
					% Background material: it might be helpful to have a note that says whether a given document was presented as a speech, was written in a journal, etc.
					% Also a comment here about how reliable the source might be, or what shape it is in, if there are large omissions or parts that are hard to understand, and so on.
				
				\begin{description}
					\item[Text]
				\end{description}
				
					\begin{tightright}
							\hypertarget{EMS-1834-JS-19-FEB-1834-a}%
						Kirtland
							\marginnote{a}
						Feb. 19. 1834,
					\end{tightright}
					
					The Council assembled pursuant to adjournment. Joseph Smith Jn\supers{r.} opened the council by reading the 3\supers{rd.} Chap of of Joel's prophecy, and prayer. After which he arose before the Council, and said, that he had laboured the day before with all the strength and wisdom that he had given him in makeing the corrections necessary in the last council minutes which he would now read before this Council. He asked the council for their attention, that they might rightly judge upon the truth and propriety of these minutes, as all were equally interested in them \&c. He als[o] urged the necessity of prayer, that the spirit might be given, that the things of the spirit might be judged thereby; because the carnal mind cannot discern the things of God \&c. He then proceeded to read the minutes and afterwards made some remarks, when it was decided by the members of the council present, that it might be read a second time.
						\hypertarget{EMS-1834-JS-19-FEB-1834-b}%
					Sidney
						\marginnote{b}%
					Ridgon then proceeded to read the minutes or Constitution of the high Council the second time, remarking at the time, that it could not be justly urged to be read at this time, as the hour was passed which was appointed for the Council to assemble. An impropriety by some was discovered in the commencement of the minutes, as it says, a council of high priests, and afterwards says, that elders, priests and private members acted in said council.
						\hypertarget{EMS-1834-JS-19-FEB-1834-c}%
					Said
						\marginnote{c}%
					objections were corrected, and the minutes read the third time by Oliver Cowd[e]ry. The questions were then asked, whether the present Council acknowledged the same, and receive them for a form, and constitution of the high Council of the Church of Christ hereafter: The document was received by the unanimous voice of the Council, with this provision, that, if the president should hereafter discover any lack in the same he should be privileged to fill it up.

					The number present who received the above named document, was twenty six high priests, eighteen Elders, three priests, one teacher and fourteen private members, makeing in all, sixty two

						\hypertarget{EMS-1834-JS-19-FEB-1834-d}%
					After
						\marginnote{d}%
					much good instruction, Joseph, the president, laid his [hands?] upon the heads of the two assistant presidents and pronounced a blessing upon them, that they might have wisdom to magnify their office, and power over all the power of the adversary. He also laid his hands upon the twelve counsellors and commanded a blessing to rest upon them, that they might have wisdom and power to counsel in righteousness upon all subjects that might be laid before them. He also prayed that they might be delivered from those evils to which they were most exposed, and that their lives might be prolonged on the earth.

						\hypertarget{EMS-1834-JS-19-FEB-1834-e}%
					Joseph
						\marginnote{e}%
					Smith Sen. then laid his hands upon the head of his son, Joseph, and said: Joseph, I lay my hands upon thy head, and pronounc[e] the blessings of thy progenitors upon thee, that thou mayest hold the keys of the mysteries of the Kingdom of heaven until the coming of the Lord, Amen. He, also, laid his hands upon the head of his son Samuel [Smith] and said. Sam\supers{l.}, I lay my hands upon thy head and pronounce the blessing of thy progenitors upon thee, that thou mayest remain a priest of the most high God, and like Samuel of old, hear his voice, saying, Samuel, Samuel, Amen.

						\hypertarget{EMS-1834-JS-19-FEB-1834-f}%
					John
						\marginnote{f}%
					Johnson, also, laid his hands upon the head of his son Luke [Johnson] and said, my Father in Heaven, I ask thee to bless this my son according to the blessings of his forefathers, that he may be strengthened in his ministry according to his holy calling, Amen.

						\hypertarget{EMS-1834-JS-19-FEB-1834-g}%
					The
						\marginnote{g}%
					president then gave the assistant presidents a Solem charge to do their duty in righteousness and in the fear of God. He also charged the twelve Counsellors in a similar manner, all in the name of Jesus Christ. We then, all raised our hands to heaven in token of the everlasting Covenant, and the Lord blessed us with his spirit. He then said the Council was organized according to the ancient order, and also according to the mind of the Lord...

		% -----
		\subsubsection{Minutes, 20 February 1834}
		% -----
			\label{EMS-1834-JS-20-FEB-1834}
			% \label{EMS-1830-JS-DAY-3LETTERMONTH-YEAR}
			
			\begin{description}
					\item[Print Sources]
				\end{description}

					\begin{tabular}{p{0.5in}p{1in}p{0.5in}p{1in}}
					JSP	 	& D3:448--452		& JSR 		& --- \\
					DC	 	& --- 		& HC 		& 2:34--35 \\
					HDDC 	& ---		& TCHC 		& 2:24--26 \\
					\end{tabular}
					% HDDC = woodford's DC diss; JSR = Marquardt's JS Revelations; TCHC = Vogel HC
				
				\begin{description}
					\item[Online Access] \href{https://www.josephsmithpapers.org/paper-summary/minutes-20-february-1834/1}{Here}
					% \item[Analysis] \hyperref[XX]{Here}
				\end{description}
				
				\begin{description}
					\item[Background]
				\end{description}
				
					% It might also be useful to give a two sentence context of the period: e.g., "This speech was given in Nauvoo to a small congregation one month prior to the destruction of the Nauvoo Expositor. These and other tensions may have been on the prophet's mind when he gave the speech."
				
				\begin{description}
					\item[Original Source]
				\end{description}
				
					% brief description of original source material, with reference to JSP or somewhere else for more info
					% Background material: it might be helpful to have a note that says whether a given document was presented as a speech, was written in a journal, etc.
					% Also a comment here about how reliable the source might be, or what shape it is in, if there are large omissions or parts that are hard to understand, and so on.
				
				\begin{description}
					\item[Text]
				\end{description}
				
					\begin{tightright}
							\hypertarget{EMS-1834-JS-20-FEB-1834-a}%
						Kirtland,
							\marginnote{a}
						20 Feb'\supers{y.} 1834.
					\end{tightright}
					
					High council met this evening accord[i]ng to appointment to determine concerning the Elders going out to preach \&c. The presidnet opened the council by prayer.---

					At a church meeting held in Pennsylvania, Erie Co. and Springfield Township by Orson Pratt \& Lyman Johnson, high priests; Some of the members of that Church refused to partak[e] of the sacrament because the Elder administering it did not observe the words of wisdom to obey them. Lyman argued that they were justified in so doing because the Elder was in transgression. Orson argued that the Church was bound to receive the supper under the administration of an Elder so long as he retained his office, or licence. Voted that six counsellors should speak upon the subject, or case.

						\hypertarget{EMS-1834-JS-20-FEB-1834-b}%
					The
						\marginnote{b}%
					Council then proceeded to try the question, whether disobedience to the word of wisdom was a transgression sufficient to deprive an official member from holding an office in the church, after haveing it sufficiently taught him. Samuel H Smith, Luke Johnson, John S Carter, Sylvester Smith, John Johnson and Orson Hyde were called to speak upon the case then before the council. After the counsellors had spoken, the President proceeded to give a decision: ``That no official member in this church is worthy to hold an office after haveing the word\sout{s} of wisdom properly taught to him, and he, the official member, neglecting to comply with, or obey them;[''] after which the Counsellors voted according to the same.

						\hypertarget{EMS-1834-JS-20-FEB-1834-c}%
					The
						\marginnote{c}%
					president then asked if there were any Elders present who would go to Canada and preach the gospel to that people; for they have written, said he, a number of letters for help; and the whole council felt as though the Spirit required brethren to go there. It was, therefore, decided by the Council that Lyman Johnson and Milton Holmes should travel together into Canada, and also that Zebedee Coultrin [Coltrin] and Henry Herriman travel together into Can[a]da. It was also decided that Jared Carter and Phineas Youngs [Young] travel together if they can arrange their affairs at home so as to be liberated.
						\hypertarget{EMS-1834-JS-20-FEB-1834-d}%
					It
						\marginnote{d}%
					was also decided that Bro. Oliver Granger should travel eastward as soon as his circumstances will permit, and that he should travel alone on account of his age. It was also decided that bro. Martin Harris should travel alone whenever he travels. Bro's. John S Carter \& Jesse Smith travel east together as soon as they can. The council also decided that bro. Brig[h]am Young should travel alone, it being his own choice. Decided also, that James Durfee and Edmund Marvin should travel together eastward. Also that Sidney Rigdon and John P Green[e] go to Strongsville.
						\hypertarget{EMS-1834-JS-20-FEB-1834-e}%
					Also
						\marginnote{e}%
					that bro's. Orson Pratt and Harrison Sagers travel together for the time being, and that there should be a general conference held in Saco in the State of Maine on the 13 day of [p. 40] June 1834. It was furthermore voted that bro. Orson Hyde, accompanied by bro. Orson Pratt, go east to obtain donations for Zion, and means to redeem the farm on which the house of the Lord stands. The Church and Council then prayed with uplifted hands that they might be prospered in their Mission. Conference adjourned after the usual form
					
					\begin{tightcenter}
					by Order of the Conference
					\end{tightcenter}

					\begin{tabular}{p{3in}p{1.5in}p{1in}}
					 & Orson Hyde \&) & \\
					 &  & Clks \\
					 & Oliver Cowdery,) & \\
					 \end{tabular}

		% -----
		\subsubsection{Minutes, 24 February 1834}
		% -----
			\label{EMS-1834-JS-24-FEB-1834-MIN}
			% \label{EMS-1830-JS-DAY-3LETTERMONTH-YEAR}
			
			\begin{description}
					\item[Print Sources]
				\end{description}

					\begin{tabular}{p{0.5in}p{1in}p{0.5in}p{1in}}
					JSP	 	& D3:453--457		& JSR 		& --- \\
					DC	 	& --- 		& HC 		& 2:39--40 \\
					HDDC 	& ---		& TCHC 		& 2:31--32 \\
					\end{tabular}
					% HDDC = woodford's DC diss; JSR = Marquardt's JS Revelations; TCHC = Vogel HC
				
				\begin{description}
					\item[Online Access] \href{https://www.josephsmithpapers.org/paper-summary/minutes-24-february-1834/1}{Here}
					% \item[Analysis] \hyperref[XX]{Here}
				\end{description}
				
				\begin{description}
					\item[Background]
				\end{description}
				
					% It might also be useful to give a two sentence context of the period: e.g., "This speech was given in Nauvoo to a small congregation one month prior to the destruction of the Nauvoo Expositor. These and other tensions may have been on the prophet's mind when he gave the speech."
				
				\begin{description}
					\item[Original Source]
				\end{description}
				
					% brief description of original source material, with reference to JSP or somewhere else for more info
					% Background material: it might be helpful to have a note that says whether a given document was presented as a speech, was written in a journal, etc.
					% Also a comment here about how reliable the source might be, or what shape it is in, if there are large omissions or parts that are hard to understand, and so on.
				
				\begin{description}
					\item[Text]
				\end{description}
				
						\hypertarget{EMS-1834-JS-24-FEB-1834-MIN-a}%
					Kirtland,
						\marginnote{a}
					Feb.' 24, 1834.
					
					The high council of the Church met this day at the house of Joseph Smith Jun\supers{r.} for the purpose of giveing an audience or hearing to Lyman Wight and Parley [P.] Pratt, representatives from Zion, to represent to us the state of the church in that place.
					
						\hypertarget{EMS-1834-JS-24-FEB-1834-MIN-b}%
					Joseph,
						\marginnote{b}%
					the president, opened the council by prayer. Two of the standing counsellors were absent, namely, Joseph Coe and John Smith. Hyrum Smith was chosen to act in the place of John Smith and John P Green[e] to act in the place of Joseph Coe. Thus the high council was organiz[e]d and six of the counsellors were appointed to speak. Bro's. P. Pratt and L. Wight, messengers from Zion, arose and laid their business before the council and delivered their message. the substance of which, was, an inquiry when, how and by what means Zion was to be redeemed from our enemies.
						\hypertarget{EMS-1834-JS-24-FEB-1834-MIN-c}%
					They
						\marginnote{c}%
					said that our brethern who had been driven away from their lands and scattered abroad had found so much favour in the eyes of the people that they could obtain food and raiment of them for their labour insomuch that they were comfortable. But the idea of being driven away from the land of Zion pained their very souls and they desired of God, by earnest prayer, to return with songs of everlasting joy as said Isaiah, the Prophet.
					
						\hypertarget{EMS-1834-JS-24-FEB-1834-MIN-d}%
					They
						\marginnote{d}%
					also said that none of their lands were sold into the hands of our enemies except a piece owned by bro. W\supersu{m}\supers{.} E. M\supers{c.}Lellin of thirty acres which he sold into the hands of the enemy, and seven acres more which he would have sold to the enemy if a brother had not come forward \& purchased it and paid him his money Bro. Joseph then arose and said that he was going to Zion to assist in redeeming it. He then called for the voice of the Council to sanction his going which was given without a dissenting voice.
						\hypertarget{EMS-1834-JS-24-FEB-1834-MIN-e}%
					He
						\marginnote{e}%
					then called for volunteers to go with him, when some thirty or forty volunteered to go who were then present at the council. It was a question whether we should go by water or by land, and after a short investigation it was decided unanimously that we go by land. Joseph Smith Jun. was nominated and seconded to be the Commander in Chief of the Armies of Israel and the leader of those who volunteered to go and assist in the redemption of Zion, and carried by the vote of all present. Council then adjourned by prayer and thanksgiveing.
					
					\begin{tabular}{p{3in}p{1.5in}p{1in}}
					 & Orson Hyde and) & \\
					 &  & C'lks. \\
					 & Oliver Cowdery) & \\
					 \end{tabular}
					 
		% -----
		\subsubsection{Revelation, 24 February 1834 [D\&C 103]}
		% -----
			\label{EMS-1834-JS-24-FEB-1834}
			% \label{EMS-1830-JS-DAY-3LETTERMONTH-YEAR}
			
			\begin{description}
					\item[Print Sources]
				\end{description}

					\begin{tabular}{p{0.5in}p{1in}p{0.5in}p{1in}}
					JSP	 	& D3:457--463		& JSR 		& 253--255 \\
					DC	 	& Section 103 		& HC 		& 2:36--39 \\
					HDDC 	& 1327--1347		& TCHC 		& 2:29--31 \\
					\end{tabular}
					% HDDC = woodford's DC diss; JSR = Marquardt's JS Revelations; TCHC = Vogel HC
				
				\begin{description}
					\item[Online Access] \href{https://www.josephsmithpapers.org/paper-summary/revelation-24-february-1834-dc-103/1}{Here}
					% \item[Analysis] \hyperref[DC103]{Here}
				\end{description}
				
				\begin{description}
					\item[Background]
				\end{description}
				
					% It might also be useful to give a two sentence context of the period: e.g., "This speech was given in Nauvoo to a small congregation one month prior to the destruction of the Nauvoo Expositor. These and other tensions may have been on the prophet's mind when he gave the speech."
				
				\begin{description}
					\item[Original Source]
				\end{description}
				
					% brief description of original source material, with reference to JSP or somewhere else for more info
					% Background material: it might be helpful to have a note that says whether a given document was presented as a speech, was written in a journal, etc.
					% Also a comment here about how reliable the source might be, or what shape it is in, if there are large omissions or parts that are hard to understand, and so on.
				
				\begin{description}
					\item[Text]
				\end{description}
				
					Verily I say unto you my Friends behold, I will give unto you a revelation \& commandment, that you may know how to act in the discharge of your duties concerning the salvation \& redemption of your brethren who have been scattered from the land of Zion: being driven \& smitten by the\sout{r} hands of mine enemies on whom I will pour out of my wrath without measure in mine own time; for I have suffered them thus far. that they might fill up the measure of their iniquities that their cup might be full, \& that those who call themselves after my name might be chastened for a little season, with a sore \& grievous chastisement; because they did not hearken all together unto the precepts \& commandments which I gave unto them. But verily I say unto you, that I have decreed a decree which my people shall realize inasmuch as they hearken from this hour unto the counsel which <I> the Lord their God shall give unto them. Behold, they shall, for I have decreed it, begin to prevail against mine enemies; from this verry hour, \& by hearkening to observe all the words which I the Lord their God shall speak unto them, they shall never cease to prevail untill the kingdoms of the world, are subdued under my feet, \& the earth is given unto the saints to possess it forever, \& ever. But inasmuch as they keep not my commandments \& hearken not to observe all my words, the kingdoms of the world shall prevail against them; for they were set to be a light unto the world, \& to be the Saviours of men: \& \sout{a} inasmuch as they are not the Saviours of men they are as salt that has lost its \sout{its} savor, \& is thenceforth good for nothing but to be cast out \& to be trodden under the feet of men. But verily, I say unto you, I have decreed that your brethren who have been scattered shall return to the lands of their inheritances \& build up the waste places of Zion; for after much tribulation, as I have said unto you in a former commandment, commeth the blessing: behold this is the blessing which I promised after your tribulations, \& the tribulations, of your brethren; your redemption \& the redemption of your brethren, even their restoration to the land of Zion, to be established no more to be thrown down. Nevertheless if they shall pollute their inheritances they shall be thrown down; for I will not spare them if they shall pollute their inheritances. Behold I say unto you, that the redemtion of Zion must needs come by power; therefore, I will raise up unto my people a man who shall lead them like as Moses led the children of Israel; for ye are the children of Israel, \sout{\&} \& of the seed of Abraham; \& ye must needs be led out of bondage by power, \& with a stretched out arm. And as your Fathers were led at the first even so shall the redemtion of Zion be. Therefore let not your hearts \sout{be} faint; for I say not unto you as I said unto your fathers, mine Angel shall go up before you, but not my presence; but I say unto you, mine Angel shall go up before you, \& also my presence. And in time ye shall possess the goodly land. verily, verily, I say unto you, that my servant Joseph is the man to whom I likened the servant to whom the Lord of the vinyard spake in the parable which I have given unto you. Therefore, let my servant Joseph, say unto the strength of my house, my young men, \& the middle aged, Gather ye together unto the land of Zion, upon the lands which I have bought with moneys that have been consecrated unto me; \sout{a} \& let all the churches send up wise men with their moneys \& purchase lands, even as I have commanded them. And inasmuch as mine enemies come against you to drive you from my goodly land which I have consecrated to be the land of Zion; even from your own lands after these testimonies which ye have brought before me against them, ye shall curse them; \& whomsoever ye curse I will curse, And ye shall avenge me of mine enemies; and my presence shall be with you even in avenging me of mine enemies, unto the third \& fourth generation of them that hate me. Let no man be afraid to lay down his life for my sake; for whoso layeth down his life for my \sout{my} sake \sout{sake} shall find it again; \& whoso is not willing to lay down his life for my sake is not my dissiple. It is my will that my servant Sidney [Rigdon] should lift up his voice in the congregations in the eastern countries in preparing the churches to keep the commandments which I have given unto them concerning the restoration \& redemption of Zion, It is my will that my servant Parley [P. Pratt], \& my servant Lyman [Wight] should not return to the land of their brethren until they have obtained companies to go up unto the land of Zion, by tens, or by twenties or by fifties, or by <a> hundred, until they have obtained unto the number of five hundred of the strength of my house. Behold, this is my will; ask \& ye shall receive; but men do not always do my will; therefore if ye cannot obtain five hundred, seek diligently, that peradventure ye may obtain three; \& if ye cannot obtain three hundred, seek diligently that peradventure ye may obtain one hundred: But verily I say unto you, a commandment I give unto you that you shall not go up unto the Land of Zion until you have obtained a hundred of the strength of my house, to go up with you unto the land of Zion. Therefore as I said unto you, ask \& you shall receive; pray earnestly, that peradventure my servant Joseph may go up with you \& preside in the midst of my people \& organize my kingdom upon the consecrated land, \& establish the children of Zion upon the laws \& commandments which have been given \& which shall be given unto you. All victory and glory is brought to pass unto you through your diligence, faithfulness, \& prayers of faith. Let my Servant Parley journey with my servant Joseph; let my servant Lyman Journey with my servant Sidney: let my servant Hyrum [Smith] journey with my servant Frederick [G. Williams]: let my servant Orson Hyde journey with my servant Orson Pratt, whithersoever my servant Joseph shall counsel them in obtaining the fulfillment of these commandments which I have given unto you, \& leave the residue in \sout{in} my hands; even so, Amen.
					
		% -----
		\subsubsection{Letter to the Church, circa March 1834}
		% -----
			\label{EMS-1834-JS-CIR-MAR-1834}
			% \label{EMS-1830-JS-DAY-3LETTERMONTH-YEAR}
			
			\begin{description}
					\item[Print Sources]
				\end{description}

					\begin{tabular}{p{0.5in}p{1in}p{0.5in}p{1in}}
					JSP	 	& D3:472--484		& JSR 		& --- \\
					DC	 	& --- 		& HC 		& 2:13--22 \\
					HDDC 	& ---		& TCHC 		& --- \\
					TPJS 	& 56--66			& 	 		&  \\
					\end{tabular}
					% HDDC = woodford's DC diss; JSR = Marquardt's JS Revelations; TCHC = Vogel HC
				
				\begin{description}
					\item[Online Access] \href{https://www.josephsmithpapers.org/paper-summary/letter-to-the-church-circa-march-1834/1}{Here}
					% \item[Analysis] \hyperref[XX]{Here}
				\end{description}
				
				\begin{description}
					\item[Background]
				\end{description}
				
					% It might also be useful to give a two sentence context of the period: e.g., "This speech was given in Nauvoo to a small congregation one month prior to the destruction of the Nauvoo Expositor. These and other tensions may have been on the prophet's mind when he gave the speech."
				
				\begin{description}
					\item[Original Source]
				\end{description}
				
					% brief description of original source material, with reference to JSP or somewhere else for more info
					% Background material: it might be helpful to have a note that says whether a given document was presented as a speech, was written in a journal, etc.
					% Also a comment here about how reliable the source might be, or what shape it is in, if there are large omissions or parts that are hard to understand, and so on.
				
				\begin{description}
					\item[Text]
				\end{description}
				
					\begin{tightcenter}
							\hypertarget{EMS-1834-JS-CIR-MAR-1834-a}%
						THE 
							\marginnote{a}
						ELDERS OF THE CHURCH IN KIRTLAND, TO THEIR BRETHREN ABROAD.
					
						\textit{(Continued from our last.)}
					\end{tightcenter}

					\begin{tightright}
						\textit{Dear brethren in Christ, and companions in tribulation.}
					\end{tightright}

					HAVING in a former number of the Star, written you quite lengthy on some few items connected with the religion which we profess, we deem it of importance to the cause in which all our united efforts ought, with an eye single to the glory of God, to be engaged, that we may escape the corruptions of the world, and not only show ourselves approved in his sight, but may be instruments in the order of his providence in convincing some of our fellow-travellors to eternity of the importance of turning from error to righteousness, and embracing the fulness of the everlasting gospel---to continue this letter of instruction and exhortation, believing, (as we have previously remarked,) that on your part it will be received in brotherly fellowship.
						\hypertarget{EMS-1834-JS-CIR-MAR-1834-b}%
					We
						\marginnote{b}%
					would remind you, brethren, of the fateagues, trials, privations, and persecutions, which the ancient saints endured for the only purpose of persuading men of the excellency and propriety of the faith of Christ, were it in our opinion necessary, or would serve in any respect to stimulate you to labor in the vineyard of the Lord with any more diligence; but we have reason to believe, (if you make the holy scriptures a sufficient part of your studies,) that their perseverance is known to you all; and that they were willing to sacrifice the present honors and pleasures of this world, that they might obtain an assurance of a crown of life from the hand of our Lord; and their excellent examples in labor, which manifests their zeal to us in the cause which they embraced, you are daily striving to pattern.
						\hypertarget{EMS-1834-JS-CIR-MAR-1834-c}%
					And
						\marginnote{c}%
					not only these, but the commandments of our Lord, we hope, are constantly revolving in your hearts, teaching you, not only his will in proclaiming his gospel, but his meekness and perfect walk before all, even in those times of severe persecutions and abuse which were heaped upon him by a wicked and adulterous generation. Remember, brethren, that he has called you unto holiness; and need we say, to be like him in purity? How wise; how holy; how chaste, and how perfect, then, you ought to conduct yourselves in his sight; and remember too, that his eyes are continually upon you.
						\hypertarget{EMS-1834-JS-CIR-MAR-1834-d}%
					Viewing
						\marginnote{d}%
					these facts in a proper light, you cannot be insensible, that without a strict observance of all his divine requirements, you may, at last, be found wanting; and if so, you will admit, that your lot will be cast among the unprofitable servants.--- We beseech you, therefore brethren, to improve upon all things committed to your charge, that you lose not your reward!

						\hypertarget{EMS-1834-JS-CIR-MAR-1834-e}%
					No
						\marginnote{e}%
					doubt, the course which we pursued in our last to you, is yet familiar to your minds; that we there endeavored to show, as far as our limits would extend, the propriety, in part of adhering to the law of heaven; and also, the consistency in looking to heaven for a law or rule to serve us as a guide in this present state of existence, that we may be prepared to meet that which inevitably awaits us, as well as all mankind.--- There is an importance, perhaps, attached to this subject, which the world has not as fully examined as the importance of it requires.
						\hypertarget{EMS-1834-JS-CIR-MAR-1834-f}%
					Think
						\marginnote{f}%
					for a moment, of the greatness of the Being who created the universe; and ask, Could he be so inconsistant with his own character, as to leave man without a law or rule to regulate his conduct, after placing him here, where, according to the formation of his nature he must in a short period sink into the dust? Is there nothing further; is there no existence beyond this vail of death which is so suddenly to be cast over all of us? If there is, why not that Being who had power to place us here, inform us something concerning hereafter? If we had power to place ourselves in this present existence, why not have power to know what shall follow when that dark vail is cast over our bodies? If in this life we receive our all; if when we crumble back to dust we are no more, from what source did we emanate, and what was the purpose in our existence?
						\hypertarget{EMS-1834-JS-CIR-MAR-1834-g}%
					If
						\marginnote{g}%
					this were all, we should be led to query, whether there was really any substance in existence: and we might with propriety say, ``Let us eat and drink; for tomorrow we die!'' If this were really so, then why this constant toiling, why this continual warfare, and why this unceasing trouble? But this is not the case, the voice of \textsc{reason}, the language of \textsc{inspiration}, and the Spirit of the living GOD, our Creator, teaches us, as we hold the record of truth in our hands, that this is not the case; that this is not so; for, the heavens declare the glory of a GOD, and the firmament shows his handy work;
						\hypertarget{EMS-1834-JS-CIR-MAR-1834-h}%
					and
						\marginnote{h}%
					a moment's reflection, is sufficient to teach every man of common intellect, that all these are not the mere production of \textit{chance}, nor could they be supported by any power less than by an Almighty hand: and he that can mark the power of Omnipotence inscribed upon the heavens, can also see His own hand-writing in the sacred volume; and he who reads it oftenest will like it best, and he who is acquainted with it, will know the hand wherever he can see it; and when once discovered, it will not only receive an acknowledgment, but an obedience to all its heavenly precepts. For a moment reflect, what could have been the purpose in our Father in giving to us a law? Was it that it might be obeyed, or disobeyed?
						\hypertarget{EMS-1834-JS-CIR-MAR-1834-i}%
					And
						\marginnote{i}%
					think further too, not only the propriety, but the importance of attending to his laws in every particular. If, then, there is an importance in this respect, is there not a responsibility of great weight resting upon those who are called to declare these truths to men? Could we, or were we capable of laying any thing before you as a just comparison, we would cheerfully do it; but in this our capacity fails, and we are inclined to think, that man is unable, without an assistance beyond what has been given to those before us, of expressing in words the greatness of this important office.
						\hypertarget{EMS-1834-JS-CIR-MAR-1834-j}%
					We
						\marginnote{j}%
					can only say, that if an anticipation of the joys of the celestial glory, as witnessed to the hearts of the humble is not sufficient, we will leave with yourselves the result of your own diligence; for God ere long, will call \textit{all} his servants before him, and there from His own hand they will receive a just recompense and a righteous reward for all their labors.

						\hypertarget{EMS-1834-JS-CIR-MAR-1834-k}%
					So
						\marginnote{k}%
					much by way of introduction, and we shall now proceed to examine still further the subject of law. However little may have been heretofore thought, or said upon the subject of law, does not diminish in the least the propriety nor the design of it, since it emanated from God; and though it may have been, and may be at this day a subject untouched by the professors of christianity, that does not lessen its value, neither does it diminish its power in judging men from their actions according to it, at the last day, those who have, or may have come to a knowledge of it.
						\hypertarget{EMS-1834-JS-CIR-MAR-1834-l}%
					It
						\marginnote{l}%
					may be supposed, and we think with a degree of propriety, that man had given to him in the beginning, from the hand of his Maker, every necessary law and instruction, for his peace, happiness, and future comfort; and if not, living as he did in the immediate presence, and walking under the inspection of heaven, if he needed more, he could yet ask it, and that wise Hand which had formed him of the dust was sufficient; not only sufficient, but knowing all things, knew whether man needed more or not, and if he did, it would be bestowed. To suppose that the Maker of the universe never gave to man any law after he had formed him, would, in our opinion, be offering an insult to his glorious character, and be comparing him beneath, even an earthly parent!
						\hypertarget{EMS-1834-JS-CIR-MAR-1834-m}%
					For
						\marginnote{m}%
					where, we ask, is the kind humane father to be found, who would, for any consideration whatever, suffer his children to grow up to manhood without giving them instruction, and instruction too, which would be wisely calculated to benefit them, even in ripened years? Should he teach them virtue in their youth, (a principle too much neglected with most parents,) if observed in age it certainly would be virtue still; and the more it was observed the more honorable would be the gray hairs, until his spirit took its welcome exit to mingle with its kindred spirits, and rejoice in the salvation of that God from whom came the first principles of virtue.
						\hypertarget{EMS-1834-JS-CIR-MAR-1834-n}%
					Should
						\marginnote{n}%
					the great Author of our being, after he had made all things, and even man, and pronounced them \textit{all} good, leave man without a law, we might well suppose that here was a contradiction in terms, indeed; for he had pronounced all things which he had made \textsc{good}, and yet there was no good in man, consequently he was not worthy to receive a law whereby his conduct might be governed; but must be left without any principles or directions from the hand of his Maker to guide him in the least particular.

						\hypertarget{EMS-1834-JS-CIR-MAR-1834-o}%
					From
						\marginnote{o}%
					these facts, in short, and the further knowledge contained in the scriptures, it is reasonable to suppose, that man departed from the first teachings, or instructions which he received from heaven in the first age, and refused by his disobedience to be governed by them. Consequently, he formed such laws as best suited his own mind, or as he supposed, best adapted to his situation.
						\hypertarget{EMS-1834-JS-CIR-MAR-1834-p}%
					But
						\marginnote{p}%
					that God has influenced man more or less since that time in the formation of law for his benefit we have no hesitancy in believing; for, as before remarked, he being the source of all good, every just and equitable law was in a greater or less degree influenced by him. And though man in his own supposed wisdom would not admit the influence of a power superior to his own, yet for wise and great purposes, for the good and happiness of his creatures, God has instructed man to form wise and wholesome laws, since he had departed from him and refused to be governed by those laws which he had given by his own voice from on high in the beginning.
						\hypertarget{EMS-1834-JS-CIR-MAR-1834-q}%
					But
						\marginnote{q}%
					notwithstanding this trangression, by which man had cut himself off from an immediate intercourse with his Maker without a Mediator, it appears that the great and glorious plan of his redemption was previously meditated; the sacrifice prepared; the atonement wrought out in the mind and purpose of God, even in the person of the Son, through whom man was now to look for acceptance, and through whose merits he was now taught that he alone could find redemption, since the word had been pronounced, Unto dust thou shalt return!

						\hypertarget{EMS-1834-JS-CIR-MAR-1834-r}%
					But
						\marginnote{r}%
					that man was not sufficient of himself to erect a system, or plan with power sufficient to free him from a destruction which awaited him, is evident from the fact, that God, as before remarked, prepared a sacrifice in the gift of his own Son which should be sent in due time, in his own wisdom, to prepare a way, or open a door through which man might enter into his presence, from whence he had been cast for disobedience.---
						\hypertarget{EMS-1834-JS-CIR-MAR-1834-s}%
					From
						\marginnote{s}%
					time to time these glad tidings were sounded in the ears of men in different ages of the world down to the time of his coming. By faith in this atonement or plan of redemption, Abel offered to God a sacrifice that was accepted, which was the firstlings of the flock. Cain offered of the fruit of the ground, and was not accepted, because he could not do it in faith: he could have no faith, or could not exercise faith contrary to the plan of heaven.
						\hypertarget{EMS-1834-JS-CIR-MAR-1834-t}%
					It
						\marginnote{t}%
					must be the shedding of the blood of the Only Begotten to atone for man; for this was the plan of redemption; and without the shedding of blood was no remission; and as the sacrifice was instituted for a type, by which man was to discern the great Sacrifice which God had prepared; to offer a sacrifice contrary to that, no faith could be exercised, because redemption was not purchased in that way, nor the power of atonement instituted after that order; consequently, Cain could have no faith: and whatsoever is not of faith is sin. But Abel offered an acceptable sacrifice, by which he obtained witness that he was righteous, God himself testifying of his gifts.
						\hypertarget{EMS-1834-JS-CIR-MAR-1834-u}%
					Certainly,
						\marginnote{u}%
					the shedding of the blood of a beast could be beneficial to no man, except it was done in imitation, or as a type, or explanation of what was to be offered through the gift of God himself; and this performance done with an eye looking forward in faith on the power of that great Sacrifice for a remission of sins. But however various may have been, and may be at the present time the opinions of men respecting the conduct of Abel, and the knowledge which he had on the subject of atonement, it is evident in our minds, that he was instructed more fully into the plan than what the bible speaks; for how could he offer a sacrifice in faith, looking to God for a remission of his sins in the power of the great Atonement, without having been previously instructed into that plan?
						\hypertarget{EMS-1834-JS-CIR-MAR-1834-v}%
					And
						\marginnote{v}%
					further, if he was accepted of God, what were the ordinances performed further than the offering of the firstlings of the flock? It is said by Paul in his letter to his Hebrew brethren, that Abel obtained witness that he was righteous, God testifying of his gifts. To whom did God testify of the gifts of Abel, was it to Paul? We have very little on this important subject in the fore part of the bible. But it is said, that Abel himself obtained witness that he was righteous. Then certainly God spoke to him: indeed, it is said that God talked with him; and if he did, would he not, seeing he was righteous, deliver to him the whole plan of the gospel? And is not the gospel the news of redemtion?
						\hypertarget{EMS-1834-JS-CIR-MAR-1834-w}%
					How
						\marginnote{w}%
					could Abel offer a sacrifice and look forward with faith on the Son of God for a remission of his sins, and not understand the gospel? The mere shedding the blood of beasts or offering any thing else in sacrifice, could not procure a remission of sins, except it were performed in faith of something to come, if it could, Cain's offering must have been as good as Abel's. And if Abel was taught of the coming of the Son of God, was he not taught of his ordinances? We all admit that the gospel has ordinances, and if so, had it not always ordinances, and were not its ordinances always the same?
						\hypertarget{EMS-1834-JS-CIR-MAR-1834-x}%
					Perhaps,
						\marginnote{x}%
					our friends will say, that the gospel and its ordinances were not known till the days of John the son of Zecharias, in the days of Herod the king of Judea. But we will here look at this point: For our own part, we cannot believe, that the ancients in all ages were so ignorant of the system of heaven as many suppose, since all that were ever saved, were saved through the power of this great plan of redemption, as much so before the coming of Christ as since; if not, God has had different plans in operation, (if we may so express it,) to bring men back to dwell with himself;
						\hypertarget{EMS-1834-JS-CIR-MAR-1834-y}%
					and
						\marginnote{y}%
					this we cannot believe, since there has been no change in the constitution of man since he fell; and the ordinance or institution of offering blood in sacrifice, was only designed to be performed till Christ was offered up and shed his blood, as said before, that man might look forward with faith to that time. It will be noticed that according to Paul, -[see Gal. \textsc{iii} 8.]- the gospel was preached to Abraham.
						\hypertarget{EMS-1834-JS-CIR-MAR-1834-z}%
					We
						\marginnote{z}%
					would like to be informed in what name the gospel was then preached, whether it was in the name of Christ or some other name? If in any other name, was it the gospel? And if it was the gospel, and that preached in the name of Christ, had it any ordinances? If not, was it the gospel? And if it had, what were they? Our friends may say, perhaps, that there were never any ordinances except those of offering sacrifices, before the coming of Christ, and that it could not be possible for the gospel to have been administered while the sacrifices of blood were. But we will recollect, that Abraham offered sacrifice, and notwithstanding this, had the gospel preached to him.
						\hypertarget{EMS-1834-JS-CIR-MAR-1834-aa}%
					That
						\marginnote{aa}%
					the offering of sacrifice was only to point the mind forward to Christ, we infer from these remarkable words of his to the Jews, Your father Abraham rejoiced to see my day: and he saw it and was glad. -[See John \textsc{viii} 56.]- So, then, because the ancients offered sacrifice it did not hinder their hearing the gospel; but served, as we said before, to open their eyes, and enabled them to look forward to the time of the coming of the Savior, and to rejoice in his redemption.
						\hypertarget{EMS-1834-JS-CIR-MAR-1834-bb}%
					We
						\marginnote{bb}%
					find also, that when the Israelites came out of Egypt they had the gospel preached to them, according to Paul in his letter to the Hebrews, which says, For unto us was the gospel preached, as well as unto them: but the word preached did not profit them, not being mixed with faith in them that heard it. -[See Heb. \textsc{iv} 2.]- It is said again, in Gal. \textsc{iii} 19, that the law -[of Moses, or the Levitical law]- was added because of transgression. What, we ask, was this law added to, if it was not added to the gospel? It must be plain that it was added to the gospel, since we learn that they had the gospel preached to them.
						\hypertarget{EMS-1834-JS-CIR-MAR-1834-cc}%
					From
						\marginnote{cc}%
					these few facts, we conclude, that whenever the Lord revealed himself to men in ancient days, and commanded them to offer sacrifice to him, that it was done that they might look forward in faith to the time of his coming, and rely upon the power of that atonement for a remission of their sins. And this they have done, thousands who have gone before us, whose garments are spotless, and who are, like Job, waiting with an assurance like his, that they will see him in the \textit{latter day} upon the earth, even in their flesh.
		
						\hypertarget{EMS-1834-JS-CIR-MAR-1834-dd}%
					We
						\marginnote{dd}%
					may conclude, that though there were different dispensations, yet all things which God communicated to his people, were calculated to draw their minds to the great object, and to teach them to rely upon him alone as the Author of their salvation, as contained in his law. From what we can draw from the scriptures relative to the teachings of heaven we are induced to think, that much instruction has been given to man since the beginning which we have not. This may not agree with the opinions of some of our friends, who are bold to say, that we have every thing written in the bible which God ever spake to men since the world began, and that if he had ever said any thing more we should certainly have received it.
						\hypertarget{EMS-1834-JS-CIR-MAR-1834-ee}%
					But
						\marginnote{ee}%
					we ask, does it remain for a people who never had faith enough to call down one scrap of revelation from heaven, and for all they have now, are indebted to the faith of another people who lived hundreds and thousands of years before them, to say how much God has spoken and how much he has not spoken? We have what we have, and the bible contains what it does contain; but to say that God never said any thing more to man than is there recorded, would be saying at once, that we have at last received a revelation;
						\hypertarget{EMS-1834-JS-CIR-MAR-1834-ff}%
					for
						\marginnote{ff}%
					it must be one to advance thus far, because it is no where said in that volume by the mouth of God, that he would not, after giving what is there contained, speak again; and if any man has found out that for a fact, he has ascertained it by an immediate revelation, other than has been previously written by the prophets and apostles. But through the kind providence of our Father a portion of his word which he delivered to his ancient saints, has fallen into our hands, and they are presented to us with a promise of a reward if obeyed, and with a penalty if disobeyed; and that all are deeply interested in these laws, or teachings, must be admitted by all who acknowledge their divine authenticity.

						\hypertarget{EMS-1834-JS-CIR-MAR-1834-gg}%
					It
						\marginnote{gg}%
					may be proper for us to notice in this place, a few of the many blessings held out in this law of heaven as a reward to those who obey its teachings. God has appointed a day in which he will judge the world, and this he has given an assurance of in that he raised up his Son Jesus Christ from the dead; the point on which the hope of all who believe the inspired record is founded for their future happiness and enjoyment: because, if Christ is not risen, said Paul to the Corinthians, your faith is vain; ye are yet in your sins: and those who have fallen asleep in him have perished. -[See 1 Cor. \textsc{xv}.]-
						\hypertarget{EMS-1834-JS-CIR-MAR-1834-hh}%
					If
						\marginnote{hh}%
					the resurrection from the dead is not an important point, or item in our faith, we must confess that we know nothing about it; for if there is no resurrection from the dead, then Christ has not risen; and if Christ has not risen he was not the Son of God; and if he was not the Son of God there is not nor cannot be a Son of God, if the present book called the scriptures is true; because the time has gone by when, according to that book he was to make his appearance.
						\hypertarget{EMS-1834-JS-CIR-MAR-1834-ii}%
					On
						\marginnote{ii}%
					this subject, however, we are reminded of the words of Peter to the Jewish Sanhedrim, when speaking of Christ, he says, that God raised him from the dead, and we -[the apostles]- are his witnesses of these things, and so is also the Holy Ghost, whom God hath given to them that obey him. -[See Acts \textsc{v}.]- So that after the testimony of the scriptures on this point, the assurance is given by the Holy Ghost, bearing witness to those who obey him, that Christ himself has assuredly risen from the dead; and if he has risen from the dead, he will, by his power, bring all men to stand before him; for if \textit{he} has risen from the dead the bands of the temporal death are broken that the grave has no victory.
						\hypertarget{EMS-1834-JS-CIR-MAR-1834-jj}%
					If
						\marginnote{jj}%
					then, the grave has no victory, those who keep the sayings of Jesus and obey his teachings have, not only a promise of a resurrection from the dead; but an assurance of being admitted into his glorious kingdom; for, he himself says, Where I am, there shall also my servant be. -[see John \textsc{xii}.]- In the twenty second chapter of Luke's account of the Messiah, we find the kingdom of heaven likened unto a king who made a marriage for his son. That this son was the Messiah will not be disputed, since it was the \textit{kingdom of heaven} that was represented in the parable;
						\hypertarget{EMS-1834-JS-CIR-MAR-1834-kk}%
					and
						\marginnote{kk}%
					that the saints, or those who are found faithful to the Lord, are the individuals who will be found worthy to inherit a seat at the marriage-supper, is evident from the sayings of John in the Revelations, where he represents the sound which he heard in heaven to be like a great multitude, or like the voice of mighty thunderings, saying, The Lord God Omnipotent reigneth. Let us be glad and rejoice, and give honor to him; for the marriage of the Lamb is come, and his wife hath made herself ready. And to her was granted that she should be arrayed in fine linen, clean and white: for the fine linen is the righteousness of saints. -[Rev. \textsc{xix}.]-

						\hypertarget{EMS-1834-JS-CIR-MAR-1834-ll}%
					That
						\marginnote{ll}%
					those only are the individuals who keep the commandments of the Lord and walk in his statutes to the end, that are permitted to set at this glorious feast, is evident from the following items: In Paul's last letter to Timothy, which was written just previous to his death, he says, I have fought a good fight, I have finished my course, I have kept the faith: henceforth there is laid up for me a crown of righteousness which the Lord, the righteous Judge shall give me at that day: and not to me only, but unto all them also that love his appearing. No one who believes the account, will doubt for a moment this assertion of Paul which was made, as he knew, just before he was to take his leave of this world.
						\hypertarget{EMS-1834-JS-CIR-MAR-1834-mm}%
					Though
						\marginnote{mm}%
					he once, according to his own word, persecuted the church of God and wasted it, yet after embracing the faith, his labors were unceasing to spread the glorious news; and like a faithful soldier, when called to give his life in the cause which he had espoused, he laid it down, as he says, with an assurance of an eternal crown. Follow the labors of this apostle from the time of his conversion to the time of his death, and you will have a fair sample of industry and patience in promulgating the gospel of Christ: Whipped, stoned, and derided, the moment he escaped the hands of his persecutors, he, as zealously as ever, proclaimed the doctrine of the Savior. And all may know, that he did not embrace the faith for the honor of this life, nor for the gain of earthly goods.
						\hypertarget{EMS-1834-JS-CIR-MAR-1834-nn}%
					What
						\marginnote{nn}%
					then could have induced him to undergo all this toil? It was, as he said, that he might obtain that crown of righteousness from the hand of God. No one, we presume, will doubt the faithfulness of Paul to the end: None will say, that he did not keep the faith, that he did not fight the good fight, that he did not preach and persuade to the last: And what was he to receive? A crown of righteousness. And what shall others receive who do not labor faithfully, and continue to the end? We leave such to search out their own promises if any they have; and if they have any they are welcome to them, on our part, for the Lord says, that every man is to recive according to his works.
						\hypertarget{EMS-1834-JS-CIR-MAR-1834-oo}%
					Reflect
						\marginnote{oo}%
					for a moment, brethren, and enquire, whether you would consider yourselves worthy a seat at the marriage feast with Paul and others like him, if you had been unfaithful? Had you not fought the good fight, and kept the faith, could you expect to receive; have you a promise of receiving a crown of righteousness from the hand of the Lord, with the church of the first born? Here then, we understand, that Paul rested his hope in Christ because he had kept the faith, and loved his appearing; and from his hand he had a promise of receiving a crown of righteousness. If the saints are not to reign, for what purpose are they crowned?
						\hypertarget{EMS-1834-JS-CIR-MAR-1834-pp}%
					In
						\marginnote{pp}
					an exhortation of the Lord to a certain church in Asia, which was built up in the days of the apostles, unto whom he communicated his word on that occasion by his servant John, he says, Behold I come quickly: hold that fast which thou hast, that no man take thy crown. And again, To him that overcometh will I grant to sit with me in my throne, even as I also overcame, and am set down with my Father in his throne. -[see Rev \textsc{iii}.]- And again, it is written, Behold, now are we the sons of God, and it doth not appear what we shall be: but we know, that when he shall appear, we shall be like him; for we shall see him as he is.--- And he that hath this hope in him, purifieth himself, even as he is pure. -[1 John. \textsc{iii}, 2 \& 3.]-
						\hypertarget{EMS-1834-JS-CIR-MAR-1834-qq}%
					How
						\marginnote{qq}%
					is it that these old apostels should say so much on the subject of the coming of Christ? He certainly had once came; but Paul says, To all who love his appearing, shall be given the crown: and John says, When he shall appear, we shall be like him; for we shall see him as he is. Can we mistake such language as this? Do we not offer violence to our own good judgment when we deny the second coming of the Messiah? When has he partook of the fruit of the vine new with his ancient apostles in his Father's kingdom, as he said, just before he was crucified?
						\hypertarget{EMS-1834-JS-CIR-MAR-1834-rr}%
					In
						\marginnote{rr}%
					Paul's epistle to the Philippians, \textsc{iii}, 20 \& 21, he says, For our conversation is in heaven; from whence also we look for the Savior, the Lord Jesus Christ; who shall change our vile body, that it may be fashioned like unto his glorious body, according to the working whereby he is able even to subdue all things unto himself. We find another promise to individuals living in the church at Sardis, -[see Rev. \textsc{iii}. 4 \& 5.]- which will also show something of the blessings held out to the ancients who walked worthily before the Lord, which says, Thou hast a few names even in Sardis which have not defiled their garments; and they shall \textsc{walk with me in white}; for they are worthy.
						\hypertarget{EMS-1834-JS-CIR-MAR-1834-ss}%
					He
						\marginnote{ss}%
					that overcometh, the same shall be clothed in \textsc{white} raiment and I will not blot out his name out of the book of life; but I will confess his name before my Father, and before his angels. John represents the sound which he heard from heaven, as giving thanks and glory to God, saying that the Lamb was worthy to take the book, and to open its seals; beacuse he was slain, and had by his blood redeemed them out of every kindred and tongue, and people, and nation; and had made them kings and priests unto God: and they should reign on the earth. -[see Rev. \textsc{v}.]- In the twentieth chapter we find a length of time specified, during which Satan is to be confined in his own place, and the saints reign in peace.
						\hypertarget{EMS-1834-JS-CIR-MAR-1834-tt}%
					All
						\marginnote{tt}%
					these promises and blessings we find contained in the law of the Lord, which the righteous are to enjoy; and we might enumerate many more places where the same or similar promises are made to the faithful, but we do not deem it of importance to rehearse them here, as this letter is now lengthy; and our brethren no doubt, are familiar with them all. Most assuredly it is, however, that the ancients, though persecuted and afflicted by men, obtained from God promises of such weight and glory, that our hearts are often filled with gratitude, that we are even permitted to look upon them, while we contemplate that there is no respect of persons in \textsc{his} sight, and that in every nation, he that feareth him and worketh righteousness, is accepted with him.
						\hypertarget{EMS-1834-JS-CIR-MAR-1834-uu}%
					But
						\marginnote{uu}%
					from the few items previously quoted, we can draw a conclusion, that there is to be a day when all will be judged of their works, and rewarded according to the same; that those who have kept the faith will be crowned with a crown of righteousness; be clothed in white raiment; be admitted to the marriage-feast; be free from every affliction, and reign with Christ on the earth, where, according to the ancient promise, they will partake of the fruit of the vine new in the glorious kingdom with him: at least we find that such promises were made to the ancient saints.
						\hypertarget{EMS-1834-JS-CIR-MAR-1834-vv}%
					And
						\marginnote{vv}%
					though we cannot claim these promises which were made to the ancients, or that they are not our property merely because they were made to them, yet if we are the children of the most High, and are called with the same calling with which they were called, and embrace the same covenant that they embraced, and are faithful to the testimony of our Lord as they were, we can approach the Father in the name of Christ as they approached him, and for ourselves obtain the same promises. These promises, when obtained, if ever by us, will not be because Peter, John, and the other apostles, with the churches at Sardis, Purgamos, Philadelphia, and elsewhere, walked in the fear of God and had power and faith to prevail and obtain them;
						\hypertarget{EMS-1834-JS-CIR-MAR-1834-ww}%
					but
						\marginnote{ww}%
					it will be because \textit{we}, \textit{ourselves}, have faith and approach him in the name of his Son Jesus Christ, even as they did; and when these promises are obtained, they will be promises directly to us, or they will do us no good: communicated for \textit{our} benefit; being our own property, (through the gift of God,) earned by our own diligence in keeping his commandments, and walking uprightly before him. If not, to what end serves the gospel of our Lord Jesus Christ, and why was it ever communicated to us?

					Previous to commencing this letter we designed giving you some instruction upon the regulation of the church; but that will be given hereafter. -[\textsc{to be continued}.]-
					
		% -----
		\subsubsection{Minutes, 17 March 1834}
		% -----
			\label{EMS-1834-JS-17-MAR-1834}
			% \label{EMS-1830-JS-DAY-3LETTERMONTH-YEAR}
			
			\begin{description}
					\item[Print Sources]
				\end{description}

					\begin{tabular}{p{0.5in}p{1in}p{0.5in}p{1in}}
					JSP	 	& D3:484--488		& JSR 		& --- \\
					DC	 	& --- 		& HC 		& 2:44 \\
					HDDC 	& ---		& TCHC 		& 2:39--40 \\
					\end{tabular}
					% HDDC = woodford's DC diss; JSR = Marquardt's JS Revelations; TCHC = Vogel HC
				
				\begin{description}
					\item[Online Access] \href{https://www.josephsmithpapers.org/paper-summary/minutes-17-march-1834/1}{Here}
					% \item[Analysis] \hyperref[XX]{Here}
				\end{description}
				
				\begin{description}
					\item[Background]
				\end{description}
				
					% It might also be useful to give a two sentence context of the period: e.g., "This speech was given in Nauvoo to a small congregation one month prior to the destruction of the Nauvoo Expositor. These and other tensions may have been on the prophet's mind when he gave the speech."
				
				\begin{description}
					\item[Original Source]
				\end{description}
				
					% brief description of original source material, with reference to JSP or somewhere else for more info
					% Background material: it might be helpful to have a note that says whether a given document was presented as a speech, was written in a journal, etc.
					% Also a comment here about how reliable the source might be, or what shape it is in, if there are large omissions or parts that are hard to understand, and so on.
				
				\begin{description}
					\item[Text]
				\end{description}
				
					\begin{tightright}
							\hypertarget{EMS-1834-JS-17-MAR-1834-a}%
						Avon
							\marginnote{a}%
						Livingston Co. N. York.
						
						March 17, 1834.
					\end{tightright}
					
					This day a Conference of Elders assembled at the house Alvah Beeman [Beman]. Joseph Smith Jun. Sidney Rigdon Parley [P.] Pratt Lyman Wight John Murdock Orson Pratt and Orson Hyde, high priests, and Roger Orton, Isaac Mc. Withey [McWithy], Joseph Young, Harry Brown, Freeman Nickerson and Henry Shibly, Elders, were present. Bro. J Smith Jun. opened the conference by prayer. He then arose and introduced the object of our meeting, which was to obtain young men and middle aged to go and assist in the redemption of Zion according to the commandment, and for the church to gather up their riches and send them to purchase lands according to the commandments of the Lord.
						\hypertarget{EMS-1834-JS-17-MAR-1834-b}%
					Also,
						\marginnote{b}%
					to devise means, or obtain moneys for the relief of the brethren in Kirtland; say Two Thousand Dollars, which sum will deliver Kirtland from Debt for the present, and also to determine the course which the several shall pursue or journey when they leave this place. It was proposed by bro. Joseph that father [Edmund] Bosley and bro. Mc. Withey go with him to bro. Perry's and see if, by their united efforts, they could not raise Two Thousand Dollars for the relief of Kirtland. It was voted, that, bro. R. Orton, father Bosley, father Nickerson and father Mc. Withey, should exert themselves to obtain the said Two Thousand Dollars for the present relief of Kirtland.
						\hypertarget{EMS-1834-JS-17-MAR-1834-c}%
					They
						\marginnote{c}%
					all agreed to do all they could to obtain it; and they firmly believed that they could obtain the amount by the first of April. It was also voted that I should tarry and preach in the regions round about until the money could be obtaid, and bring it immediately to Kirtland. Voted that bro's. Joseph Sidney \& Lyman Wight go to Kirtland soon. Bro's. John Murdock and Orson Pratt were appointed to journey to Kirtland and preach by the way. Bro's. Parley Pratt \& Harry Brown\sout{, go} were appointed to visit the Churches in Black River country and obtain all the means they could to help Zion.
					
					\begin{tightright}
						Orson Hyde, Cl'k.
					\end{tightright}
					
		% -----
		\subsubsection{Record, 20 March 1834}
		% -----
			\label{EMS-1834-JS-20-MAR-1834}
			% \label{EMS-1830-JS-DAY-3LETTERMONTH-YEAR}
			
			\begin{description}
					\item[Print Sources]
				\end{description}

					\begin{tabular}{p{0.5in}p{1in}p{0.5in}p{1in}}
					JSP	 	& J1:35		& JSR 		& --- \\
					DC	 	& --- 		& HC 		& --- \\
					HDDC 	& ---		& TCHC 		& --- \\
					\end{tabular}
					% HDDC = woodford's DC diss; JSR = Marquardt's JS Revelations; TCHC = Vogel HC
				
				\begin{description}
					\item[Online Access] \href{https://www.josephsmithpapers.org/paper-summary/journal-1832-1834/64}{Here}
					% \item[Analysis] \hyperref[XX]{Here}
				\end{description}
				
				\begin{description}
					\item[Background]
				\end{description}
				
					% It might also be useful to give a two sentence context of the period: e.g., "This speech was given in Nauvoo to a small congregation one month prior to the destruction of the Nauvoo Expositor. These and other tensions may have been on the prophet's mind when he gave the speech."
				
				\begin{description}
					\item[Original Source]
				\end{description}
				
					% brief description of original source material, with reference to JSP or somewhere else for more info
					% Background material: it might be helpful to have a note that says whether a given document was presented as a speech, was written in a journal, etc.
					% Also a comment here about how reliable the source might be, or what shape it is in, if there are large omissions or parts that are hard to understand, and so on.
				
				\begin{description}
					\item[Text]
				\end{description}
				
					\textbf{Thursday 20\supers{th} Started on <our> Journy at noon took dinner at Brother Joseph Holbrooks, and at night tryed three times to git keept in the name of Deciples, and could not be keept, after night we found a man who would keep us for mony thus we see that there <is> more place for mony than for Jesus <Deciples or> the Lamb of God, the name of the man is \sout{Wilson Rauben Wilson} Reuben Wilson that would not keep us without mony <he lived in China> \&c.}
					
		% -----
		\subsubsection{Letter to Edward Partridge and Others, 30 March 1834}
		% -----
			\label{EMS-1834-JS-30-MAR-1834}
			% \label{EMS-1830-JS-DAY-3LETTERMONTH-YEAR}
			
			\begin{description}
					\item[Print Sources]
				\end{description}

					\begin{tabular}{p{0.5in}p{1in}p{0.5in}p{1in}}
					JSP	 	& D3:488--498		& JSR 		& --- \\
					DC	 	& --- 		& HC 		& --- \\
					HDDC 	& ---		& TCHC 		& --- \\
					\end{tabular}
					% HDDC = woodford's DC diss; JSR = Marquardt's JS Revelations; TCHC = Vogel HC
				
				\begin{description}
					\item[Online Access] \href{https://www.josephsmithpapers.org/paper-summary/letter-to-edward-partridge-and-others-30-march-1834/1}{Here}
					% \item[Analysis] \hyperref[XX]{Here}
				\end{description}
				
				\begin{description}
					\item[Background]
				\end{description}
				
					% It might also be useful to give a two sentence context of the period: e.g., "This speech was given in Nauvoo to a small congregation one month prior to the destruction of the Nauvoo Expositor. These and other tensions may have been on the prophet's mind when he gave the speech."
				
				\begin{description}
					\item[Original Source]
				\end{description}
				
					% brief description of original source material, with reference to JSP or somewhere else for more info
					% Background material: it might be helpful to have a note that says whether a given document was presented as a speech, was written in a journal, etc.
					% Also a comment here about how reliable the source might be, or what shape it is in, if there are large omissions or parts that are hard to understand, and so on.
				
				\begin{description}
					\item[Text]
				\end{description}
				
						\hypertarget{EMS-1834-JS-30-MAR-1834-a}%
					To
						\marginnote{a}
					Edward [Partridge], Williams [William W. Phelps], and others of the [United] firm.
				
					\begin{tightright}
						Kirtland, March 30, 1834.
					\end{tightright}

					Dear Brethren:

					We have received several communications from you of late; but the most of us being absent, brother Oliver [Cowdery] laid them over till council could be had; and I now seat myself to dictate to answer them all in one. Since brothers, Parley [P. Pratt] \& Lyman [Wight], arrived I have written a few lines with my own hand in letters which have already gone: one from this place, and one from Freedom N.Y. but was not able to write the more weighty matters, and did not think to say anything more than to comfort your hearts if possible, and keep you from fainting, while God, in his wisdom, and in the order of his Providence, is preparing all things before his face for the redemption of Zion. We rejoic greatly on learning that you and the brethren, so many of them, are yet spared in the midst of those who wear the form of human beings, but are less merciful than the prowling beast of the wilderness.
						\hypertarget{EMS-1834-JS-30-MAR-1834-b}%
					We
						\marginnote{b}%
					would inform you that with very few exceptions the Church in this place are all well; and every man, woman \& child, that belongs to the Church, as far as I have any knowledge of the matter, are crying day \& night to God for the deliverance and prosperity of Zion; and many are preparing with all zeal to do all that lies in their power to accomplish the great work, and it will be seen in due time, that the saints in this region are not slack towards you considering their circumstances, \& their great poverty, \& afflictions \& persecutions with which they are called to suffer in this part as well as you in that region; for the more we try to live Godly in Christ Jesus, the more we are made to feel the weight of persecution, inflicted by those who are under the influence of the enemy of the souls of men.
						\hypertarget{EMS-1834-JS-30-MAR-1834-c}%
					But
						\marginnote{c}%
					let this suffice: I shall proceed first to answer some of the most important items contained in your last communications, the more part which gave us much satisfaction. We admire the confidence \& love which our brethren have manifested in them, in giving us \uline{sharp}, \uline{piercing}, \& \uline{cutting} reproofs, which are calculated to wake us up \& make us search about ourselves, \& put a double watch over ourselves in all things that we do. And we acknowledge that it is our duty to receive all reproofs \& chastisements given of the spirit of the most Holy One.
						\hypertarget{EMS-1834-JS-30-MAR-1834-d}%
					And
						\marginnote{d}%
					if being chastised and reproved of what we are guilty, seems not to be joyous for the present but grievous, O, how wounding, \& how poignant must it be to receive chastisements \& reproofs, for things that we are not guilty of from a source we least expect them, arising from a distrustful, a fearful, \& jealous spirit. However, we feel to make all allowances, \& reflect seariously \& consider upon all sides before we make an effort to throw off the yoke, lest we should be found in anywise blamable before God.
						\hypertarget{EMS-1834-JS-30-MAR-1834-e}%
					There
						\marginnote{e}%
					are some items contained in bro. William's letters by the way of reproof, that we feel to give, we think some reasonable excuses, that you may <know> how far you have reasons to give reproof, that you may not have wrong feelings concerning those to whom you are espoused in Christ Jesus who always will be found \uline{true} to all confidence that shall be imposed in them.

						\hypertarget{EMS-1834-JS-30-MAR-1834-f}%
					Firstly,
						\marginnote{f}%
					you have given us to understand that there are glaring errors in the Revelation, or rather, have shown us the most glaring ones, which are not calculated to suit the refinement of the age in which we live, of the great men, \&c. We would say, by way of excuse, that we did not think so much of the orthography, or the manner, as as we did of the subject matter; as the word of God means what it says; \& it is the word of God, as much as Christ was God, although \uline{he} was born in a stable, \& was rejected by the manner of his birth, notwithstanding he was God. \uline{What a mistake}! the manner of his birth, \& the source from which he sprang caused him to be rejected \& cast out, \& to be taken \& put to death.

						\hypertarget{EMS-1834-JS-30-MAR-1834-g}%
					Whereas
						\marginnote{g}%
					<had> he pleased the great men, the high priests, the lawyers, \& the learned, he might have escaped. But supposing we should happen to make as great a mistake as the Lord did, \& come under the censure of big men \& fall in the same way, what would be the consequence? The fact was, there was no room in the \uline{Inn}; \& when man cannot do as they would, they must do as they can; for God set the example before them. \uline{For there was no room in the Inn}! but there was room found ``in the \uline{stable}; \& here was utterly a fault in the eyes of the laughing philosophers;'' but it is not given to us to understand that he altered his course to please any man.
						\hypertarget{EMS-1834-JS-30-MAR-1834-h}%
					And
						\marginnote{h}%
					who was it that triumphed? was it the ``laughing philosophers,'' or \uline{him} who never deviated from the will of him who sent him? Now the fact is, if we have made any mistakes in punctuation, or spelling, it has been done in consequence of brother Oliver, having come from Zion in great afflictions, through much fatigue and anxiety, and being sent contrary to his expectations to New York, and obtain[in]g press and Types, and \uline{hauling} them up in the midst of \uline{mobs}, when he and I, and all the church in Kirtland had to lie every night for a long time upon our arms to keep off \uline{mobs}, of forties, of eighties, \& of hundreds to save our lives and the press, and that we might not be scattered \& driven to the four winds!
						\hypertarget{EMS-1834-JS-30-MAR-1834-i}%
					And
						\marginnote{i}%
					all this in the midst of every kind of confusion \& calamity \& in the sorrowful tale of Zion, for the sake of Zion, that the word of God might be printed \& sent forth by confidential brethren to the different churches; for the churches are just like you--- they will not receive anything but by revelation! for when you \uline{hint} they will ask a question, and if by any means in the heat of zeal you would hit them a \uline{kick} it never fails to turn over the \uline{dish}.
						\hypertarget{EMS-1834-JS-30-MAR-1834-j}%
					Therefore,
						\marginnote{j}%
					when we give them a hint, and they ask a question, we sometimes answer them plainly; but all this is a wonder and a mystery; but it wont do to \uline{kick}, therefore to unfold the mystery we must of necessity send out the word of \sout{the mystery we must of} God unto the different churches, or they could not be made to understand, that they, with their moneys, and their young men \& their middle aged, must, in order to do the will of God, redeem the land which had been purchased, \& the children of Zion--- and if by chance in doing all this <we> should have to suffer peril by false brethren.
						\hypertarget{EMS-1834-JS-30-MAR-1834-k}%
					For
						\marginnote{k}%
					men are as liable in this generation to turn aside from the holy commandments, as were the children of Israel when Aaron bought the golden calf at the expense of all the jewelry, \& riches of the children of Israel, while Moses tarried yet forty days in the mount, that he might receive the law of the everlasting gospel upon tables of stone, written by the finger of of God, while they, the children of Israel, were delivered over, \& bowed down and worshiped the dumb idol, and said, These be our Gods that brought us up out of the land of Egypt. And Moses being angry destroyed the tables of Stone, and the golden calf and made the children of Isarel drink the substance of their God, which they said brought them up out of the land of Egypt. Therefore, I say, if we should suffer peril among false brethren, should it be accounted a strange thing?
						\hypertarget{EMS-1834-JS-30-MAR-1834-l}%
					But
						\marginnote{l}%
					here comes up another question, a great mystery! How did the revelation come to be garbled by the printers of the day, published and sent to Jackson Co. and elsewhere? But if all these things, upon a little reflection had been rightly considered and understood, there would have been no mystery, nor any question asked. Not a sigh--- not a lingering thought--- not a grief, or a single reflection cast upon the innocent, a virgin, \uline{the spouse of Zion}! Suffice it to say, that the revelation went into the hands of the world by stealth, through the means of \uline{false} brethren, and lest it should reach the ears of the President and Governor, with a false coloring, being misrepresented, wisdom dictated that we should send it in its own proper light.
						\hypertarget{EMS-1834-JS-30-MAR-1834-m}%
					And
						\marginnote{m}%
					if truth, and the word of God will not bear off the Palms and bring us the victory, shall we, who profess to be men of God condescend to folly? Shall we turn aside from the word of God and seek to save our lives, and that we may please men? If men will seek occasion against the truth, will they not seek occasion even if we should shun the truth? The fact is, beloved brethren, we seek not gold or silver or this world's goods, nor honors nor the applause of men; but we seek to please him, and to do the will of him who hath power not only to destroy the body; but to cast the soul into hell! Ah! men should not attempt to steady the ark of God! But enough on this subject.
			
						\hypertarget{EMS-1834-JS-30-MAR-1834-n}%
					Now
						\marginnote{n}%
					concerning employing Mr. [Robert W.] Wells of Jefferson C'ty. as Counsellor \&c. We think it would be advisable. You may consider that you have our consent: We speak to wise men! Judge ye what we say! Employ, then, Mr. Wells, and although we have neither gold nor silver, we have run into debt for the press, and also to obtain money to pay the New York debt for Zion, and have received but a very few dollars for the Star and printing as yet, no means of speculation to gain or make money, yet we think that the money can be had, and that there will be no difficulty on this subject: and this, while you are writing to us to reprove us, and telling us, that your dependence for money is on your eastern brethren, and at the same time saying ``\uline{Dont} \uline{buy} \uline{your} \uline{gold} \uline{too dear}!'' this is the way that we buy our gold!
						\hypertarget{EMS-1834-JS-30-MAR-1834-o}%
					Now,
						\marginnote{o}%
					brethren, let me tell you, that it is my disposition to give and \uline{forgive}, and to bear \& to \uline{forbear}, with all long suffering and patience, with the foibles, follies, weaknesses, \& wickedness of my brethren and all the world of mankind; and my confidence and love toward you is not slackened, nor weakened. And now, if you should be called upon to bear with us a little in any of our weaknesses and follies, and should, with us, receive a rebuke to \uline{yourselves}, dont be offended, dont in anywise let it \uline{hit} you, so as to turn over the dish! And when you \& I meet face to face, I anticipate, without the least doubt, that all matters between us will be fairly understood, and perfect love prevail; and sacred covenant by which we are bound together, have the uppermost seat in our hearts.

						\hypertarget{EMS-1834-JS-30-MAR-1834-p}%
					We expect that a number of our able
						\marginnote{p}%
					brethren will come on soon and go to Zion; and should you have no other way of obtaining moneys, you can sell them your lands, let them go on to them, protect them on the same, till your suits are determined, and then, (if you succeed) you will have means to purchase more, and if not they will receive you into their bosoms. We see no other way now; but the Lord may open other ways in time. Brs. Parley and Lyman are both in the east; but we expect they will leave here for the west by the first of may, and go as soon as they can, so should you be organized by the time they arrive, perhaps it would be well.
						\hypertarget{EMS-1834-JS-30-MAR-1834-q}%
					You
						\marginnote{q}%
					must act wisdom for yourselves in many things, as you are better prepared to judge in many things than we are. Many things are familiar with some of us, which we \uline{cannot} communicate by letter; but will be brought about in their times. You ought to be prepared to go back at a moment's warning, and we are inclined to think that it was a wise step in employing the Att'y Gen. for he will investigate and learn the truth, and then Governor will investigate also.

						\hypertarget{EMS-1834-JS-30-MAR-1834-r}%
					Once
						\marginnote{r}%
					more I design coming unto [you?]; but when, it has not been revealed: whether it will be with Parley \& Lyman I cannot now say; but \uline{once more} I design to come \uline{mob} or no \uline{mob}, \uline{enemy} or no \uline{enemy}! There needs be no difficulty in relation to the revelations; for they show plainly from the face of them, that no blood is to be shed except in self-defense; and that the law of God as well as man gives us a privilege. If you make yourselves acquainted with the revelations, you will see that this is the case, though we should not publish any more than we are obliged to of necessity for the Church's sake.
						\hypertarget{EMS-1834-JS-30-MAR-1834-s}%
					We
						\marginnote{s}%
					have nothing to fear if we are faithful: God will strike through kings in the day of his wrath but what he will deliver his people; and what do you suppose he could do with a few \uline{mobbers} in Jackson County, where, ere long, he will set his feet, when earth \& heaven shall tremble!

					Be united, brethren, in all your moves, and stand by each other even unto death that you may prevail.

					I remain your brother in the new covenant.

					\begin{tightright}
						Joseph Smith Jun...
					\end{tightright}
					
		% -----
		\subsubsection{Record, 1 April 1834}
		% -----
			\label{EMS-1834-JS-1-APR-1834}
			% \label{EMS-1830-JS-DAY-3LETTERMONTH-YEAR}
			
			\begin{description}
					\item[Print Sources]
				\end{description}

					\begin{tabular}{p{0.5in}p{1in}p{0.5in}p{1in}}
					JSP	 	& J1:37		& JSR 		& --- \\
					DC	 	& --- 		& HC 		& --- \\
					HDDC 	& ---		& TCHC 		& --- \\
					\end{tabular}
					% HDDC = woodford's DC diss; JSR = Marquardt's JS Revelations; TCHC = Vogel HC
				
				\begin{description}
					\item[Online Access] \href{https://www.josephsmithpapers.org/paper-summary/journal-1832-1834/68}{Here}
					% \item[Analysis] \hyperref[XX]{Here}
				\end{description}
				
				\begin{description}
					\item[Background]
				\end{description}
				
					% It might also be useful to give a two sentence context of the period: e.g., "This speech was given in Nauvoo to a small congregation one month prior to the destruction of the Nauvoo Expositor. These and other tensions may have been on the prophet's mind when he gave the speech."
				
				\begin{description}
					\item[Original Source]
				\end{description}
				
					% brief description of original source material, with reference to JSP or somewhere else for more info
					% Background material: it might be helpful to have a note that says whether a given document was presented as a speech, was written in a journal, etc.
					% Also a comment here about how reliable the source might be, or what shape it is in, if there are large omissions or parts that are hard to understand, and so on.
				
				\begin{description}
					\item[Text]
				\end{description}
				
					...\textbf{<this> is \sout{this} Aprel 1\supers{st} Tusday my Soul delighteth in the Law of the Lord for he forgiveth my sins and <will> confound mine Enimies the Lord shall destroy him who has lifted his heel against me even that wicked man Docter P Hrlbert [Doctor Philastus Hurlbut] he <will> deliver him to the fowls of heaven and his bones shall be cast to the blast of the wind <for> he lifted his <arm> against the Almity therefore the Lord shall destroy him}
					
		% -----
		\subsubsection{Letter to Orson Hyde, 7 April 1834}
		% -----
			\label{EMS-1834-JS-7-APR-1834}
			% \label{EMS-1830-JS-DAY-3LETTERMONTH-YEAR}
			
			\begin{description}
					\item[Print Sources]
				\end{description}

					\begin{tabular}{p{0.5in}p{1in}p{0.5in}p{1in}}
					JSP	 	& D4:5--9		& JSR 		& --- \\
					DC	 	& --- 		& HC 		& 2:48--49 \\
					HDDC 	& ---		& TCHC 		& 2:45--46 \\
					\end{tabular}
					% HDDC = woodford's DC diss; JSR = Marquardt's JS Revelations; TCHC = Vogel HC
				
				\begin{description}
					\item[Online Access] \href{https://www.josephsmithpapers.org/paper-summary/letter-to-orson-hyde-7-april-1834/1}{Here}
					% \item[Analysis] \hyperref[XX]{Here}
				\end{description}
				
				\begin{description}
					\item[Background]
				\end{description}
				
					% It might also be useful to give a two sentence context of the period: e.g., "This speech was given in Nauvoo to a small congregation one month prior to the destruction of the Nauvoo Expositor. These and other tensions may have been on the prophet's mind when he gave the speech."
				
				\begin{description}
					\item[Original Source]
				\end{description}
				
					% brief description of original source material, with reference to JSP or somewhere else for more info
					% Background material: it might be helpful to have a note that says whether a given document was presented as a speech, was written in a journal, etc.
					% Also a comment here about how reliable the source might be, or what shape it is in, if there are large omissions or parts that are hard to understand, and so on.
				
				\begin{description}
					\item[Text]
				\end{description}
				
					\begin{tightright}
							\hypertarget{EMS-1834-JS-7-APR-1834-a}%
						Kirtland
							\marginnote{a}
						April 7th 1834
					\end{tightright}
					
					Dear bro Orson [Hyde]

					We received yours of the 31st ult in due course of Mail and were much grieved on learning that you were not like to succeed according to our expectations, Myself bro Newel [K. Whitney] Frederick [G. Williams] and Oliver [Cowdery] retired to the \uline{Translating} \uline{room} when prayer was wont to be made and unbosomed our feelings before God and cannot but exercise faith yet that you in the meraculus providence of God will succeed in obtaining help the fact is \sout{that} unless we can obtain help I myself cannot go to Zion and if I do not go it will be impossable to get my brethren in Kirtland any of them to go and if we do not go it is in vain for our eastern brethren to think of going up to better themselves by obtaining so goodly a land which now can be obtained for one dollar and a quarter pr acre and stand against that wicked \uline{Mob} for unless they do the will of God, God will not help them and if God does not help them all is \uline{vain}.
						\hypertarget{EMS-1834-JS-7-APR-1834-b}%
					Now
						\marginnote{b}%
					the fact is this is the head \sout{is} of the church and the life of the body and threu able men as members of the body God has appointd to be hands to administer to the necessities of the body, Now if a mans hands refuse to administer to the necessity of his body it must perish of hunger and if the body perish all the members perish with it and if the head fails the whole body is sickened, the heart faints and the body dies the spirit takes its exit and the carcass remains to be devoured of worms

						\hypertarget{EMS-1834-JS-7-APR-1834-c}%
					Now
						\marginnote{c}%
					bro Orson if this Church which is assaying to be the church of Christ will not help us when they can do it without sacrifice with those blessings which God has bestowed upon them I proph[es]y I speak the truth I Lie not God shall take away their tallant and give it to those who have no tallant and shall prevent them from ever obtaining a place of reffuge or an inheritance upon the Land of Zion
						\hypertarget{EMS-1834-JS-7-APR-1834-d}%
					Therefore
						\marginnote{d}%
					they may tarry for they might as well be overtaken where they are as to incur the displeasur of God and fall under his wrath by the wayside or to fall into the hands of a merciless \uline{mob} where there is no God to deliver as salt that has lost its savour and is thenceforth \sout{god} good for nothing but to be troden undr feet of men I therefor beseech you to conjure [adjure] them in the name of the Lord by the Love of God to lend us a helping hand and if all this will not soften their hearts to administer to our necessity for Zion sake
						\hypertarget{EMS-1834-JS-7-APR-1834-e}%
					turn
						\marginnote{e}%
					your back up on them and return speedily to Kirtland and the blood of Zion be upon their heads even as upon the heads of her enimies and let their reccompence be as the reccompence of her enemes for thus shall it come to pass saith the Lord of hosts who has the cattle upon a thousand hills who has put forth his Almyhty [Almighty] hand to bring to pass his strang[e] act and what man shall put forth his hand to steady the ark of God or be found turning a deaf ear to the voice of his servant God shall speak in due time and all will be declared amen

						\hypertarget{EMS-1834-JS-7-APR-1834-f}%
					Your
						\marginnote{f}%
					broth[er] in the new covenant Joseph

					PS I am much disappointed on learni[n]g about my horse but if you cannot obtain him bring the mare and please do not [ride?] her very fast in a day and be very car[e]ful that I may not loose her and perhaps I may dispose of her to good advantage for the benefit of Zion--- I am yours by every sacred and holy tie that can bind up souls

					\begin{tightright}
					Joseph Smith Jr

					F G Williams

					Oliver Cowde[r]y
					\end{tightright}
					
		% -----
		\subsubsection{Record, 18--19 April 1834}
		% -----
			\label{EMS-1834-JS-18-19-APR-1834}
			% \label{EMS-1830-JS-DAY-3LETTERMONTH-YEAR}
			
			\begin{description}
					\item[Print Sources]
				\end{description}

					\begin{tabular}{p{0.5in}p{1in}p{0.5in}p{1in}}
					JSP	 	& J1:41--42		& JSR 		& --- \\
					DC	 	& --- 		& HC 		& --- \\
					HDDC 	& ---		& TCHC 		& --- \\
					\end{tabular}
					% HDDC = woodford's DC diss; JSR = Marquardt's JS Revelations; TCHC = Vogel HC
				
				\begin{description}
					\item[Online Access] \href{https://www.josephsmithpapers.org/paper-summary/journal-1832-1834/77}{Here}
					% \item[Analysis] \hyperref[XX]{Here}
				\end{description}
				
				\begin{description}
					\item[Background]
				\end{description}
				
					% It might also be useful to give a two sentence context of the period: e.g., "This speech was given in Nauvoo to a small congregation one month prior to the destruction of the Nauvoo Expositor. These and other tensions may have been on the prophet's mind when he gave the speech."
				
				\begin{description}
					\item[Original Source]
				\end{description}
				
					% brief description of original source material, with reference to JSP or somewhere else for more info
					% Background material: it might be helpful to have a note that says whether a given document was presented as a speech, was written in a journal, etc.
					% Also a comment here about how reliable the source might be, or what shape it is in, if there are large omissions or parts that are hard to understand, and so on.
				
				\begin{description}
					\item[Text]
				\end{description}
				
						\hypertarget{EMS-1834-JS-18-19-APR-1834-a}%
					...We
						\marginnote{a}%
					soon retired to the wilderness where we united in prayer and suplication for the blessings of the Lord to be given unto his church: We called upon the Father in the name of Jesus to go with the breth[r]en who were going up to the land of Zion, to give brother Joseph strength, and wisdom, and understanding sufficient to lead the people of the Lord, and to gather back, and establish the saints upon the land of their inheritances, and organize them according to the will of heaven, that they be no more cast down forever.
						\hypertarget{EMS-1834-JS-18-19-APR-1834-b}%
					We
						\marginnote{b}%
					then united and laid on hands: Brothers Sidney, Oliver, and Zebedee laid hands upon bro. Joseph, and confirmed upon him all the blessings necessary to qualify him to \sout{do} <stand> before the Lord in his high calling; and he return again in peace and triumph, to enjoy the society of his breth[r]en. Brothers Joseph, Sidney, and Zebedee then laid hands upon bro. Oliver, and confirmed upon him the blessings of wisdom and understanding sufficient for his station; that he be qualified to assist brother Sidney in arranging the church covenants which are to be soon published; and to have intelligence in all things to do the work of printing.
						\hypertarget{EMS-1834-JS-18-19-APR-1834-c}%
					Brother
						\marginnote{c}%
					Joseph, Oliver, Zebedee then laid hands upon bro. Sidney, and confirmed upon him the blessings of wisdom and knowledge to preside over the church in the abscence of brother Joseph, and to have the spirit to assist bro. Oliver in conducting the Star, and to arrange the Church covenants, and the blessing of old age and peace, till Zion is built up \& Kirtland established, till all his enemies are under his feet, and of a crown of eternal life \sout{at the} <in the> Kingdom of God with us.
						\hypertarget{EMS-1834-JS-18-19-APR-1834-d}%
					We,
						\marginnote{d}%
					Joseph, Sidney, and Oliver then laid hands upon bro. Zebedee, and confirmed the blessing of wisdom to preach the gospel, even till \sout{is} <it> spreads to the islands of the seas, and to be spared to see th[r]eescore years and ten, and see Zion built up and Kirtland established forever, and even at last to receive a crown of life. Our hearts rejoiced, and we were comforted with the Holy Spirit, Amen.
					
		% -----
		\subsubsection{Minutes and Discourse, 21 April 1834}
		% -----
			\label{EMS-1834-JS-21-APR-1834}
			% \label{EMS-1830-JS-DAY-3LETTERMONTH-YEAR}
			
			\begin{description}
					\item[Print Sources]
				\end{description}

					\begin{tabular}{p{0.5in}p{1in}p{0.5in}p{1in}}
					JSP	 	& D4:13--19		& JSR 		& --- \\
					DC	 	& --- 		& HC 		& 2:52--54 \\
					HDDC 	& ---		& TCHC 		& 2:57--61 \\
					TPJS 	& 70--71			& 	 		&  \\
					\end{tabular}
					% HDDC = woodford's DC diss; JSR = Marquardt's JS Revelations; TCHC = Vogel HC
				
				\begin{description}
					\item[Online Access] \href{https://www.josephsmithpapers.org/paper-summary/minutes-and-discourse-21-april-1834/1}{Here}
					% \item[Analysis] \hyperref[XX]{Here}
				\end{description}
				
				\begin{description}
					\item[Background]
				\end{description}
				
					% It might also be useful to give a two sentence context of the period: e.g., "This speech was given in Nauvoo to a small congregation one month prior to the destruction of the Nauvoo Expositor. These and other tensions may have been on the prophet's mind when he gave the speech."
				
				\begin{description}
					\item[Original Source]
				\end{description}
				
					% brief description of original source material, with reference to JSP or somewhere else for more info
					% Background material: it might be helpful to have a note that says whether a given document was presented as a speech, was written in a journal, etc.
					% Also a comment here about how reliable the source might be, or what shape it is in, if there are large omissions or parts that are hard to understand, and so on.
				
				\begin{description}
					\item[Text]
				\end{description}
				
					\begin{tightright}
							\hypertarget{EMS-1834-JS-21-APR-1834-a}%
						Norton
							\marginnote{a}
						Medina Co. Ohio April 21, 1834.
					\end{tightright}
					
					This day a conference of the elders of the church of Christ assembled at the dwelling house of bro. Carpenters at 10 o'clock A.M.
					
					Opened by singing, ``How firm a foundation, \&c.'' Bro. Joseph Smith Jun\supers{r.} read the 2\supers{nd.} chapter of the prophecy of Joel \& took the lead in prayer; after which, he commenced addressing the congregation, as follows. It is very difficult for us to communicate to the churches all that God has revealed to us, in consequence of tradition; for we are differently situated from any other people that ever existed upon this earth: Consequently those former revelations cannot be suited to our condition, because they were given to other people who were before us; but in the last days, God was to call a remnant, in which was to be deliverance, as well as in Jerusalem, and Zion.
						\hypertarget{EMS-1834-JS-21-APR-1834-b}%
					Now,
						\marginnote{b}%
					if God should give no more revelations, where will we find Zion and this remnant? He said that the time was near when desolation was to cover the Earth, and then God would have a place of deliverance in his remnant, and in Zion, \&c. He then gave a relation of obtaining and translating the Book of Mormon, the revelation of the priesthood of Aaron, the organization of the Church in the year 1830, the revelation of the high priesthood, and the gift of the Holy Spirit poured out upon the church, \&c.
						\hypertarget{EMS-1834-JS-21-APR-1834-c}%
					Take
						\marginnote{c}%
					away the book of Mormon, and the revelations, and where is our religion? We have none; for without a Zion and a place of deliverance, we must fall, because the time is near when the sun will be darkened, the moon turn to blood, the stars fall from heaven and the earth reel to and fro; then if this is the case, if we are not sanctified and gathered to the places where God has appointed, our former professions and our great love for the bible, we must fall, we cannot stand, we cannot be saved; for God will gather out his saints from the gentiles and then comes desolation or destruction and none can escape except the pure in heart who are gathered, \&c.
				
						\hypertarget{EMS-1834-JS-21-APR-1834-d}%
					...The
						\marginnote{d}%
					time was occupied for a few minutes by two or three others upon the same subject. Bro. Joseph Smith Jn\supers{r.} then deliverd a short prophecy, that if Zion was not deliverd the time was near when all of this church, whereever they might be found, would be persecuted and destroyed in like manner
					
						\hypertarget{EMS-1834-JS-21-APR-1834-e}%
					Bro.
						\marginnote{e}%
					Sidney Rigdon then took up the second item, viz:--- The endowment of the Elders with power from on high. He gave an account of the endowment of the ancient apostles and laid before the conference the dimensions of the house to be built in Kirtland and rehearsed the promise to the elders in the last days which they were to realize after the house of the Lord was built.
						\hypertarget{EMS-1834-JS-21-APR-1834-f}%
					Bro.
						\marginnote{f}%
					J. Bosworth then related a few items of a vision which he gave as a testimony of those things contained in the revelations read by Bro. Sidney, and his remarks upon that part relative to the endowment of the elders with power from on high. Bro. Joseph then occupied a few minutes by way of explanation of the revelation concerning the building of the house of the Lord...
					
					...Bro. Joseph then laid hands upon certain children \& blessed them in the name of the Lord....
					
		% -----
		\subsubsection{Revelation, 23 April 1834 [D\&C 104]}
		% -----
			\label{EMS-1834-JS-23-APR-1834}
			% \label{EMS-1830-JS-DAY-3LETTERMONTH-YEAR}
			
			\begin{description}
					\item[Print Sources]
				\end{description}

					\begin{tabular}{p{0.5in}p{1in}p{0.5in}p{1in}}
					JSP	 	& D4:19--31			& JSR 		& 255--260 \\
					DC	 	& Section 104 		& HC 		& 2:54--60 \\
					HDDC 	& 1348--1371		& TCHC 		& 2:61--64 \\
					\end{tabular}
					% HDDC = woodford's DC diss; JSR = Marquardt's JS Revelations; TCHC = Vogel HC
				
				\begin{description}
					\item[Online Access] \href{https://www.josephsmithpapers.org/paper-summary/revelation-23-april-1834-dc-104/1}{Here}
					% \item[Analysis] \hyperref[DC104]{Here}
				\end{description}
				
				\begin{description}
					\item[Background]
				\end{description}
				
					% It might also be useful to give a two sentence context of the period: e.g., "This speech was given in Nauvoo to a small congregation one month prior to the destruction of the Nauvoo Expositor. These and other tensions may have been on the prophet's mind when he gave the speech."
				
				\begin{description}
					\item[Original Source]
				\end{description}
				
					% brief description of original source material, with reference to JSP or somewhere else for more info
					% Background material: it might be helpful to have a note that says whether a given document was presented as a speech, was written in a journal, etc.
					% Also a comment here about how reliable the source might be, or what shape it is in, if there are large omissions or parts that are hard to understand, and so on.
				
				\begin{description}
					\item[Text]
				\end{description}
					
					\begin{tightcenter}
					April 23, 1834
					\end{tightcenter}
					
					Verily, I say unto you my friends, I give unto you counsel \& a commandment concerning all the properties which belong to the Firm, which I commanded to be organized \& established to be \sout{a} a United Firm, \& an everlasting Firm, for the benifit of my church, \& for the salvation of men until I come, with promise immutible \& unchangeable, that inasmuch as those whom I commanded were faithful, they should be blessed with multiplicity of blessings; but inasmuch as they were not faithful, they were nigh unto cursing. Therefore inasmuch as some of my servants have not kept the commandment but have broken the covenant, by coveteousness \& with feigned words, I have cursed them with a verry sore \& grievous curse; for I the Lord have decreed in my heart, that inasmuch as any <man> belonging to the Firm, shall be found a transgressor, or in other words, shall \sout{break} brake the covenant with which ye are bound, he shall be cursed in his life \& shall be trodden down by whom I will; for I the Lord am not to be mocked in these things; \& all this that the inocent among you may not be condemned with the unjust, \& that the guilty among you may not escape because I the Lord have promised unto you a crown of glory at my right hand. Therefore, inasmuch as ye are found transgressors, ye cannot escape my wrath in your lives; \& inasmuch as ye are cut off by transgression ye cannot escape the buffetings of Satan unto the day of Redemption. And I now give unto you power from this verry hour, that if any man among you, of the Firm, is found a transgressor, \& repenteth not of the evil, that ye shall deliver him over unto the buffetings of Satan, \& he shall have no more power to bring evil upon you; but as long as ye hold communion with transgressors, behold, they bring evil upon you. It is wisdom in me, therefore, a commandment I give unto you, that ye shall organize yourselves, \& appoint every man his stewardship, that every man may give an account unto me of the stewardship which is appointed unto him; for it is exped[i]ent, that I the Lord, should make every man accountable, as stewards over earthly Blessings, which I have made \& prepared for my creatures, I the Lord stretched out the heavens, \& builded the earth as a verry handy work, \& all things therein are mine, \& it is my business to provide for my saints, for all things are mine; but it must needs be done in mine own way: \& behold, this is the way that I the Lord hath decreed to provide for my saints, that the poor shall be exalted in that the rich are made low; for the earth is full, \& there is enough \& to spare; yea, I have prepared all things, \& have given unto the children of men to be agents unto themselves. Therefore if any man shall take of the abundance which I have made, \& impart not his portion according to the law of my gospel unto the poor \& the needy, he shall with Dives\sout{es} lift up his eyes in hell, being in torment. And now verily, I say unto you concerning the properties of the Firm, Let my servant Sidney [Rigdon] have appointed unto him the place where he now resides, \& the lot of the Tan[n]ery for his stewardship for his support while he is labouring in my vinyard, even as I will, when I shall command him; \& let all things be done according to counsel of the Firm, \& united consent, or voice of the Firm which dwells in the land of Kirtland. And this stewardship \& blessing, I the Lord confer upon my servant Sidney for a blessing upon him, \& upon his seed after him, \& I will multiply blessings upon him \& upon his seed after him inasmuch as he shall be humble before me. And again let my servant Martin [Harris] have appointed unto him for his stewardship the lot of land which my servant John [Johnson] obtained in exchange for his farm, for him \& his seed after him; \& inasmuch as he is faithful I will multiply blessings upon him \& his seed after him. And let my servant Martin devote his moneys for the printing of my word, according as my servant Joseph shall direct.

					And again let my servant Frederick [G. Williams] have the place upon which he now dwells; \& let my servant Oliver [Cowdery] have the lot which is set off joining the house which is to be for the printing office which is lot number one; \& also the lot upon which his father resides; \& let my servants Frederick \& Oliver have the printing office \& all things that pertain unto it; \& this shall be their stewardship which shall be appointed unto them; \& inasmuch as they are faithful, behold, I will bless them, \& multiply blessings upon them, \& this is the beginning of the stewardship which I have appointed unto them, for them \& their seed after them; \& inasmuch as they are faithful I will multiply blessings upon them \& their seed after them, even a multiplicity of blessings.

					And again, let my servant John have the house in which he lives, \& the farm, all, save the ground which has been reserved for the building of my houses, which pertains to that farm, \& those lots which have been named for my servant Oliver; \& inasmuch as he is faithful I will multiply blessings upon him. And it is my will that he should sell the lots that are laid off for the building up of the city of my saints, inasmuch as it shall be made known to him by the voice of the spirit \& according to the counsel of the Firm, \& by the voice of the Firm, \& this is the beginning of the stewardship which I have appointed unto him, for a blessing unto him \& his seed after him; \& inasmuch as he is faithful I will multiply a multiplicity of blessings upon him.

					And again, let my servant Newel [K. Whitney] have appointed unto him the houses \& lot where he now resides, \& the lot \& building on which the store stands, \& the lot also which is on the corner south of the store, \& also the lot on which the Ashery is situated. And all this I have appointed unto my servant Newel for his stewardship, for a blessing upon him \& his seed after him, for the benefit of the mercantile establishment of my Firm, which I have established for my Stake in the land of Kirtland; yea, verily, this is the stewardship which I have appointed unto my servant Newel, even this whole mercantile establishment, him \& his agent, \& his seed after him, \& inasmuch as he is faithful in keeping the commandments which I have given unto him, I will multiply blessings upon him, \& his seed after him, even a multiplicity of blessings.

					And again let my servant Joseph have appointed unto him the lot which is laid off for the building of my houses, which is forty rods long and twelve wide, \& also the farm upon which his father now resides, \& this is the beginning of the Stewardship which I have appointed unto him, for a blessing upon him \& upon his father; for behold, I have reserved an inheritance for his father, for his support; therefore he shall be reckoned in the house of my servant Joseph: \& I will multiply blessings upon the house of my servant Joseph inasmuch as he is faithful, even a multiplicity of blessings.

					And now a commandment I give unto you concerning Zion, that you shall no longer be bound as a United Firm, to your brethren of Zion, only on this wise: after you are organized, you shall be called, The United Firm of the Stake of Zion, the city of Kirtland, among yourselves. And your brethren, after they are organized, shall be called The United Firm of the City of Zion, \& they shall be organized in their own names, \& in their own name; \& they shall do their business in their own name, \& in their own names; \& you shall do your business in your own name, \& in your own names--- And this I have commanded to be done for your salvation, as also for their salvation, in consequence of their being driven out, and that which is to come. The covenant being broken through transgression, by coveteousness \& feigned words, therefore, you are dissolved as a United Firm with your brethren, that you are not bound only up to this hour unto them, only on this wise, as I said, By loan, as shall be agreed by this Firm in counsel as your circumstances will admit, \& the voice of the council direct.

					And again, a commandment I give unto you concerning your Stewardship which I have appointed unto you, behold, all these properties are mine, or else, your faith is vain, \& ye are found hypocrites, \& the covenants which you have made unto me are broken, \& if these properties are mine, then, ye are stewards, otherwise ye are no stewards. But, verily, I say unto you, I have appointed unto you to be Stewards over mine house, even stewards indeed, \& for this purpose have I commanded you to organize yourselves, even to print my word, the fulness of my scriptures, the revelations which I have given unto you, \& which I shall hereafter <from> time to time give unto you, for the purpose of building up my church \& kingdom on the earth, \& to prepare my people for the time of my coming which is nigh at hand. Therefore a commandment I give unto you that ye shall take the books of Mormon, \& also the copy-right, \& also the copy-right which shall be secured of the articles \& covenants, in which covenants, all my commandments, which it is my will should be printed, shall be printed, as it shall be made known unto you; \& also the copy-right to the new translation of the scriptures; \& this I say that others may not take the blessings away from you which I have conferred upon you. And ye shall prepare for yourselves a place for a Treasury, \& consecrate it unto my name, \& ye shall appoint one among you to keep the treasury \& he shall be ordained unto this blessing; \& there shall be a seal upon the Treasury, \& all these sacred things shall be delivered into the Treasury, \& no man among you shall call it his own or any part of it; for it shall belong to you all with one accord, \& I give it unto you from this very hour; \& now see to it, that ye go to \& make use of the stewardship which I have appointed unto you, exclusive of these sacred things, for the purpose of printing these sacred things, according as I have said; \& the avails of these sacred things shall be had in the Treasury, \& a seal shall be upon it, \& it shall not be used or taken out of the Treasury by any one, neither shall the seal be loosed which shall be placed upon it only by the voice of the Firm, or by commandment. And thus shall ye preserve all the avails of these sacred things in the Treasury, for sacred \& holy purposes, \& this shall be called, The Sacred Treasury of the Lord, \& a seal shall be kept upon it, that it may be holy \& consecrated unto the Lord.

					And again, there shall be another Treasury prepared \& a Treasurer appointed to keep the Treasury, \& a seal shall be placed upon it, \& all monies that you receive in your stewardships by improving upon the properties which I have appointed unto you, in houses, or in lands, or in cattle, \& in all things save it be the holy \& sacred writings, which I have reserved unto myself for holy \& sacred purposes, shall be cast into the Treasury as fast as you receive monies, by hundreds, or by fifties, or by twenties, or by tens, or by fives, or in other words, if any man among you, obtain five dollars, let him cast it into the Treasury, or if he obtain ten, or twenty, or fifty or a hundred, let him do likewise; \& let not any man among you say that it is his own; for it shall not be called his, nor any part of it, \& there shall not any part of it be ussed, or taken out of the Treasury only by the voice \& common consent of the Firm.

					And this shall be the voice \& common consent of the Firm. that any man among you, say unto the Treasurer, I have need of this to help me in my stewardship, if it be five dollars, or if it be ten dollars, or twenty, or fifty, or a hundred. The Treasurer shall give unto him the sum which he requires, to help him in his stewardship, until he be found a transgressor, \& it is manifest before the counsel of the Firm, [p. [39]] plainly that he is an unfaithful \& an unwise steward; but so long as he <is> in full fellowship \& is faithful \& wise in his stewardship, this shall be his token unto the Treasurer, that the Treasurer shall not withhold; but in case of transgression the Treasurer shall be subject unto the counsel \& voice of the Firm, \& in case the Treasurer is found an unfaithful \& an unwise steward, he shall be subject to the counsel \& voice of the Firm, \& shall be removed out of his place \& another shall be appointed in his stead. And again, verily I say unto you concerning your debts, behold, it is my will that you should pay all your debts; \& it is my will that you should humble yourselves before me, \& obtain this blessing by your diligence, \& humility \& the prayer of faith; \& inasmuch as you are diligent \& humble, \& exercise the prayer of faith, behold, I will soften the hearts of those to whom you are in debt, until I shall send means unto you for your deliverance. Therefore, write speedily unto New York, \& write according to that which shall be dictated by my Spirit, \& I will soften the hearts of those to whom you are in debt, that it shall be taken away out of their minds to bring affliction upon you. And inasmuch as ye are humble \& faithful \& call on my name, behold, I will give you the victory; I give unto you a promise, that you shall be delivered this once, out of your bondage. Inasmuch as you obtain a chance to loan money by hundreds, or by thousands, even until you shall loan enough to deliver yourselves from bondage, it is your privilege, \& pledge the properties which I have put into your hands this once by giving your names by common consent, or otherwise as it shall seem good unto you, I give unto you the privilege this once, \& behold, if you proceed to do the things which I have laid before you, according to my commandment, all these things are mine, \& ye are my stewards, \& the Master will not suffer his house \uline{to be broken up}; \uline{even so}, \uline{Amen}.

					copied from the original by O[rson] Pratt [\textit{8 blank lines}]
					
		% -----
		\subsubsection{Revelation, 28 April 1834}
		% -----
			\label{EMS-1834-JS-28-APR-1834}
			% \label{EMS-1830-JS-DAY-3LETTERMONTH-YEAR}
			
			\begin{description}
					\item[Print Sources]
				\end{description}

					\begin{tabular}{p{0.5in}p{1in}p{0.5in}p{1in}}
					JSP	 	& D4:33--35			& JSR 		& 260--261 \\
					DC	 	& --- 				& HC 		& --- \\
					HDDC 	& ---				& TCHC 		& --- \\
					\end{tabular}
					% HDDC = woodford's DC diss; JSR = Marquardt's JS Revelations; TCHC = Vogel HC
				
				\begin{description}
					\item[Online Access] \href{https://www.josephsmithpapers.org/paper-summary/revelation-28-april-1834/1}{Here}
					% \item[Analysis] \hyperref[XX]{Here}
				\end{description}
				
				\begin{description}
					\item[Background]
				\end{description}
				
					% It might also be useful to give a two sentence context of the period: e.g., "This speech was given in Nauvoo to a small congregation one month prior to the destruction of the Nauvoo Expositor. These and other tensions may have been on the prophet's mind when he gave the speech."
				
				\begin{description}
					\item[Original Source]
				\end{description}
				
					% brief description of original source material, with reference to JSP or somewhere else for more info
					% Background material: it might be helpful to have a note that says whether a given document was presented as a speech, was written in a journal, etc.
					% Also a comment here about how reliable the source might be, or what shape it is in, if there are large omissions or parts that are hard to understand, and so on.
				
				\begin{description}
					\item[Text]
				\end{description}
				
					\begin{tightright}
					Kirtland 28.--- April 1834.
					\end{tightright}
					
					Verily thus saith the Lord concerning the division and settlement of the United Firm: Let there be reserved three Thousand Dollars for the right and claim of the Firm in Kirtland for inheritances in due time, even when the Lord will; and with this claim to be had in remembrance when the Lord shall reveal it for a right of inheritance. Ye are made free from the Firm of Zion and the Firm \sout{of} in Zion is made free from the firm in Kirtland: \uline{Thus Saith} \uline{the Lord. Amen}
					
					Copied from the original by \uline{Orson Hyde}
					
		% -----
		\subsubsection{Letter to the Church, circa April 1834}
		% -----
			\label{EMS-1834-JS-CIR-APR-1834}
			% \label{EMS-1830-JS-DAY-3LETTERMONTH-YEAR}
			
			\begin{description}
					\item[Print Sources]
				\end{description}

					\begin{tabular}{p{0.5in}p{1in}p{0.5in}p{1in}}
					JSP	 	& D4:35--39		& JSR 		& --- \\
					DC	 	& --- 		& HC 		& 2:22--24 \\
					HDDC 	& ---		& TCHC 		& --- \\
					TPJS 	& 66-68 	&  &  \\
					\end{tabular}
					% HDDC = woodford's DC diss; JSR = Marquardt's JS Revelations; TCHC = Vogel HC
				
				\begin{description}
					\item[Online Access] \href{https://www.josephsmithpapers.org/paper-summary/letter-to-the-church-circa-april-1834/1}{Here}
					% \item[Analysis] \hyperref[XX]{Here}
				\end{description}
				
				\begin{description}
					\item[Background]
				\end{description}
				
					% It might also be useful to give a two sentence context of the period: e.g., "This speech was given in Nauvoo to a small congregation one month prior to the destruction of the Nauvoo Expositor. These and other tensions may have been on the prophet's mind when he gave the speech."
				
				\begin{description}
					\item[Original Source]
				\end{description}
				
					% brief description of original source material, with reference to JSP or somewhere else for more info
					% Background material: it might be helpful to have a note that says whether a given document was presented as a speech, was written in a journal, etc.
					% Also a comment here about how reliable the source might be, or what shape it is in, if there are large omissions or parts that are hard to understand, and so on.
				
				\begin{description}
					\item[Text]
				\end{description}
					
					\begin{tightcenter}
							\hypertarget{EMS-1834-JS-CIR-APR-1834-a}%
						THE 
							\marginnote{a}%
						ELDERS OF THE CHURCH IN KIRTLAND, TO TEHIR [THEIR] BRETHREN ABROAD.
					
						\textit{(Continued from our last.)}
					\end{tightcenter}
					
					\begin{tightright}
						\textit{Dear brethren in Christ, and companions in tribulation.}
					\end{tightright}
					
					IN our own country, surrounded with blessings innumerable, to which thousands of our fellow men are strangers, enjoying unspeakable benefits, and inexpressible comforts, when once our situation is compared with the ancient saints, as followers of the Lamb of God who has taken away our sins by his own blood, we are bound to rejoice and give thanks to him always. Since the organization of the church of Christ, or the church of the LATTER DAY SAINTS, which was on the 6th of April, 1830, we have had the satisfaction of witnessing the spread of the truth into various parts of our land, notwithstanding its enemies have exerted their unceasing diligence to stop its course and prevent its progress.
						\hypertarget{EMS-1834-JS-CIR-APR-1834-b}%
					Though
						\marginnote{b}%
					evil and designing men have been combined to destroy the innocent, because their own craft was in danger, and have been assisted in raising \textit{mobs} and circulating falsehoods by a miserable set of apostates, who have, for wicked and unbecoming conduct, been expelled from the body of which they were once members, yet the glorious gospel in its fulness is spreading and daily gaining converts, and our prayer to God is, that it may continue, and numbers be added of such as \textit{shall} be saved.
					
						\hypertarget{EMS-1834-JS-CIR-APR-1834-c}%
					The
						\marginnote{b}%
					Messiah's kingdom on earth is of that kind of government, that there has always been numerous apostates, for this very fact, that it admits of no sins unrepented of without excluding the individual from its fellowship. Our Lord said, Strive to enter in at the strait gate: for many, I say unto you, will seek to enter in, and shall not be able. And again, many are called, but few chosen.
						\hypertarget{EMS-1834-JS-CIR-APR-1834-d}%
					Paul
						\marginnote{d}%
					said to the elders of the church at Ephesus, after he had labored three years with them, that he knew, that some of their own number would turn away from the faith, and seek to lead away disciples after them. None, we presume, in this generation will pretend that they have the experience of Paul, in building up the church of Christ; and yet, after his departure from the church at Ephesus, many, even of the elders, turned away from the truth; and what is almost always the case, sought to lead away disciples after them. Strange as it may appear, at first thought, yet it is no less so than true, that with all the professed determination to live godly, after turning from the faith of Christ, apostates have, unless they have speedily repented, sooner or later, fallen into the snares of the wicked one and been left destitute of the Spirit of God, to manifest their wickedness in the eyes of multitudes.
						\hypertarget{EMS-1834-JS-CIR-APR-1834-e}%
					From
						\marginnote{e}%
					apostates the faithful have received the severest persecutions: Judas was rebuked, and immediately betrayed his Lord into the hands of his enemies, because \textit{satan} entered into him. There is a supreme intelligence bestowed upon such as obey the gospel with full purpose of heart, which, if sinned against, the apostate is left \textit{naked} and destitute of the Spirit of God, and they are in truth, nigh unto cursing, and their end is to be burned. When once that light which was in them is taken from them, they become as much darkened as they were previously enlightened. And then, no marvel, if all their power should be enlisted against the truth, and they, Judas like, seek the destruction of those who were their greatest benefactors!
						\hypertarget{EMS-1834-JS-CIR-APR-1834-f}%
					What
						\marginnote{f}%
					nearer friend on earth, or in heaven, had Judas, than the Savior? and his first object was to destroy him! Who, among all the saints in these last days, can consider himself as good as our Lord? Who is as perfect, who is as pure, and who as holy as he was? Are they to be found? He never transgressed or broke a commandment or law of heaven---no deceit was in his mouth, neither was guile found in his heart! and yet one that ate with him, who had often supped of the same cup, was the first to lift up his heel against him! Where is there one like him? He cannot be found on earth.
						\hypertarget{EMS-1834-JS-CIR-APR-1834-g}%
					Then
						\marginnote{g}%
					why should his followers complain, if from those whom they once called brethren, and considered in the nearest relation in the everlasting covenant, they should receive persecution? From what source emanated the principle which has ever been manifested by apostates from the true church, to persecute with double diligence, and seek with double perseverance, to destroy those whom they once professed to love, with whom they once communed, and with whom they once covenanted to strive, with every power, in righteousness, to obtain the rest of God? Perhaps, our brethren will say, The same that caused satan to seek to overthrow the kindom of God, because he himself was evil, and God's kingdom is holy.
					
						\hypertarget{EMS-1834-JS-CIR-APR-1834-h}%
					Being
						\marginnote{h}%
					limited to a short space in this number of the Star, we have advanced these few items, though in short, in stead of persuing our subject as in former numbers. The great plan of salvation is a theme which ought to occupy our strictest attention, and be regarded as one of heaven's best gifts to mankind. No consideration whatever ought to deter us from approving ourselves in the sight of God, according to his divine requirement. Men not unfrequently forget, that they are dependent upon heaven for every blessing which they are permitted to enjoy, and that for every opportunity, granted them, they are to give an account.
						\hypertarget{EMS-1834-JS-CIR-APR-1834-i}%
					You
						\marginnote{i}%
					know, brethren, that when the Master called his servants, he gave them their several benefits to improve only while he should tarry for a little season, and then he will call each to render his account; and where five tallents were bestowed, ten will be required, and he that has made no improvement will be cast out as an unprofitable servant, and the faithful are to enjoy everlasting honors.--- Therefore, we earnestly emplore the grace of our Father to rest upon you, through Jesus Christ his Son, that you may not faint in the hour of temptation, nor be overcome in the time of persecution. \textsc{To be continued.}
					
		% -----
		\subsubsection{Letter to Emma Smith, 4 June 1834}
		% -----
			\label{EMS-1834-JS-4-JUN-1834}
			% \label{EMS-1830-JS-DAY-3LETTERMONTH-YEAR}
			
			\begin{description}
					\item[Print Sources]
				\end{description}

					\begin{tabular}{p{0.5in}p{1in}p{0.5in}p{1in}}
					JSP	 	& D4:52--59		& JSR 		& --- \\
					DC	 	& --- 		& HC 		& --- \\
					HDDC 	& ---		& TCHC 		& --- \\
					\end{tabular}
					% HDDC = woodford's DC diss; JSR = Marquardt's JS Revelations; TCHC = Vogel HC
				
				\begin{description}
					\item[Online Access] \href{https://www.josephsmithpapers.org/paper-summary/letter-to-emma-smith-4-june-1834/1}{Here}
					% \item[Analysis] \hyperref[XX]{Here}
				\end{description}
				
				\begin{description}
					\item[Background]
				\end{description}
				
					% It might also be useful to give a two sentence context of the period: e.g., "This speech was given in Nauvoo to a small congregation one month prior to the destruction of the Nauvoo Expositor. These and other tensions may have been on the prophet's mind when he gave the speech."
				
				\begin{description}
					\item[Original Source]
				\end{description}
				
					% brief description of original source material, with reference to JSP or somewhere else for more info
					% Background material: it might be helpful to have a note that says whether a given document was presented as a speech, was written in a journal, etc.
					% Also a comment here about how reliable the source might be, or what shape it is in, if there are large omissions or parts that are hard to understand, and so on.
				
				\begin{description}
					\item[Text]
				\end{description}
				
						\hypertarget{EMS-1834-JS-4-JUN-1834-a}%
					...The
						\marginnote{a}%
					whole of our journey, in the midst of so large a company of social honest \sout{men} and sincere men, wandering over the plains of the Nephites, recounting occasionaly the history of the Book of Mormon, roving over the mounds of that once beloved people of the Lord, picking up their skulls \& their bones, as a proof of its divine authenticity, and gazing upon a country the fertility, the splendour and the goodness so indescribable, all serves to pass away time unnoticed, and in short were it not at every now and then our thoughts linger with inexpressible anxiety for our wives and our children our kindred according to the flesh who are entwined around our hearts;
						\hypertarget{EMS-1834-JS-4-JUN-1834-b}%
					And
						\marginnote{b}%
					also our brethren and friends; our whole journey would be as a dream, and this would be the happiest period of all our lives. We learn this journey how to travel, and we look with pleasing anticipation for the time to come, when we shall retrace our steps, and take this journey again in the enjoyment and embrace of that society we so much love, which society can only cause us to have any desire or lingering thoughts of that which is below.
						\hypertarget{EMS-1834-JS-4-JUN-1834-c}%
					We
						\marginnote{c}%
					have not as yet heard any thing from Lyman [Wight] and Hyrum [Smith] and do not expect to till we get to salt river Church, which is only fifty miles from this place. Tell Father Smith [Joseph Smith Sr.] and all the family, and brother Oliver [Cowdery] to be comforted and look forward to the day when the trials and tribulations of this life will be at an end, and we all enjoy the fruits of our labour if we hold out faithful to the end which I pray may be the happy lot of us all...
					
		% -----
		\subsubsection{Revelation, 22 June 1834 [D\&C 105]}
		% -----
			\label{EMS-1834-JS-22-JUN-1834}
			% \label{EMS-1830-JS-DAY-3LETTERMONTH-YEAR}
			
			\begin{description}
					\item[Print Sources]
				\end{description}

					\begin{tabular}{p{0.5in}p{1in}p{0.5in}p{1in}}
					JSP	 	& D4:69--77			& JSR 		& 263--266 \\
					DC	 	& Section 105 		& HC 		& 2:108--111 \\
					HDDC 	& 1372--1391		& TCHC 		& 2:121--122 \\
					\end{tabular}
					% HDDC = woodford's DC diss; JSR = Marquardt's JS Revelations; TCHC = Vogel HC
				
				\begin{description}
					\item[Online Access] \href{https://www.josephsmithpapers.org/paper-summary/revelation-22-june-1834-dc-105/1}{Here}
					% \item[Analysis] \hyperref[DC105]{Here}
				\end{description}
				
				\begin{description}
					\item[Background]
				\end{description}
				
					% It might also be useful to give a two sentence context of the period: e.g., "This speech was given in Nauvoo to a small congregation one month prior to the destruction of the Nauvoo Expositor. These and other tensions may have been on the prophet's mind when he gave the speech."
				
				\begin{description}
					\item[Original Source]
				\end{description}
				
					% brief description of original source material, with reference to JSP or somewhere else for more info
					% Background material: it might be helpful to have a note that says whether a given document was presented as a speech, was written in a journal, etc.
					% Also a comment here about how reliable the source might be, or what shape it is in, if there are large omissions or parts that are hard to understand, and so on.
				
				\begin{description}
					\item[Text]
				\end{description}
				
					Revelation given \sout{22} June 22, 1834. Clay Co. Mo.
					
					Verily I say unto you, who have assembled together that you may learn my will, concerning the redemption of mine afflicted people; behold I say unto you, were it not for the transgression of my people speaking concerning the church and not individuals, they might have been redeemed even now; but behold, they have not learned to be obedient to the things which I require at their hands, but are full of all manner of evil and do not impart of their substanc[e] as becometh saints; to the poor and afflicted among them and are not united, according to the union \sout{of} required by the law of the celestial kingdom and Zion cannot be built up unless it is by the principoles of the law of the Celestial kingdom otherwise, I cannot receive her unto my-self and my people must needs be chastened, until they learn obedience if it must needs be by the [things] which they suffer. I speak not concerning those who are appointed to lead my people who are the first elders of my church for they are not all under this condemnation, but I speak concerning the church abroad, there are many who will say where is their god, Behold, he will deliver in time of trouble, otherwise we will not go up unto Zion, and will keep our monies.
					
					Therefore in consequence of the transgression of my people, it is expedient in me that mine elders should wait for a little season for the redemption of Zion, that they themselves may be prepared and that my people may be taught more perfectly, and have experience and know more perfectly concerning their duty, and the things which I require at their hands, and this cannot be brought to pass until mine elders are endowed with power from on high, for behold I have prepared a great endowment and blessing to be poured out upon them--- in as much as they are faithful and continue in humility before me. Therefore it is expedient in me, that mine elders should wait for a little season for the redemption of Zion, For behold I do not require at their hands, to fight the battles of Zion. for as I have said in a former commandment even so I will fulfil. I will fight your battles, behold the destroyer I have already sent forth to destroy and lay waste mine enemies, and not many years hence they shall not be left to pollute mine heritage, to blaspheme my name. upon the land which I have \sout{consecranted} consecrated for the gathering together of my Saints.
					
					Behold I have commanded my servant Joseph, to say to the strength of my house; even my wariers my young men and middle aged to gather \sout{to} together for the redemption of my people and throw down the tower of mine enemies, and scatter their watchmen--- but the strength of my house has not hearkened unto \sout{unto} my words, but in as much as there are those that have hearkened unto my words, I have prepared a blessing and an endowment for them, I have heard their prayers and will accept their offering, and it is expedient in me that they should be brought thus far, for a trial of their faith.
					
					And now verily I say unto you, a commandment I give unto you, that as many as have come up hither that can stay in the region round about, Let them stay, and those who cannot stay, who have families in the east, let them tarry for a little season, in as much as my servant Joseph shall appoint unto them for I will counsel him concerning this matter, And all things whatsoever \sout{is} he shall appoint unto them shall be fulfilled and let all my people who dwell in the region round about, be very faithful and prayerful and humb[l]e before me, and reveal not the things which I have revealed unto them, Talk not of Judgment boast not of faith nor of mighty works, but carefully gather together in one region as can be consistently with the feeling of the people. And Behold I will give unto you favor and grace in their eyes, that you may rest in peace and safety, whilst you are saying unto the people execute judgment Justice for us according to law, and redress us of our wrongs.
					
					Now behold, I say unto you my friends, in this way you may find favor in the eyes of the people, until the armies of Israel become very great, and I will soften the hearts of the people as I did \sout{I} the heart of Pharioh from time to time, until my servant Joseph and mine elders whom he shall appoint shall have time to gather up the strength of my house; and to have sent wise men to fulfil that which I have commanded concerning the purchasing of all the lands in Jackson County that can be purchased and in the adjoining Counties round about, for it is my will, that these lands should be purchased that my saints should possess them according to the law of consecration, which I have given. And after these lands are purchased I will hold the armies of Israel guiltless in taking possessions of their own lands and of throw[i]ng down the tower of mine enemies that may be upon them, and scattering their watchmen, and avengeing me of mine enemies, unto the third and fourth generation of them that hate me. But firstly let my army become very great, and let it be sanctified before me, that it may become fair as the Sun and clear as the moon, and that her banners may be terrable unto all nations, that the kingdom of this world may be constrainned to acknowledge that the kingdom of Zion, is in very deed the kingdom of our God, and his Christ,
					
					Therefore let us become subject unto her laws, Verily I say unto you it is expedient in me, that the first elders of my church should receive their endowment from on high in mine house which I have commanded to be built unto my name in the land of Kirtland, and let those commandments which I have given concerning Zion, and her law be executed and fulfilled after her redemption. There has been a day of calling. but the time has come for a day of choosing, and let those be chosen that are worthy and it shall be manifest unto my servant Joseph by the voice of the Spirit, those who are chosen, and they shall be sanctified, and in as much as they, follow the counsels which they receive they shall have power after many days to accomplish all things partaining to Zion.
					
					<And> again I say unto you, sue for peace not only the people that have smitten you, but also to all people. and lift up an ensign of peace, and make a proclamation \sout{of} for peace unto the ends of the earth and make proposals \sout{of} <for> peace unto those who have smitten you, according to the voice of the spirit which is in you, and all thing shall work together for your good, and be faithful and behold and low I am with you even unto the end. even so Amen------
					
		% -----
		\subsubsection{Minutes, 23 June 1834}
		% -----
			\label{EMS-1834-JS-23-JUN-1834}
			% \label{EMS-1830-JS-DAY-3LETTERMONTH-YEAR}
			
			\begin{description}
					\item[Print Sources]
				\end{description}

					\begin{tabular}{p{0.5in}p{1in}p{0.5in}p{1in}}
					JSP	 	& D4:80--84		& JSR 		& --- \\
					DC	 	& --- 		& HC 		& 2:112--113 \\
					HDDC 	& ---		& TCHC 		& 2:125--127 \\
					\end{tabular}
					% HDDC = woodford's DC diss; JSR = Marquardt's JS Revelations; TCHC = Vogel HC
				
				\begin{description}
					\item[Online Access] \href{https://www.josephsmithpapers.org/paper-summary/minutes-23-june-1834/1}{Here}
					% \item[Analysis] \hyperref[XX]{Here}
				\end{description}
				
				\begin{description}
					\item[Background]
				\end{description}
				
					% It might also be useful to give a two sentence context of the period: e.g., "This speech was given in Nauvoo to a small congregation one month prior to the destruction of the Nauvoo Expositor. These and other tensions may have been on the prophet's mind when he gave the speech."
				
				\begin{description}
					\item[Original Source]
				\end{description}
				
					% brief description of original source material, with reference to JSP or somewhere else for more info
					% Background material: it might be helpful to have a note that says whether a given document was presented as a speech, was written in a journal, etc.
					% Also a comment here about how reliable the source might be, or what shape it is in, if there are large omissions or parts that are hard to understand, and so on.
				
				\begin{description}
					\item[Text]
				\end{description}
				
					...A council of High Priests met according to revelation in order to choose some of the first Elders to receive their endowments.--- being appointed by the voice of the Spirit through Br. Joseph Smith jr. President of the Church of Christ...%
						\footnote{%
							\label{EMS-1834-JS-23-JUN-1834-NOTE-explanation}%
							Talk about what follows in this document, a bunch of times it saying that people are going to receive their endowment.
							}
							
		% -----
		\subsubsection{Minutes and Discourse, circa 7 July 1834}
		% -----
			\label{EMS-1834-JS-7-JUL-1834}
			% \label{EMS-1830-JS-DAY-3LETTERMONTH-YEAR}
			
			\begin{description}
					\item[Print Sources]
				\end{description}

					\begin{tabular}{p{0.5in}p{1in}p{0.5in}p{1in}}
					JSP	 	& D4:90--96		& JSR 		& --- \\
					DC	 	& --- 		& HC 		& 2:124--126 \\
					HDDC 	& ---		& TCHC 		& 2:136--138 \\
					\end{tabular}
					% HDDC = woodford's DC diss; JSR = Marquardt's JS Revelations; TCHC = Vogel HC
				
				\begin{description}
					\item[Online Access] \href{https://www.josephsmithpapers.org/paper-summary/minutes-and-discourse-circa-7-july-1834/1}{Here}
					% \item[Analysis] \hyperref[XX]{Here}
				\end{description}
				
				\begin{description}
					\item[Background]
				\end{description}
				
					% It might also be useful to give a two sentence context of the period: e.g., "This speech was given in Nauvoo to a small congregation one month prior to the destruction of the Nauvoo Expositor. These and other tensions may have been on the prophet's mind when he gave the speech."
				
				\begin{description}
					\item[Original Source]
				\end{description}
				
					% brief description of original source material, with reference to JSP or somewhere else for more info
					% Background material: it might be helpful to have a note that says whether a given document was presented as a speech, was written in a journal, etc.
					% Also a comment here about how reliable the source might be, or what shape it is in, if there are large omissions or parts that are hard to understand, and so on.
				
				\begin{description}
					\item[Text]
				\end{description}
				
						\hypertarget{EMS-1834-JS-7-JUL-1834-a}%
					...After
						\marginnote{a}%
					which br. Joseph Smith jun proceeded to give the Council such instruction (relative to their high calling) as would enable them to proceed and minister in their office--- agreeable to the pattern given heretofore;--- also read to them the Revelation speaking on the subject--- He also informed them if he should now be taken away that he had accomplished the great work which the Lord had laid before him, and that which he had desired of the Lord, and that he now had done his duty in organizing the High Council, through which Council the will of the Lord might be known on all importent occasions in the building up of Zion, and establishing truth in the earth.
					
						\hypertarget{EMS-1834-JS-7-JUL-1834-b}%
					A
						\marginnote{b}%
					vote was then taken whether the brethren that were appointed on the 3\supers{rd} of July should stand according to appoinment. The vote was clear in their behalf.
					
					Br. Joseph Smith jr. then proceeded and ordained the three Presidents. David Whitmer as President%
						\footnote{%
							\label{EMS-1834-JS-7-JUL-1834-NOTE-david}%
							See historical introduction, note on David Whitmer being chosen as JS's successor?
							}
					and William W. Phelps \& John Whitmer assistants and their twelve Counsellors...
					
		% -----
		\subsubsection{Minutes, 11 August 1834}
		% -----
			\label{EMS-1834-JS-11-AUG-1834}
			% \label{EMS-1830-JS-DAY-3LETTERMONTH-YEAR}
			
			\begin{description}
					\item[Print Sources]
				\end{description}

					\begin{tabular}{p{0.5in}p{1in}p{0.5in}p{1in}}
					JSP	 	& D4:97--101		& JSR 		& --- \\
					DC	 	& --- 		& HC 		& 2:142--144 \\
					HDDC 	& ---		& TCHC 		& 2:155--157 \\
					\end{tabular}
					% HDDC = woodford's DC diss; JSR = Marquardt's JS Revelations; TCHC = Vogel HC
				
				\begin{description}
					\item[Online Access] \href{https://www.josephsmithpapers.org/paper-summary/minutes-11-august-1834/1}{Here}
					% \item[Analysis] \hyperref[XX]{Here}
				\end{description}
				
				\begin{description}
					\item[Background]
				\end{description}
				
					% It might also be useful to give a two sentence context of the period: e.g., "This speech was given in Nauvoo to a small congregation one month prior to the destruction of the Nauvoo Expositor. These and other tensions may have been on the prophet's mind when he gave the speech."
				
				\begin{description}
					\item[Original Source]
				\end{description}
				
					% brief description of original source material, with reference to JSP or somewhere else for more info
					% Background material: it might be helpful to have a note that says whether a given document was presented as a speech, was written in a journal, etc.
					% Also a comment here about how reliable the source might be, or what shape it is in, if there are large omissions or parts that are hard to understand, and so on.
				
				\begin{description}
					\item[Text]
				\end{description}
				
						\hypertarget{EMS-1834-JS-11-AUG-1834-a}%
					Kirtland
						\marginnote{a}%
					August 11\supersu{th} 1834

					This day a number of high priests and elders of the church of the Latter-Day Saints assembled in the new school-house in this place, for the purpose of investigating a matter of difficulty growing out of certain reports or statements made by brother Sylvester Smith; one of the High counsellors of this Church, accusing brother Joseph Smith Jun\supers{r.} with criminal conduct during his journey to and from Missouri this Spring \& Summer.
						\hypertarget{EMS-1834-JS-11-AUG-1834-b}%
					After
						\marginnote{b}%
					coming to order, brother Joseph commenced and spake to a considerable length upon the circumstances of their journey to and from Missouri, and very minutely laid open the causes out of which those jealousies of brother Sylvesters and others had grown. He made a satisfactory statement concerning his rebukes and chastisements upon Sylvester \& others, and also concerning the distribution of monies and other properties, calling on brethren present who accompanied him to attest to the same. All of which, was satisfactory to the brethren present as appeared by their own remarks afterward.
						\hypertarget{EMS-1834-JS-11-AUG-1834-c}%
					After
						\marginnote{c}%
					brother Joseph had closed his lengthy remarks, brother Sylvester made some observations relative to the subject of their difficulties, and begun to make a partial confession for his previous conduct, asking forgiveness for accusing brother Joseph publicly on the Saturday previous, of proph[e]sying lies in the name of the Lord, and for abusing (as he had said) his (Sylvester's) character, before the brethern while journeying to the West.
						\hypertarget{EMS-1834-JS-11-AUG-1834-d}%
					Brother
						\marginnote{d}%
					Sidney Rigdon made some remarks by way of reproff upon the conduct of brother Sylvester, J. P. Green [John P. Greene] [spoke], and others, followed by the Clerk, after which by motion of brother S. Rigdon the assembly arranged itself into a council, brother N[ewel] K. Whitney presiding, and proceeded to discuss how this difficulty should be disposed of. Brother John Smith thought that for brother Sylvester, to make a public confession in the Star would be the way to heal the wound. Brother [Reynolds] Cahoon followed with nearly the same remarks Brother Isaac Hill, thought it ought to be quashed here, and go no farther. followed with the same from brother I[saac] Bishop.
						\hypertarget{EMS-1834-JS-11-AUG-1834-e}%
					Brother
						\marginnote{e}%
					Samuel H. Smith said that it was his opinion, that brother Sylvester ought to make a more public confession, and send by letter to those who are in the same transgression, with himself and inform them of this decision, and then if necessary make it public in the Star. Brother Orson Hyde thought that the confession ought to be as liberal as the accusation, or that it ought to be written and published Brother J. P. Green, said that if brother Sylvester would view this thing in its proper light, he would be willing to make a public confession and send forth; and he advised him to do this for the salvation of the Churches abroad.
						\hypertarget{EMS-1834-JS-11-AUG-1834-f}%
					Brother
						\marginnote{f}%
					Isaac Story, said that it was his opinion the plaster ought to be as large as the wound, that a proper statement ought to be published abroad. The clerk then proposed that this council send a certificate or resolution, informing the churches abroad, that the conduct of brother Joseph, has been investigated, and that he has acted in a proper manner and in every respect has conducted himself to the satisfaction of the church in Kirtland, and also let brother Sylvester make a proper confession following the same minutes. Brother A[sa] Lyman, P[eter] Shirts, T[ruman] Wait, R[oswell] Evans. A[lpheus] Cutler and T[homas] Burdick, made remarks to the same effect. Brother S. Rigdon, made a few remarks upon the attitude in which, he, Sylvester now stands before the world, in endeavoring to preach the gospel.

						\hypertarget{EMS-1834-JS-11-AUG-1834-g}%
					Brother
						\marginnote{g}%
					O. Hyde moved for a decision, relative to the first question, (viz.) What is to be done to arrest the evil? The moderator then proceeded, after a few remarks, to give a decision according, to a motion previously made, (viz.) that an article be published in the Evening \& the morning Star, by the direction of the council, that the church in Kirtland has investigated the conduct of brother Joseph Smith \supersu{Junr}\supers{.} while journeying to the West and returning, and that we find that he has acted in every respect in \sout{the} an honorable and proper Manner, with all monies and other properties entrusted to his charge. After which a vote was taken and carried. A. motion was then made by brother O. Hyde and seconded by brother S. Rigdon, that a committee of three, be appointed to write the article for the Star agreeably to the decision.%
						\footnote{%
							\label{EMS-1834-JS-11-AUG-1834-NOTE-decision}%
							Note to minutes of 23 August 1834 (not in here), where Sylvester reverses his decision. (There is more material to add in the footnotes here if I want, from other sources, including 28--29 August 1834.)
							}

						\hypertarget{EMS-1834-JS-11-AUG-1834-h}%
					Brethren,
						\marginnote{h}%
					The Clark [clerk], Thomas Burdick and Orson Hyde were nominated and appointed by unanimous vote.

					Brother Sylvester, then \sout{made} said that he was willing to publish a confession in the Star. A motion was then made by brother J. P. Green, and seconded by the moderator, that the above named committe, be appointed to write letters in the name of the council to bro--- Z[erubbabel] Snow, C. Snelling [Cyrus Smalling] \& J. P. [John D.] Parker.

					Prayer by brother S. Rigdon.

					\begin{tightright}
						Oliver Cowdery)

						Clerk of Council)
					\end{tightright}
					
		% -----
		\subsubsection{Letter to Lyman Wight and Others, 16 August 1834}
		% -----
			\label{EMS-1834-JS-16-AUG-1834}
			% \label{EMS-1830-JS-DAY-3LETTERMONTH-YEAR}
			
			\begin{description}
					\item[Print Sources]
				\end{description}

					\begin{tabular}{p{0.5in}p{1in}p{0.5in}p{1in}}
					JSP	 	& D4:102--108		& JSR 		& --- \\
					DC	 	& --- 		& HC 		& 2:144--146 \\
					HDDC 	& ---		& TCHC 		& 2:157--160 \\
					\end{tabular}
					% HDDC = woodford's DC diss; JSR = Marquardt's JS Revelations; TCHC = Vogel HC
				
				\begin{description}
					\item[Online Access] \href{https://www.josephsmithpapers.org/paper-summary/letter-to-lyman-wight-and-others-16-august-1834/1}{Here}
					% \item[Analysis] \hyperref[XX]{Here}
				\end{description}
				
				\begin{description}
					\item[Background]
				\end{description}
				
					% It might also be useful to give a two sentence context of the period: e.g., "This speech was given in Nauvoo to a small congregation one month prior to the destruction of the Nauvoo Expositor. These and other tensions may have been on the prophet's mind when he gave the speech."
				
				\begin{description}
					\item[Original Source]
				\end{description}
				
					% brief description of original source material, with reference to JSP or somewhere else for more info
					% Background material: it might be helpful to have a note that says whether a given document was presented as a speech, was written in a journal, etc.
					% Also a comment here about how reliable the source might be, or what shape it is in, if there are large omissions or parts that are hard to understand, and so on.
				
				\begin{description}
					\item[Text]
				\end{description}
				
						\hypertarget{EMS-1834-JS-16-AUG-1834-a}%
					Kirtland
						\marginnote{a}
					August 16th 1834
					
					Copy from Joseph S to the brethr[e]n in Zion
					
					Dear Brethren Lyman Wight Edward Partri[d]ge John Carrill [Corrill] Isaac Morley and others the high council of the Church of the Latter Day Saints------
					
						\hypertarget{EMS-1834-JS-16-AUG-1834-b}%
					After
						\marginnote{b}%
					so long a time I dictate a few lines to you to let you know that I am in Kirtland and that I found all well on my arival, as pertaining to health \&c but \sout{found} our common advisary had taken the advantage of our brothe[r] Sylvest[er] Smith and others who gave a false colloring to allmost every transaction from the time that we left Kirtland untill we returned, and thereby Stirred up a great difficulty in the Church against me accordingly I was met in the face and eyes as soon as I had got home with a catalogue that was as black as the author himself and the cry was Tyrant! Pope!! King!!! Usurper!!!! Abuser of men!!!!! Ange!!!!!! False proph[e]t\sout{!!!!!} Prophecying Lies in the name of the Lord\sout{!!!!!!!} and taking Consecrated monies!!!!!!!! and every other lie to fill up and complete the cattelogue that was necissary to perfect the Church to be meet for the devourer the shaft of the \sout{devouring} <distroying> Angel!
						\hypertarget{EMS-1834-JS-16-AUG-1834-b}%
					and
						\marginnote{b}%
					in consequence of having to combat all these I have not been able to regulate my mind so as to write to give you council and the information that you needed, but that God who rules on high and thunders Judgments upon Israel when they transgress has given me power from the time that I was born (into this Kingdom) to stand and I have succeeded in putting all gainsayers and enemies to flight unto the present time and not withstanding the advisary Laid a plan which was more subtle than all others, I now swim in good \uline{clean} \sout{pure} water with my \uline{head} out! as you will see by the next star
					
						\hypertarget{EMS-1834-JS-16-AUG-1834-c}%
					I
						\marginnote{c}%
					shall now procede to give you such council as the spirit of the Lord may dictate you will reccollect that your business must be done by your high council: you will recollect that the first elders are to receive their endowment in Kirtland before the redemption of Zion you will reccollect that your high council will have power to say who of the first Elders among the Children of Zion are accounted worthy; and you will also reccollect that you have my testamony in be half of certain ones previously to my departure you will reccollect that the sooner that these ambassadors of the most high are dispatched to bear testamony to lift up a warning voice and to proclaim the everlasting gospel and to use every convincing proof and facculty with this generation while on their Journey. <to Kirtland*>
						\hypertarget{EMS-1834-JS-16-AUG-1834-d}%
					<*The
						\marginnote{d}%
					better it shall be for them and for Zion inasmuch as the indignation of the people sleepeth for a while our time should be employed to the best advantage altho it is not the will of God that any one of these ambassador should hold their peace afte[r] they have startd upon their Journey.> They should awaken <the> sympathy of the people. I would reccommend to brother [William W.] Phelps (If he is yet there) to write a petition such as will be approved of by the high council and let there be every signer obtained that can be in the State of Missouri and while they are on their Journey to this country that paradventure we may learn whithe [whether] we have friends or not in these United States,
					
						\hypertarget{EMS-1834-JS-16-AUG-1834-e}%
					This
						\marginnote{e}%
					petition to be sent to the Govenor of Missouri to solicit him to call on the President of the United States for a guard to protect our brethren in Jackson County upon their own Lands from the insults and abuses of the \uline{Mob}
					
					And I would reccomend to brother Wight to enter complaints to the Govonor as of ten as he receves any insults or injury, and in case that they procede to endeaver to take life or tear down homes, and if the citizens of Clay co, do not befriend us to gather up the little army and be set over Immediately into Jackson County and trust in God and do the \sout{worst} <best> he can in \sout{defending} maintaining the ground, but in case the excitement continues to be allayed and peace prevails use every effort to prevail on the churches to gather to those regions and situate themselves to be in readiness to move into Jackson Co. in two years from the Eleventh of September next which is the appointed time for the redemption of Zion,
						\hypertarget{EMS-1834-JS-16-AUG-1834-f}%
					\phantom{ }\uline{If}
						\marginnote{f}%
					Verely \sout{If} I say unto you \uline{If} the church with one united effort perform their duties If they do this the work shall be complete If they do not this in all humility making preperation from this time forth like Joseph in Egypt laying up store against the time of famine every man having his tent, his horses, his charrots [chariots] his armory his cattle his family and his whole substance in readiness against the time <when> it shall be said \uline{To your tents O Isreal}!! and let not this be noised abroad let every heart beat in silence and every mouth be shut
					
						\hypertarget{EMS-1834-JS-16-AUG-1834-g}%
					Now
						\marginnote{g}%
					my beloved brethren you will learn by this we have a great work to do, and but little time to \sout{do} it in and if we dont exert ourselves to the utmost in gathering up the strength of the Lords house that this thing may be accomplished behold their remaineth a scorge* <*for the Church even that they shall be driven from City to City and but few shall remain to receive an inheritence if these things are not kept there remaineth a scorge> Also, Therefore be wise this once O ye children of Zion! and give heed to my council saith the Lord!...
					
					...so long as unrighteous\sout{ness} acts are suffered in the church it cannot [be] sanctified neither Zion be redeemed...
					
		% -----
		\subsubsection{Minutes, 8 September 1834}
		% -----
			\label{EMS-1834-JS-8-SEP-1834}
			% \label{EMS-1830-JS-DAY-3LETTERMONTH-YEAR}
			
			\begin{description}
					\item[Print Sources]
				\end{description}

					\begin{tabular}{p{0.5in}p{1in}p{0.5in}p{1in}}
					JSP	 	& D4:164--168		& JSR 		& --- \\
					DC	 	& --- 		& HC 		& 2:162--163 \\
					HDDC 	& ---		& TCHC 		& 2:176--178 \\
					\end{tabular}
					% HDDC = woodford's DC diss; JSR = Marquardt's JS Revelations; TCHC = Vogel HC
				
				\begin{description}
					\item[Online Access] \href{https://www.josephsmithpapers.org/paper-summary/minutes-8-september-1834/1}{Here}
					% \item[Analysis] \hyperref[XX]{Here}
				\end{description}
				
				\begin{description}
					\item[Background]
				\end{description}
				
					% It might also be useful to give a two sentence context of the period: e.g., "This speech was given in Nauvoo to a small congregation one month prior to the destruction of the Nauvoo Expositor. These and other tensions may have been on the prophet's mind when he gave the speech."
				
				\begin{description}
					\item[Original Source]
				\end{description}
				
					% brief description of original source material, with reference to JSP or somewhere else for more info
					% Background material: it might be helpful to have a note that says whether a given document was presented as a speech, was written in a journal, etc.
					% Also a comment here about how reliable the source might be, or what shape it is in, if there are large omissions or parts that are hard to understand, and so on.
				
				\begin{description}
					\item[Text]
				\end{description}
				
						\hypertarget{EMS-1834-JS-8-SEP-1834-a}%
					New-Portage,
						\marginnote{a}
					Ohio Sept 8\supersu{th}\supers{.} 1834

					Minutes of a conference of the Elders of the Church of the Latter-Day Saints, held at New-Portage Ohio. September 8\supersu{th} 1834, After prayer by brother Joseph Smith J\supersu{unr}\supers{.} he, brother Joseph and Oliver Cowdery, united in anointing with oil and laying hands upon a sick sister, who said she was healed, and requested us to pray that her faith fail not, saying, if she did not doubt, she should not be afflicted any more.

						\hypertarget{EMS-1834-JS-8-SEP-1834-b}%
					Brother
						\marginnote{b}%
					Joseph then made a few introductory remarks upon the subject of false spirits and other items.

					Brother A[mbrose] Palmer made a few observations, and proceeded to present a case which had previously occasioned some difficulty in the church. It was that a certain brother Carpenter had been tried for a fault before the church, when the Church gave him a certain time to reflect whether he would acknowledge or not. Brother Gordon at the time, spake in tongues and declared that brother Carpenter should not have any lenity. He, brother Palmer, wished instruction on this point, whether they had proceeded right or not, as brother Carpenter was dissatisfied \&c.

						\hypertarget{EMS-1834-JS-8-SEP-1834-c}%
					Brother
						\marginnote{c}%
					Joseph then proceeded to give an explanation of the gift of tongues: that it was particularly instituted for the preaching of the Gospel to other nations and languages, but it was not given for the government of the Church. He further said, if brother Gordon introduced the Gift of tongues, as a testimony against brother Carpenter, that it was contrary to the rules and regulations of the Church, because, in all our decisions we must judge from actual testimony. Brother Gordon said the testimony was had and the decision given before the gift of tongues was manifested. Brother Joseph advised that <we> speak in our own language in all such matters, and then the adversary cannot lead our minds astray.
						\hypertarget{EMS-1834-JS-8-SEP-1834-d}%
					Brother
						\marginnote{d}%
					Palmer then gave a relation of a certain difficulty which took place in a conference. He, brother Palmer presided, when several of the brethren spake out of order, and brother J[oseph] B. Bosworth refused to submit to to order according to his (brother Palmer's) request. He now wished instruction on this point, whether he, or some one else should preside over this branch of the Church, and also whether such conduct could be approbated in conferences. Brother Gordon then made some remarks on the subject which was at the time before the council.

						\hypertarget{EMS-1834-JS-8-SEP-1834-e}%
					Brother
						\marginnote{e}%
					Joseph said, relative to the first question, that brother Gordon'\supers{s} tongues in the end, did operate as testimony, as, by his remarks in tongues, the former decision was set aside, and his taken. That it was his decision that brother Gordon'\supers{s} manifestation was incorrect, and from a suspicious heart. He approved the first decision, but discarded the second...

		% -----
		\subsubsection{Letter to Oliver Cowdery, 24 September 1834}
		% -----
			\label{EMS-1834-JS-24-SEP-1834}
			% \label{EMS-1830-JS-DAY-3LETTERMONTH-YEAR}
			
			\begin{description}
					\item[Print Sources]
				\end{description}

					\begin{tabular}{p{0.5in}p{1in}p{0.5in}p{1in}}
					JSP	 	& D4:168--171		& JSR 		& --- \\
					DC	 	& --- 		& HC 		& --- \\
					HDDC 	& ---		& TCHC 		& --- \\
					\end{tabular}
					% HDDC = woodford's DC diss; JSR = Marquardt's JS Revelations; TCHC = Vogel HC
				
				\begin{description}
					\item[Online Access] \href{https://www.josephsmithpapers.org/paper-summary/letter-to-oliver-cowdery-24-september-1834/1}{Here}
					% \item[Analysis] \hyperref[XX]{Here}
				\end{description}
				
				\begin{description}
					\item[Background]
				\end{description}
				
					% It might also be useful to give a two sentence context of the period: e.g., "This speech was given in Nauvoo to a small congregation one month prior to the destruction of the Nauvoo Expositor. These and other tensions may have been on the prophet's mind when he gave the speech."
				
				\begin{description}
					\item[Original Source]
				\end{description}
				
					% brief description of original source material, with reference to JSP or somewhere else for more info
					% Background material: it might be helpful to have a note that says whether a given document was presented as a speech, was written in a journal, etc.
					% Also a comment here about how reliable the source might be, or what shape it is in, if there are large omissions or parts that are hard to understand, and so on.
				
				\begin{description}
					\item[Text]
				\end{description}
				
						\hypertarget{EMS-1834-JS-24-SEP-1834-a}%
					\phantom{ }\textit{Kirtland,}
						\marginnote{a}
					\phantom{ }\textit{Ohio, September} 24, 1834.

					\textsc{Dear brother},---

					I have, of late, been perusing Mr. A[lexander] Campbell's ``Millennial Harbinger.'' I never have rejoiced to see men of corrupt hearts step forward and assume the authority and pretend to teach the ways of God---this is, and always has been a matter of grief; therefore I cannot but be thankful, that I have been instrumental in the providence of our heavenly Father in drawing forth, before the eyes of the world, the \textit{spirits} by which certain ones, who profess to be ``Reformers, and Restorers of ancient principles,'' are actuated!
						\hypertarget{EMS-1834-JS-24-SEP-1834-b}%
					I
						\marginnote{b}%
					have always had the satisfaction of seeing the truth triumph over error, and darkness give way before light, when such men were \textit{provoked} to expose the corruption of their own hearts, by crying delusion, deception, and false prophets, accusing the innocent, and condemning the guiltless, and exalting themselves to the stations of gods, to lead blind-fold, men to perdition!

						\hypertarget{EMS-1834-JS-24-SEP-1834-c}%
					I
						\marginnote{c}%
					have never been blessed, (if it may be called such,) with a personal acquaintance with Mr. Campbell, neither a personal interview; but the \textsc{great man}, not unfrequently \textit{condescends} to notice an individual of as obscure birth as myself, if I am at liberty to interpret the language of \textit{his} ``Harbinger,'' where he says, ``\textit{Joe} Smith! \textit{Joe} Smith! imposture! imposture!'' I have noticed a strange thing! I will inform you of my meaning, though I presume you have seen the same ere this.
						\hypertarget{EMS-1834-JS-24-SEP-1834-d}%
					Mr.
						\marginnote{d}%
					Campbell was very lavish of his expositions of the falsity and incorrectness of the book of Mormon, some time since, but of late, since the publication of the Evening and the Morning Star, has said little or nothing, except some of his back-handed \textit{cants}. He did, to be sure, about the time the church of Christ was established in Ohio, come out with a lengthy article, in which he undertook to prove that it was incorrect and contrary to the former revelations of the Lord.
						\hypertarget{EMS-1834-JS-24-SEP-1834-e}%
					Perhaps,
						\marginnote{e}%
					he is of opinion that he so completely overthrew the foundation on which it was based, that all that is now wanting to effect an utter downfall of those who have embraced its principles is, to continue to \textit{bark} and \textit{howl}, and cry, \textit{Joe} Smith! false prophet! and ridicule every man who may be disposed to examine the evidences which God has given to the world of its truth!

						\hypertarget{EMS-1834-JS-24-SEP-1834-f}%
					I
						\marginnote{f}%
					have never written Mr. Campbell, nor received a communication from him but a public notice in his paper:--- If you will give this short note a place in the Star you will do me a kindness, as I take this course to inform the gentleman, that while he is breathing out scurrility he is effectually showing the honest, the motives and principles by which he is governed, and often causes men to investigate and embrace the book of Mormon, who might otherwise never have perused it. I am satisfied, therefore he should continue his scurrility; indeed, I am more than gratified, because his cry of \textit{Joe} Smith! \textit{Joe} Smith! \textit{false} prophet! \textit{false} prophet! must manifest to \textit{all} men the spirit he is of, and serves to open the eyes of the people.

						\hypertarget{EMS-1834-JS-24-SEP-1834-g}%
					I
						\marginnote{g}%
					wish to inform him further, that as he has, for a length of time, smitten me upon one cheek, and I have offered no resistance, I have turned the other also, to obey the commandment of our Savior; and am content to sit awhile longer in silence and see the great work of God roll on, amid the opposition of this world in the face of every scandal and falsehood which may be invented and put in circulation.

					\begin{tightright}
						I am your brother in the testamony of the book of Mormon,
					\end{tightright}
					
					\begin{tightcenter}
						and shall ever remain.
					\end{tightcenter}

					\begin{tightright}
						JOSEPH SMITH jr.
					\end{tightright}

					\textsc{To Oliver Cowdery}.

		% -----
		\subsubsection{Revelation, 25 November 1834 [D\&C 106]}
		% -----
			\label{EMS-1834-JS-25-NOV-1834}
			% \label{EMS-1830-JS-DAY-3LETTERMONTH-YEAR}
			
			\begin{description}
					\item[Print Sources]
				\end{description}

					\begin{tabular}{p{0.5in}p{1in}p{0.5in}p{1in}}
					JSP	 	& D4:180--182		& JSR 		& 266 \\
					DC	 	& Section 106		& HC 		& 2:170--171 \\
					HDDC 	& 1392--1397		& TCHC 		& 2:185 \\
					\end{tabular}
					% HDDC = woodford's DC diss; JSR = Marquardt's JS Revelations; TCHC = Vogel HC
				
				\begin{description}
					\item[Online Access] \href{https://www.josephsmithpapers.org/paper-summary/revelation-25-november-1834-dc-106/1}{Here}
					% \item[Analysis] \hyperref[DC106]{Here}
				\end{description}
				
				\begin{description}
					\item[Background]
				\end{description}
				
					% It might also be useful to give a two sentence context of the period: e.g., "This speech was given in Nauvoo to a small congregation one month prior to the destruction of the Nauvoo Expositor. These and other tensions may have been on the prophet's mind when he gave the speech."
				
				\begin{description}
					\item[Original Source]
				\end{description}
				
					% brief description of original source material, with reference to JSP or somewhere else for more info
					% Background material: it might be helpful to have a note that says whether a given document was presented as a speech, was written in a journal, etc.
					% Also a comment here about how reliable the source might be, or what shape it is in, if there are large omissions or parts that are hard to understand, and so on.
				
				\begin{description}
					\item[Text]
				\end{description}
				
					Kirtland, November 25; 1834
					
					It is my will that my servant Warren [Cowdery] should \sout{receive} <be ap>pointed and ordained a presiding High Priest over my Church in the land of Freedom and the regions round about; and should preach my everlasting gospel, and lift up his voice and warn the people, not only in his own place, but in the adjoining countries, and devote his whole time in this high and holy calling which I now give unto him; seeking dilligently the kingdom of heaven and its righteousness, and all things necessary shall be added thereunto; for the laborer is worthy of his hire.
					
					And again, verily I say unto you, the coming of the Lord draweth nigh, and it overtaketh the world as a thief in the night: therefore, gird up your loins, that ye may be the children of the light, \sout{and} and that day shall not overtake you as a thief.
					
					And again, verily I say unto you, there was joy in heaven when my servant Warren bowed to my scepter and separated himself from the crafts of men. Therefore, blessed is my servant Warren, for I will have mercy on him, and notwithstanding the vanity of his heart, I will lift him up, \sout{and} inasmuch as he will humble himself before me; <and> I will give <unto> him grace and assurance wherewith he may stand; and if he continues to be a faithful witness, and a light unto the Church, I have prepared a crown for him in the mansion of my Father: even so. Amen.
					
		% -----
		\subsubsection{Covenant, 29 November 1834}
		% -----
			\label{EMS-1834-JS-29-NOV-1834}
			% \label{EMS-1830-JS-DAY-3LETTERMONTH-YEAR}
			
			\begin{description}
					\item[Print Sources]
				\end{description}

					\begin{tabular}{p{0.5in}p{1in}p{0.5in}p{1in}}
					JSP	 	& D4:188--191; J1:46--47		& JSR 		& --- \\
					DC	 	& --- 		& HC 		& 2:174--175 \\
					HDDC 	& ---		& TCHC 		& 2:189 \\
					TPJS 	& 70			& 	 		&  \\
					\end{tabular}
					% HDDC = woodford's DC diss; JSR = Marquardt's JS Revelations; TCHC = Vogel HC
				
				\begin{description}
					\item[Online Access] \href{https://www.josephsmithpapers.org/paper-summary/covenant-29-november-1834/1}{Here}
					% \item[Analysis] \hyperref[XX]{Here}
				\end{description}
				
				\begin{description}
					\item[Background]
				\end{description}
				
					% It might also be useful to give a two sentence context of the period: e.g., "This speech was given in Nauvoo to a small congregation one month prior to the destruction of the Nauvoo Expositor. These and other tensions may have been on the prophet's mind when he gave the speech."
				
				\begin{description}
					\item[Original Source]
				\end{description}
				
					% brief description of original source material, with reference to JSP or somewhere else for more info
					% Background material: it might be helpful to have a note that says whether a given document was presented as a speech, was written in a journal, etc.
					% Also a comment here about how reliable the source might be, or what shape it is in, if there are large omissions or parts that are hard to understand, and so on.
					
					% The second part after the signatures is not in Documents (though I should check again); it's only in Journals.
				
				\begin{description}
					\item[Text]
				\end{description}
				
						\hypertarget{EMS-1834-JS-29-NOV-1834-a}%
					After conversing and rejoicing before
						\marginnote{a}%
					the Lord on this occasion we agreed to enter into the following covenant with the Lord, viz:---
					
					\begin{quote}
					That if the Lord will
					
					prosper us in our business, and open the way before <us> that we may obtain means to pay our debts, that we be not troubled nor brought into disrepute before the world nor his people, that after that of all that he shall give us we will give a tenth, to be bestowed upon the poor in his church, or as he shall command, and that we will be faithful over that which he has entrusted to our care \sout{and} that we	may obtain much: and that our children after us shall remember to observe this sacred and holy covenant: \sout{after} \sout{us:} And that our children and our children's may know of the same we here subscribe our names with our own hands before the Lord:
					\end{quote}

					\begin{tightright}
					\textbf{Joseph Smith Jr}
					
					Oliver Cowdery.
					\end{tightright}
					
						\hypertarget{EMS-1834-JS-29-NOV-1834-b}%
					And now, O Father, as thou didst prosper
						\marginnote{b}%
					our father Jacob, and bless him with protection and prosperity where ever he went from the time he made a like covenant before and with thee; and as thou didst,--- even the same night, open the heavens unto him and manifest great mercy and favor, and give him promises, so wilt thou do by us his sons; and as his blessings prevailed above the blessings of his Progenitors into the utmost bounds of the everlasting hills, even so may our blessings prevail \sout{above} <like> his; and may thy servants be preserved from the power and influence of wicked and unrighteous men; may every weapon formed against us fall upon the head of him who shall form it; may we be blessed with a name and a place among thy saints here, and thy sanctified when they shall rest. Amen.
					
		% -----
		\subsubsection{Revelation, 5 December 1834}
		% -----
			\label{EMS-1834-JS-5-DEC-1834}
			% \label{EMS-1830-JS-DAY-3LETTERMONTH-YEAR}
			
			\begin{description}
					\item[Print Sources]
				\end{description}

					\begin{tabular}{p{0.5in}p{1in}p{0.5in}p{1in}}
					JSP	 	& D4:197			& JSR 		& 266--267 \\
					DC	 	& --- 				& HC 		& 2:176--177\footnote{Note that this is contained in a footnote.} \\
					HDDC 	& ---				& TCHC 		& --- \\
					\end{tabular}
					% HDDC = woodford's DC diss; JSR = Marquardt's JS Revelations; TCHC = Vogel HC
				
				\begin{description}
					\item[Online Access] \href{https://www.josephsmithpapers.org/paper-summary/revelation-5-december-1834/1}{Here}
					% \item[Analysis] \hyperref[XX]{Here}
				\end{description}
				
				\begin{description}
					\item[Background]
				\end{description}
				
					% It might also be useful to give a two sentence context of the period: e.g., "This speech was given in Nauvoo to a small congregation one month prior to the destruction of the Nauvoo Expositor. These and other tensions may have been on the prophet's mind when he gave the speech."
				
				\begin{description}
					\item[Original Source]
				\end{description}
				
					% brief description of original source material, with reference to JSP or somewhere else for more info
					% Background material: it might be helpful to have a note that says whether a given document was presented as a speech, was written in a journal, etc.
					% Also a comment here about how reliable the source might be, or what shape it is in, if there are large omissions or parts that are hard to understand, and so on.
				
				\begin{description}
					\item[Text]
				\end{description}
				
					Verily, condemnation resteth upon you, who are appointed to lead my Chu[r]ch, and to be saviors of men: and also upon the church: And there must needs be a repentance and a refor[m]ation among you, in all things, in your ensamples before the Chuch, and before the world, in all your manners, habits and customs, and salutations one toward another---rendering unto every man the respect due the office, \sout{and} calling, and priesthood, whereunto I the Lord have appointed and ordained you. Amen.
					
		% -----
		\subsubsection{Record, 5 December 1834}
		% -----
			\label{EMS-1834-JS-5-DEC-1834-REC}
			% \label{EMS-1830-JS-DAY-3LETTERMONTH-YEAR}
			
			\begin{description}
					\item[Print Sources]
				\end{description}

					\begin{tabular}{p{0.5in}p{1in}p{0.5in}p{1in}}
					JSP	 	& J1:47--48		& JSR 		& --- \\
					DC	 	& --- 		& HC 		& --- \\
					HDDC 	& ---		& TCHC 		& --- \\
					\end{tabular}
					% HDDC = woodford's DC diss; JSR = Marquardt's JS Revelations; TCHC = Vogel HC
				
				\begin{description}
					\item[Online Access] \href{https://www.josephsmithpapers.org/paper-summary/journal-1832-1834/94}{Here}
					% \item[Analysis] \hyperref[XX]{Here}
				\end{description}
				
				\begin{description}
					\item[Background]
				\end{description}
				
					% It might also be useful to give a two sentence context of the period: e.g., "This speech was given in Nauvoo to a small congregation one month prior to the destruction of the Nauvoo Expositor. These and other tensions may have been on the prophet's mind when he gave the speech."
				
				\begin{description}
					\item[Original Source]
				\end{description}
				
					% brief description of original source material, with reference to JSP or somewhere else for more info
					% Background material: it might be helpful to have a note that says whether a given document was presented as a speech, was written in a journal, etc.
					% Also a comment here about how reliable the source might be, or what shape it is in, if there are large omissions or parts that are hard to understand, and so on.
					
					% The second part after the signatures is not in Documents (though I should check again); it's only in Journals.
				
				\begin{description}
					\item[Text]
				\end{description}
				
					Friday Evening, December 5, 1834. According to the directions of the Holy Spirit breth[r]en Joseph Smith jr. Sidney [Rigdon], Frederick G. Williams, and Oliver Cowdery, assembled to converse upon the welfare of the church, when brother Oliver Cowdery was ordained an assistant President of the High and Holy Priesthood under the hands of brother Joseph Smith jr. Saying, ``My brother, in the name of Jesus Christ who \sout{died} was crucified for the sins of the world, I lay my hands upon thee, and ordain thee an assistant President of the high and holy p[r]iesthood in the church of the Latter Day Saints
					
		% -----
		\subsubsection{Account of Meetings, Revelation, and Blessing, 5--6 December 1834}
		% -----
			\label{EMS-1834-JS-5-6-DEC-1834}
			% \label{EMS-1830-JS-DAY-3LETTERMONTH-YEAR}
			
			\begin{description}
					\item[Print Sources]
				\end{description}

					\begin{tabular}{p{0.5in}p{1in}p{0.5in}p{1in}}
					JSP	 	& D4:191--200		& JSR 		& --- \\
					DC	 	& --- 		& HC 		& 2:176 \\
					HDDC 	& ---		& TCHC 		& 2:190 \\
					PJS 	& 1:20--25			& 	 		&  \\
					\end{tabular}
					% HDDC = woodford's DC diss; JSR = Marquardt's JS Revelations; TCHC = Vogel HC
				
				\begin{description}
					\item[Online Access] \href{https://www.josephsmithpapers.org/paper-summary/account-of-meetings-revelation-and-blessing-5-6-december-1834/1}{Here}
					% \item[Analysis] \hyperref[XX]{Here}
				\end{description}
				
				\begin{description}
					\item[Background]
				\end{description}
				
					% It might also be useful to give a two sentence context of the period: e.g., "This speech was given in Nauvoo to a small congregation one month prior to the destruction of the Nauvoo Expositor. These and other tensions may have been on the prophet's mind when he gave the speech."
				
				\begin{description}
					\item[Original Source]
				\end{description}
				
					% brief description of original source material, with reference to JSP or somewhere else for more info
					% Background material: it might be helpful to have a note that says whether a given document was presented as a speech, was written in a journal, etc.
					% Also a comment here about how reliable the source might be, or what shape it is in, if there are large omissions or parts that are hard to understand, and so on.
				
				\begin{description}
					\item[Text]
				\end{description}
				
					\begin{tightcenter}
							\hypertarget{EMS-1834-JS-5-6-DEC-1834-a}%
						Chapter
							\marginnote{a}
						1.
					\end{tightcenter}

					Friday Evening, December 5, 1834. According to the direction of the Holy Spirit, Pr[e]sident Smith, assistant Presidents, [Sidney] Rigdon and [Frederick G.] Williams, assembled for the purpose of ordaining <first> High Counsellor [Oliver] Cowdery to the office of assistant President of the High and Holy Priesthood in the Church of the Latter-Day Saints.

					It is necessary, for the special benefit of the reader, that he be instructed <into, or> concerning the power and authority of the above named Priesthood.

						\hypertarget{EMS-1834-JS-5-6-DEC-1834-b}%
					First.
						\marginnote{b}%
					The office of the President is to preside over the whole Chu[r]ch; to be considered as at the head; to receive revelations for the Church; to be a Seer, \sout{and} Revelator <and Prophet---> having all the gifts of God:--- \sout{having} taking <Moses> for an ensample. Which is

					\sout{Second.} the office and station of the above President Smith, according to the calling of God, and the ordination which he has received.

						\hypertarget{EMS-1834-JS-5-6-DEC-1834-c}%
					Second.
						\marginnote{c}%
					The office of Assistant President is to assist in presiding over the whole chu[r]ch, and to officiate in the abscence of the President, according to \sout{their} <his> rank and appointment, viz: President Cowdery, first; President Rigdon Second, and President Williams Third, as they <were> \sout{are} severally called. The office of this Priesthood is also to act as Spokesman---taking Aaron for an ensample.

						\hypertarget{EMS-1834-JS-5-6-DEC-1834-d}%
					The
						\marginnote{d}%
					virtue of \sout{this} the <above> Priesthood is to hold the keys of the kingdom of heaven, or the Church militant.

					The reader may further understand, that \sout{Presidents} <the> reason why \sout{President} <High Counsellor> Cowdery was not previously ordained <to the Presidency,> was, in consequence of his necessary attendance in Zion, to assist Wm W. Phelps in conducting the printing business; but that this promise was made by the angel while in company with President Smith, at the time they recievd the office of the lesser priesthood. And further: The circumstances and situation of the Church requiring, Presidents Rigdon and Williams were previously ordained, to assist \sout{the} President Smith.

						\hypertarget{EMS-1834-JS-5-6-DEC-1834-e}%
					After
						\marginnote{e}%
					this short explination, we now proceed to give an account of the acts, promises, and blessings of this memorable Evening:

					First. After assembling, we received a rebuke for our former low, uncultivated, and disrespectful manner of communication, and salutation, with, and unto each other, by the voice of the Spirit, saying unto us:

					\begin{quote}
						Verily, condemnation resteth upon you, who are appointed to lead my Chu[r]ch, and to be saviors of men: and also upon the church: And there must needs be a repentance and a refor[m]ation among you, in all things, in your ensamples before the Chuch, and before the world, in all your manners, habits and customs, and salutations one toward another---rendering unto every man the respect due the office, \sout{and} calling, and priesthood, whereunto I the Lord have appointed and ordained you. Amen.
					\end{quote}

						\hypertarget{EMS-1834-JS-5-6-DEC-1834-f}%
					It
						\marginnote{f}%
					is only necessary to say, relative to the foregoing reproof and instruction, that, though it was given in sharpness, it occasioned gladness and joy, and we were willing to repent and reform, in every particular, according to the instruction given. It is also proper to remark, that after the reproof was given, we all confessed, voluntarily, that such had been the manifestations of the Spirit a long time\sout{s} since; in consequence of which the rebuke came with greater sharpness.

						\hypertarget{EMS-1834-JS-5-6-DEC-1834-g}%
					Not
						\marginnote{g}%
					thinking to evade the truth, or excuse, in order to escape censure, but to give proper information, a few remarks relative to the situation of the chu[r]ch previous to this date, is necessary. Many, on hearing the fulness of the gospel, embraced it with eagerness; <yet,> at the same time were unwilling to forego their former opinions and notions relative to church government, and the rules and habits proper for the good order, harmony, peace, and beauty of a people destined, with the protecting care of the Lord, to be an ensample and light of the world.
						\hypertarget{EMS-1834-JS-5-6-DEC-1834-h}%
					They
						\marginnote{h}%
					did not dispise government; but there was a disposition to organize that government according to their own notions, or feelings. For example: Every man must be subjected <to> wear a particular fashioned coat, hat, or other garment, or else an accusation was brought that we were fashioning after the world. Every one must be called by their given name, without respecting the office or ordinance to which they had been called: Thus, President Smith was called Joseph, or brother Joseph; President Rigdon, brother Sidney, or Sidney, \&c. This manner of address gave occasion to the enemies of the truth, and was a means of bringing reproach upon the Cause of God.
						\hypertarget{EMS-1834-JS-5-6-DEC-1834-i}%
					But
						\marginnote{i}%
					in consequence of former prejudices, the church, many of them, would not submit to proper and wholesome order. This proceeded from a spirit of enthusiasm, and vain ambition---a desire to compel others to come to certain rules, not dictated by the will of the Lord; or a jealous fear, that, were men called by thier respective titles, and the ordinance of heaven honored in a proper manner, some were in a way to be exalted above others, and their form of government disregarded.
						\hypertarget{EMS-1834-JS-5-6-DEC-1834-j}%
					In
						\marginnote{j}%
					fact, the true principle of honor in the church of the saints, that the more a man is exalted, the more humble he will be, if actuated by the Spirit of the Lord, seemed to have been overlooked; and the fact, that the greatest is least and servant of all; as said our Savior, never to have been thought of, by numbers. These facts, for such they were, when viewed in their proper light, were sufficient, of themselves to cause men to humble themselves before the Lord; but when communicated by the Spirit, made an impression upon our hearts not to be forgotten.

						\hypertarget{EMS-1834-JS-5-6-DEC-1834-k}%
					Perhaps,
						\marginnote{k}%
					an arrangement of this kind in a former day would have occasioned some unpleasant reflections, in the minds of many, and at an \sout{early} <earlier> period, in this church, others to have forsaken the cause, in consequence of weakness, and unfaithfulness; but that the leaders of the church should wait so long before stepping forward according <to> the manifestation of the Spirit, deserved a reproof. And that the church should be chastened, for their uncultivated manner of salutation, is also just. But to proceed with the account of the interview.

						\hypertarget{EMS-1834-JS-5-6-DEC-1834-l}%
					After
						\marginnote{l}%
					addressing the throne of mercy, President Smith laid hands upon High Counsellor Cowdery, and ordained him to the Presidency of the High priesthood in the Church, saying:

					\begin{quote}
						Brother, In the name of Jesus Christ of Nazareth, who was crucified for the sins of the world, that we through the virtue of his blood might come to the Father, I lay my hands upon thy head, and ordain thee a President of the high and holy priesthood, to assist in presiding over the Chu[r]ch, and bearing the keys of this kingdom---which priesthood is after the order of Melchizedek---which is after the order of the Son of God---
							\hypertarget{EMS-1834-JS-5-6-DEC-1834-m}%
						And
							\marginnote{m}%
						now, O Father, wilt thou bless this thy servant with wisdom, knowledge, and understanding--- give him, by the Holy Spirit, a correct understanding of thy doctrine, laws, and will--- commune with him from on high--- let him hear thy voice, and recieve the \sout{ministries} ministring of the holy angels--- deliver him from temptation, and the power of darkness--- deliver him from evil, and from those who may seek his destruction,--- be his shield, his buckler, and his great reward---
							\hypertarget{EMS-1834-JS-5-6-DEC-1834-n}%
						endow
							\marginnote{n}%
						him with power from on high, that he may write, preach, and proclaim the gospel to his fellowmen in demonstration of the Spirit and of power--- may his feet never slide--- may his heart never feint--- may his faith never fail. Bestow upon him the blessings of his fathers Abraham, Isaac, Jacob, and of Joseph--- Prolong his life to a good old age, and bring him in peace to his end, and to rejoice with thy saints, even the sanctified, in the celestial kingdom; for thine is the kingdom, the power, and the glory, forever. Amen.
					\end{quote}

					Presidents Rigdon and Willliams, confirmed the ordinance and blessings by the laying on of hands and prayer, after which each were blessed with the same blessings and prayer.

						\hypertarget{EMS-1834-JS-5-6-DEC-1834-o}%
					Much
						\marginnote{o}%
					light was communicated to our minds, and we were instructed into the order of the church of the saints, and how they ought to conduct in respecting and reverencing each other. The praise of men, or the honor of this world, is of no benefit; but if a man is respected in his calling, and considered to be a man of righteousness, the truth may have an influence, many times, by which means they may teach the gospel with success, and lead men into the kingdom of heaven.

					On Saturday, December 6, Presidents Smith, Cowdery, and Rigdon assembled with High Counsellors Joseph Smith, sen. Hyrum Smith, and Samuel H. Smith, in company with Reynolds Cahoon, counsellor to the Bishop, High Priest William Smith, and <Elder> Don C[arlos] Smith.

						\hypertarget{EMS-1834-JS-5-6-DEC-1834-p}%
					The
						\marginnote{p}%
					meeting was opened by prayer, and a lengthy conversation held upon the subject of introducing a more refined order into the Church. On further reflection, the propriety of ordaining others to the office of Presidency of the high priesthood was also discussed, after which High Counsellor Hyrum Smith was ordained <to> the Presidency under the hands of President Smith, and High Counsellor Joseph Smith sen. under the hands of President Rigdon. The others present were blessed under the hands of Presidents J. Smith jr. Cowdery, and Rigdon, and the meeting closed, after a happy season, \sout{of} and a social intercourse upon the great subject of the gospel and the work of the Lord in this day.
					
		% -----
		\subsubsection{Blessing from Joseph Smith Sr., 9 December 1834}
		% -----
			\label{EMS-1834-JS-9-DEC-1834}
			% \label{EMS-1830-JS-DAY-3LETTERMONTH-YEAR}
			
			\begin{description}
					\item[Print Sources]
				\end{description}

					\begin{tabular}{p{0.5in}p{1in}p{0.5in}p{1in}}
					JSP	 	& D4:200--208		& JSR 		& --- \\
					DC	 	& --- 		& HC 		& --- \\
					HDDC 	& ---		& TCHC 		& --- \\
					\end{tabular}
					% HDDC = woodford's DC diss; JSR = Marquardt's JS Revelations; TCHC = Vogel HC
				
				\begin{description}
					\item[Online Access] \href{https://www.josephsmithpapers.org/paper-summary/blessing-from-joseph-smith-sr-9-december-1834/1}{Here}
					% \item[Analysis] \hyperref[XX]{Here}
				\end{description}
				
				\begin{description}
					\item[Background]
				\end{description}
				
					% It might also be useful to give a two sentence context of the period: e.g., "This speech was given in Nauvoo to a small congregation one month prior to the destruction of the Nauvoo Expositor. These and other tensions may have been on the prophet's mind when he gave the speech."
				
				\begin{description}
					\item[Original Source]
				\end{description}
				
					% brief description of original source material, with reference to JSP or somewhere else for more info
					% Background material: it might be helpful to have a note that says whether a given document was presented as a speech, was written in a journal, etc.
					% Also a comment here about how reliable the source might be, or what shape it is in, if there are large omissions or parts that are hard to understand, and so on.
				
				\begin{description}
					\item[Text]
				\end{description}
				
						\hypertarget{EMS-1834-JS-9-DEC-1834-a}%
					Joseph
						\marginnote{a}%
					Smith, junior was born in Sharon, Windsor County, Vermont, December 23, 1805.

					Joseph, my son, I lay my hands upon thy head in the name of the Lord Jesus Christ, to confirm upon thee a father's blessing. The Lord thy God has called thee by name out of the heavens: thou hast heard his voice from on high from time to time, even in thy youth. The hand of the angel of his presence has been extended toward thee by which thou hast been lifted up and sustained; yea, the Lord has delivered thee from the hands of thine enemies and thou hast been made to rejoice in his salvation: thou hast sought to know his ways, and from thy childhood thou hast meditated much upon the great things of his law.
						\hypertarget{EMS-1834-JS-9-DEC-1834-b}%
					Thou
						\marginnote{b}%
					hast suffered much in thy youth, and the poverty and afflictions of thy father's family have been a grief to thy soul. Thou hast desired to see them delivered from bondage, for thou hast lov'd them with a perfect love. Thou hast stood by thy father, and like Shem, would have covered his nakedness, rather than see him exposed to shame: when the daughters of the Gentiles laughed, thy heart has been moved with a just anger to avenge thy kindred. Thou hast been an obedient son: the commands of thy father and the reproofs of thy mother, thou hast respected and obeyed--- for all these things the Lord my God will bless thee. Thou hast been called, even in thy youth to the great work of the Lord: to do a work in this generation which no other man would do as thyself, in all things according to the will of the Lo[r]d.
						\hypertarget{EMS-1834-JS-9-DEC-1834-c}%
					A
						\marginnote{c}%
					marvelous work and a wonder has the Lord wrought by thy hand, even that which shall prepare the way for the remnants of his people to come in among the Gentiles, with their fulness, as the tribes of Israel are restored. I bless thee with the blessings of thy fathers Abraham, Isaac and Jacob; and even the blessings of thy father Joseph, the son of Jacob. Behold, he looked after his posterity in the last days, when they should be scattered and driven by the Gentiles, and wept before the Lord: he sought diligently to know from whence the son should come who should bring forth the word of the Lord, by which they might be enlightened, and brought back to the true fold, and his eyes beheld thee, my son: his heart rejoiced and his soul was satisfied, and he said, As my blessings are to extend to the utmost bounds of the everlastings hills; as my father's blessing prevailed above the blessings of his progen[i]tors,
						\hypertarget{EMS-1834-JS-9-DEC-1834-d}%
					and
						\marginnote{d}%
					as my branches are to run over the wall, and my seed are to inherit the choice land whereon the Zion of God shall stand in the last days, from among my seed, scattered with the Gentiles, shall a choice Seer arise whose bowels shall be as a fountain of truth, whose loins shall be girded with the girdle of righteousness, whose hands shall be lifted with acceptance before the God of Jacob to turn away his anger from his annointed, whose heart shall meditate great wisdom, whose intelligence shall circumscribe \sout{and} <and> comprehend the deep things of God, and whose mouth shall utter the law of the just:
						\hypertarget{EMS-1834-JS-9-DEC-1834-e}%
					His
						\marginnote{e}%
					feet shall stand upon the neck of his enemies, and he shall walk upon the ashes of those who seek his destruction: with wine and oil \sout{it} shall he be sustained, and he shall feed upon the heritage of Jacob his father: the just shall desire his society, and the upright in heart shall be his companions: No weapon formed against him shall prosper, and though the wicked mar him for a little season, he shall be like one rising up in the heat of wine---
						\hypertarget{EMS-1834-JS-9-DEC-1834-f}%
					he
						\marginnote{f}%
					shall roar in his strength, and the Lord shall put to flight his persecutors: \sout{his} he shall be blessed like the fruitful olive, and his memory shall be as sweet as the choice cluster of the first ripe grapes. Like <a> shief [sheaf] fully ripe, gathered into the garner, so shall he stand before the Lord, having produced a hundred fold. Thus spake my father Joseph. Therefore, my son, I know for a surety that these things will be fulfilled, and I confirm upon thee all these blessings. Thou shalt live <to> do the work which the Lord shall command thee: thou shalt hold the keys of this ministry, even the presidency of this church, both in time and in eternity. Thy heart shall be enlarged, and thou shalt be able to fill up the measure of thy days according to the will of the Lord.
						\hypertarget{EMS-1834-JS-9-DEC-1834-g}%
					Thou
						\marginnote{g}%
					shalt speak the word of the Lord and the earth shall tremble; the mountains shall move and the rivers shall turn out of their course. Thou shalt escape the edge of the sword, and put to flight the armies of the wicked. At thy word the lame shall walk, the deaf shall hear and the blind shall see. Thou shalt be gathered to Zion and in the goodly land thou shalt enjoy thine inheritance; thy children and thy children's children to the latest generation; for thy name and the names of thy posterity shall be recorded in the book of the Lord, even in the book of blessings \sout{an} <and> genealogies, for their joy and benefit forever.
						\hypertarget{EMS-1834-JS-9-DEC-1834-h}%
					And
						\marginnote{h}%
					now, my son, what more shall I say? Thou art as a fruitful olive and a choice vine: thou shalt be laden with precious fruit. Thousands and tens of thousands shall come to a knowledge of the truth through thy ministry, and thou shalt rejoice with them in the Celestial Kingdom: Thou shalt stand upon the earth when it shall reel to and fro as a drunken man, and be removed out of its place: thou shalt stand when the mighty judgments go forth to the destruction of the wicked: thou shalt stand on mount Zion when the tribes of Jacob come shouting from the north, and with thy brethren, the sons of Ephraim, crown them in the name of Jesus Christ: Thou shalt see thy Redeemer Come in the clouds of heaven, and with the just receive the hallowed throng with shouts of hallalujahs, praise the Lord. Amen

						\hypertarget{EMS-1834-JS-9-DEC-1834-i}%
					Emma
						\marginnote{i}%
					Smith, wife of Joseph Smith jr. was born in Harmony, Susqehannah County, Pennsylvania, July 10, 1804.

					Emma, my daughter-in-law, thou art blessed of the Lord, for thy faithfulness and truth: thou shalt be blessed with thy husband, and rejoice in the glory which shall come upon him: Thy soul has been afflicted because of the wickedness of men in seeking the destruction of thy companion, and thy whole soul has been drawn out in prayer for his deliverance: rejoice, for the Lord thy God has heard thy suplication. Thou hast grieved for the hardness of the hearts of thy father's house, and thou hast longed for their salvation.
						\hypertarget{EMS-1834-JS-9-DEC-1834-j}%
					The
						\marginnote{j}%
					Lord will have respect to thy cries, and by his judgments he will cause some of them to see their folly and repent of their sins; but it will be by affliction that they will be saved. Thou shall see many days; yea, the Lord will spare thee till thou art satisfied, for thou shalt see thy Redeemer. Thy heart shall rejoice in the great work of the Lord, and no one shall take thy rejoicing from thee. Thou shalt ever remember the great condescension of thy God in permitting thee to accompany my son when the angel delivered the record of the Nephites to his care. Thou hast seen much sorrow because the Lord has taken from thee three of thy children: in this thou are not to be blamed, for he knows thy pure desires to raise up a family, that the name of my son might be blessed.
						\hypertarget{EMS-1834-JS-9-DEC-1834-k}%
					And
						\marginnote{k}%
					now, behold, I say unto thee, that thus says the Lord, if thou wilt believe, thou shalt yet be blessed in this thing <and> thou shalt bring forth other children, to the joy and satisfaction of thy soul, and to the rejoicing of thy friends. Thou shalt be blessed with understanding, and have power to instruct thy sex. Teach thy family righteousness, and thy little ones the way of life, and the holy angels shall watch over thee: and thou shalt be saved in the kingdom of God; even so. Amen.

					\begin{tightright}
					Oliver Cowdery,\} \uline{Clerk and Recorder}.
					\end{tightright}

					Kirtland, December 9, 1834.
					
		% -----
		\subsubsection{Minutes, 28 December 1834}
		% -----
			\label{EMS-1834-JS-28-DEC-1834}
			% \label{EMS-1830-JS-DAY-3LETTERMONTH-YEAR}
			
			\begin{description}
					\item[Print Sources]
				\end{description}

					\begin{tabular}{p{0.5in}p{1in}p{0.5in}p{1in}}
					JSP	 	& D4:208--211		& JSR 		& --- \\
					DC	 	& --- 		& HC 		& --- \\
					HDDC 	& ---		& TCHC 		& --- \\
					\end{tabular}
					% HDDC = woodford's DC diss; JSR = Marquardt's JS Revelations; TCHC = Vogel HC
				
				\begin{description}
					\item[Online Access] \href{https://www.josephsmithpapers.org/paper-summary/minutes-28-december-1834/1}{Here}
					% \item[Analysis] \hyperref[XX]{Here}
				\end{description}
				
				\begin{description}
					\item[Background]
				\end{description}
				
					% It might also be useful to give a two sentence context of the period: e.g., "This speech was given in Nauvoo to a small congregation one month prior to the destruction of the Nauvoo Expositor. These and other tensions may have been on the prophet's mind when he gave the speech."
				
				\begin{description}
					\item[Original Source]
				\end{description}
				
					% brief description of original source material, with reference to JSP or somewhere else for more info
					% Background material: it might be helpful to have a note that says whether a given document was presented as a speech, was written in a journal, etc.
					% Also a comment here about how reliable the source might be, or what shape it is in, if there are large omissions or parts that are hard to understand, and so on.
				
				\begin{description}
					\item[Text]
				\end{description}
				
					...Many excellent remarks were made relative to the order of offices and titles in the church by Presidents J. Smith \supersu{Jun}\supers{r.} \& S. Rigdon when a vote was called to know whether all present were satisfied with those remarks, which was also unanimous, there being no dissenting voice...
					
		% -----
		\subsubsection{Letter to Oliver Cowdery, December 1834}
		% -----
			\label{EMS-1834-JS-DEC-1834}
			% \label{EMS-1830-JS-DAY-3LETTERMONTH-YEAR}
			
			\begin{description}
					\item[Print Sources]
				\end{description}

					\begin{tabular}{p{0.5in}p{1in}p{0.5in}p{1in}}
					JSP	 	& D4:211--215		& JSR 		& --- \\
					DC	 	& --- 		& HC 		& --- \\
					HDDC 	& ---		& TCHC 		& --- \\
					PJS 	& 1:11--14			& 	 		&  \\
					\end{tabular}
					% HDDC = woodford's DC diss; JSR = Marquardt's JS Revelations; TCHC = Vogel HC
				
				\begin{description}
					\item[Online Access] \href{https://www.josephsmithpapers.org/paper-summary/letter-to-oliver-cowdery-december-1834/1}{Here}
					% \item[Analysis] \hyperref[XX]{Here}
				\end{description}
				
				\begin{description}
					\item[Background]
				\end{description}
				
					% It might also be useful to give a two sentence context of the period: e.g., "This speech was given in Nauvoo to a small congregation one month prior to the destruction of the Nauvoo Expositor. These and other tensions may have been on the prophet's mind when he gave the speech."
				
				\begin{description}
					\item[Original Source]
				\end{description}
				
					% brief description of original source material, with reference to JSP or somewhere else for more info
					% Background material: it might be helpful to have a note that says whether a given document was presented as a speech, was written in a journal, etc.
					% Also a comment here about how reliable the source might be, or what shape it is in, if there are large omissions or parts that are hard to understand, and so on.
				
				\begin{description}
					\item[Text]
				\end{description}
				
						\hypertarget{EMS-1834-JS-DEC-1834-a}%
					\phantom{ }\textsc{Brother}
						\marginnote{a}
					\phantom{ }\textsc{O[liver] Cowdery}:

					Having learned from the first No. of the Messenger and Advocate, that you were, not only about to ``give a history of the rise and progress of the church of the Latter Day Saints;'' but, that said ``history would necessarily embrace my life and character,'' I have been induced to give you the time and place of my birth; as I have learned that many of the opposers of those principles which I have held forth to the world, profess a personal acquaintance with me, though when in my presence, represent me to be another person in age, education, and stature, from what I am.

						\hypertarget{EMS-1834-JS-DEC-1834-b}%
					I
						\marginnote{b}%
					was born, (according to the record of the same, kept by my parents,) in the town of Sharon, Windsor Co. Vt. on the 23rd of December, 1805.

					At the age of ten my father's family removed to Palmyra, N.Y. where, and in the vicinity of which, I lived, or, made it my place of residence, until I was twenty one---the latter part, in the town of Manchester.

						\hypertarget{EMS-1834-JS-DEC-1834-c}%
					During
						\marginnote{c}%
					this time, as is common to most, or all youths, I fell into many vices and follies; but as my accusers are, and have been forward to accuse me of being guilty of gross and outragious violations of the peace and good order of the community, I take the occasion to remark, that, though, as I have said above, ``as is common to most, or all youths, I fell into many vices and follies,'' I have not, neither can it be sustained, in truth, been guilty of wronging or injuring any man or society of men; and those imperfections to which I alude, and for which I have often had occasion to lament, were a light, and too often, vain mind, exhibiting a foolish and trifling conversation.

						\hypertarget{EMS-1834-JS-DEC-1834-d}%
					This
						\marginnote{d}%
					being all, and the worst, that my accusers can substantiate against my moral character, I wish to add, that it is not without a deep feeling of regret that I am thus called upon in answer to my own conscience, to fulfill a duty I owe to myself, as well as to the cause of truth, in making this public confession of my former uncircumspect walk, and unchaste conversation: and more particularly, as I often acted in violation of those holy precepts which I knew came from God. But as the ``Articles and Covenants'' of this church are plain upon this particular point, I do not deem it important to proceed further. I only add, that I do not, nor never have, pretended to be any other than a man ``subject to passion,'' and liable, without the assisting grace of the Savior, to deviate from that perfect path in which \textit{all} men are commanded to walk!

						\hypertarget{EMS-1834-JS-DEC-1834-e}%
					By
						\marginnote{e}%
					giving the above a place in your valuable paper, you will confer a lasting favor upon myself, as an individual, and, as I humbly hope, subserve the cause of righteousness.

					I am, with feelings of esteem, your fellow laborer in the gospel of our Lord.
					
					\begin{tightright}
						JOSEPH SMITH jr.
					\end{tightright}
					
	% ----------------------
	\subsection{Olivery Cowdery}
	% ----------------------
		\label{EMS-1834-OC}

		% -----
		\subsubsection{Introduction to 1834}
		% -----
		
			sdf
			
		% -----
		\subsubsection{Letter to John Whitmer, 1 January 1834}
		% -----
			\label{EMS-1834-OC-1-JAN-1834}
			% \label{EMS-1830-JS-DAY-3LETTERMONTH-YEAR}
			
				% \begin{description}
					% \item[Print Sources]
				% \end{description}

					% \begin{tabular}{p{0.5in}p{1in}p{0.5in}p{1in}}
					% JSP	 	& D3:403--407; J1:25--26		& JSR 		& --- \\
					% DC	 	& --- 				& HC 		& 2:2--3 \\
					% HDDC 	& ---				& TCHC 		& 2:6--7 \\
					% \end{tabular}
					% % HDDC = woodford's DC diss; JSR = Marquardt's JS Revelations; TCHC = Vogel HC
					
				% \begin{description}
					% \item[Online Access] \href{https://www.josephsmithpapers.org/paper-summary/prayer-11-january-1834/}{Here}
					% % \item[Analysis] \hyperref[XX]{Here}
				% \end{description}
				
				% \begin{description}
					% \item[Background]
				% \end{description}
				
					% It might also be useful to give a two sentence context of the period: e.g., "This speech was given in Nauvoo to a small congregation one month prior to the destruction of the Nauvoo Expositor. These and other tensions may have been on the prophet's mind when he gave the speech."
				
				\begin{description}
					\item[Original Source]
				\end{description}
				
					The original source is the Oliver Cowdery Letterbook at the Huntington Library. I am drawing from Word document with transcriptions of some of those letters, which I got from the MPD website.
					
					My paragraphing matches what is in the transcription.
				
					% brief description of original source material, with reference to JSP or somewhere else for more info
					% Background material: it might be helpful to have a note that says whether a given document was presented as a speech, was written in a journal, etc.
					% Also a comment here about how reliable the source might be, or what shape it is in, if there are large omissions or parts that are hard to understand, and so on.
				
				\begin{description}
					\item[Text]
				\end{description}
				
						\hypertarget{EMS-1834-OC-1-JAN-1834-a}%
					Kirtland,
						\marginnote{a}
					January 1, 1834.

					Dear Brother John: 

					This morning I take an opportunity to communicate to you a few lines, in answer to several letters received from you, which as yet remain unanswered: but I have no hesitancy in persuading myself that you will forgive this delay. When you wrote, no doubt, you supposed that all would come safe through the Post Office; but I having reason to suspect that they would be purloined, dare not communicate anything that would do you any injury, even if it should be taken out and you not receive it; in this, no doubt, you will justify me.
						\hypertarget{EMS-1834-OC-1-JAN-1834-b}%
					We
						\marginnote{b}%
					received the name of our former subscribers a few days since, which were mailed on the 3d. of last month, and a letter from bro. Corril to me, directed to bro. Elliott, Sunday evening. Know assuredly, bro. John, that we are anxious to receive intelligence from our brethren who have been driven from their inheritances, and we are pleased to hear that the Governor is likely to give you aid; for we pray continually that the Lord will stir up the hearts of the Rulers \& men in authority, to avenge his children. The Law is sufficient, the constitution was established acccording to the will of Heaven: and all the lack is, for those whose duty it is to see that they are kept inviolable do their duty: pray that this may be the case; for God is able to turn the hearts of all men sufficiently to bring his purposes to pass.
						\hypertarget{EMS-1834-OC-1-JAN-1834-c}%
					You
						\marginnote{c}%
					have requested instruction on the subject of Church Records: I have just conversed with bro. Joseph concerning the same; and he has given me instruction upon his letter published in the 8th no. of the Star; but says that we shall probably receive more on the same, and then we will communicate it. I will say, however, that it is necessary to keep the names of the Saints, \& when a child is brought forward to be blessed by the Elders, it is then necessary to take their names upon the church Record. Put down the name of the man, his place of birth, and when, \&c. and also of his family.
						\hypertarget{EMS-1834-OC-1-JAN-1834-d}%
					If
						\marginnote{d}%
					he begets children after that \& they do not come into the church their names are not known with their brethren in the book of remembrance. apostatis \rule{1cm}{0.15mm} write, opposite his name that he has. If he begets children after that \& they do not come into the church their names are not known with their brethren in the book of remembrance. The names of the saints are to be kept in a book that contains the law of God; this is what is meant in bro. Joseph's letter.
						\hypertarget{EMS-1834-OC-1-JAN-1834-e}%
					In
						\marginnote{e}%
					due time you will receive all these things, and all necessary information. Each family will have its record with the law of the Lord in it; each branch of the church the same in every city; and each city one general record kept by a general clerk. Brother Joseph says, that the item in his letter that says, that the man that is called \&c. and puts forth his hand to steady the ark of God, does not mean that any one had at the time, but it was given for a caution to those in high standing to beware, lest they should fall by the shaft of death as the Lord had said. 

						\hypertarget{EMS-1834-OC-1-JAN-1834-f}%
					Since
						\marginnote{f}%
					I came down I have been informed from a proper source that the Angel Michael is no less than our father Adam, and Gabriel is Noah. I just drop this because I supposed that you would be pleased to know, and I have no disposition to keep back anything from my brethren that I am privileged to know. Brother Gilbert writes a few words concerning my wife, in his letter to bro. Whitney, and says that he was ready on his part to advance money for her expense when bro. Orson \& John came. Please say to him that I am satisfied with him on that point, and am thankful that he was willing so to do, and am under the same obligation for his kindness, as though she had actually come. But I was greatly disappointed, and my heart is often filled with sorrow. You no doubt recollect bro. Joseph's long communication by the hand of bro. Orson and John, that it was the will of the Lord, or wisdom that I should tarry here. 

						\hypertarget{EMS-1834-OC-1-JAN-1834-g}%
					Now
						\marginnote{g}%
					then bro. John could you suppose that I could come to Zion in the spring to tarry? You request the names of the authors of that letter concerning tongues. I have written them once, but conclude you did not receive my letter. Wm. Whitney, and Rhoda Mills wrote the letter referred to; but we presume they wrote nothing that did not transpire, so we did not blame them, we only cautioned you because we saw that you would get into difficulty if you depended on the gift of tongues for revelations. I have not been able to ascertain where my wafe was, although bro. Gilbert says she is well, and you say she is with father \& mother.
						\hypertarget{EMS-1834-OC-1-JAN-1834-h}%
					I
						\marginnote{i}%
					supposed of course you knew, but it would be a satisfaction to me to know also. I want you to see her \& hand her the following \& write me immediately, \& by so doing, you will confer a lasting favor of your unworthy brother. I daily think how many hours I have spent with you and brother Wm. \& I ask myself, shall I not be privileged again to enjoy your society; yea, verily I trust in the Lord that I shall, for I often see you in visions and in dreams. My love to all the family \& brethren. Oliver. P.S. I enclose a fifty dollar note in this which has been donated by the churches, that is forty of it, for the relief of the needy Saints, which you will hand over to bro. Partridge and Phelps, excepting ten dollars which I send to my wife. I feel for her wants, but it is difficult for me to get much money to send. 

						\hypertarget{EMS-1834-OC-1-JAN-1834-i}%
					My
						\marginnote{i}
					dear Elizabeth:

					Our Father in heaven only knows the feelings \& anxieties of my heart toward, \& for you. I consider that it is now more than one year since we commenced keeping house, and out of that time, we were only permitted to spend seven months in each other's society. I looked anxiously for you with bro. Orson H. \& John G. but was disappointed! and I said, O Lord, preserve thy handmaid that thy servant may see her yet in peace, \& enjoy her company here on earth in the flesh, in the name of Jesus; Amen. And what could I say more? I think of your situation daily, and my mind is filled with continual sorrow. I live at bro. Joseph's, and am treated with brotherly kindness; but that is not like living with a family of one's own, when the Lord has given him one.
						\hypertarget{EMS-1834-OC-1-JAN-1834-j}%
					I
						\marginnote{j}%
					would say something about your coming to me, but I know not how to send, who to send nor what way to send. My time is necessarily occupied here: we are poor: we exerted every possible means to pay bro. Gilbert's debts in N. Y. and we loaned money to purchase our press and types. And notwithstanding I am separated from you, and we are in these difficulties, yet I endeavor to stay myself upon the Lord, Indeed, he is my only refuge. Be faithful, be patient, be not cast down, be not discouraged, be not terrified, for God dwells in heaven, and ere long we shall meet at his right hand, where parting will not be known.
						\hypertarget{EMS-1834-OC-1-JAN-1834-k}%
					I
						\marginnote{k}%
					have seen a letter from Diana to Sister Nancy Richardson, which says that she had made an exchange with you of some beads \&c. which she has here; but I am told she has nothing here of any value scarcely, that is not a doll's worth. When you come away, if you do not bring all of your goods with you, leave all that you do not bring with mother, or some of the family, where they can be kept safely, for we shall want them when we return. I give you this, that you may be wise and I have no doubt you will. 

						\hypertarget{EMS-1834-OC-1-JAN-1834-l}%
					I
						\marginnote{l}%
					enclose fifty dollars in this for the poor brethren, but ten of it is for you. You will take it therefore \& use if for your own benefit, as you see proper. The mail starts soon, so I must close. 

					My love to father \& mother \& all of the family. 

					I am, affectionately, yours, Oliver. 

					P.S. We have sent the Star to some of the brethren, and want bro. Wm. to send us the names of more if they want. 
							
% -----------------------------------------------------------
\pagebreak{}
% -----------------------------------------------------------